\documentclass[12pt,a4paper]{amsart}
\usepackage{amsmath,amssymb,amsthm}
\usepackage{mathtools}
\usepackage{array}
\usepackage[final,expansion=false]{microtype}
\usepackage{geometry}
\geometry{margin=1in}
\raggedbottom

% Theorem environments
\newtheorem{theorem}{Theorem}[section]
\newtheorem{lemma}[theorem]{Lemma}
\newtheorem{proposition}[theorem]{Proposition}
\newtheorem{corollary}[theorem]{Corollary}
\newtheorem{definition}[theorem]{Definition}
\theoremstyle{remark}
\newtheorem{remark}[theorem]{Remark}

\title{Microlocal Spectral Bifurcation in Schwarzschild--de~Sitter Geometry}

\author{Caner Ate\c{s}}
\address{Istanbul, Turkey}
\email{canerates053@gmail.com}
\date{June 2026}

\subjclass[2020]{%
  35P25,  % Scattering theory for PDEs
  58J50,  % Spectral theory; eigenvalue problems on manifolds
  83C57,  % Black holes
  35S05,  % Pseudodifferential operators
  34E20   % Singular perturbations, turning point theory, WKB methods
}

\keywords{quasinormal modes; Jordan block; semiclassical analysis;
Schwarzschild--de~Sitter spacetime; WKB suppression; microlocal analysis;
spectral bifurcation; Regge--Wheeler operator; resolvent estimates}

\begin{document}

\begin{abstract}
We develop a complete semiclassical resolvent theory for the Regge--Wheeler
operator on Schwarzschild--de~Sitter (SdS) spacetimes. First, we demonstrate that in
the strictly physical regime ($r_s \leq r_s^{\mathrm{Nariai}}$), the photon sphere
is permanently shielded from geometric degeneracy, thereby maintaining strict normal
hyperbolicity. Second, by analytically continuing the parameter space beyond the
Nariai boundary, we identify an exact algebraic locus $P_c$ where the spectral gap
of the linearized Hamiltonian flow completely collapses. By performing a microlocal
reduction to a $k=3$ (cusp) normal form and imposing purely outgoing boundary
conditions via complex scaling, we establish the emergence of a non-diagonalizable
Jordan block of size~3. We uncover an asymptotic WKB suppression mechanism that
shields the system from generic Puiseux fractional scaling and yields the two-sided
resolvent estimate
\[
  \|R_h(P)\| \sim h^{-6/5}|P - P_c|^{-1/2},
\]
valid in the semiclassical limit $h\to 0$ at any fixed deformation $P - P_c \neq 0$.
In the coupled double scaling regime $|P - P_c| \sim h^{4/5}$, the resolvent remains
uniformly bounded by $O(h^{-6/5})$ and the resonance spacing satisfies
$\Delta\omega \sim h^{6/5}$; these two asymptotic windows are complementary and
consistent. We show that the critical frequency satisfies
$\omega_c^2 = V_{\mathrm{eff}}(r^c;r_s^c) < 0$, so $\omega_c \in i\mathbb{R}$ is
purely imaginary; the three coalescing resonances therefore lie on the imaginary axis
at $P_c$, and the splitting $\Delta\omega$ occurs in the complex plane at angle
$-\pi/10$ from the imaginary axis.
\end{abstract}

\maketitle

% -------------------------------------------------------
\section{Introduction}
% -------------------------------------------------------

\subsection*{Physical motivation}

Schwarzschild--de~Sitter (SdS) spacetimes provide the canonical arena for studying
gravitational wave emission and quasinormal mode (QNM) spectra in the presence of a
positive cosmological constant $\Lambda > 0$. The Regge--Wheeler equation governs
linearized gravitational perturbations of spin $s=2$ around the static spherically
symmetric SdS background, and its resonance spectrum, which consists of the QNMs, encodes
the characteristic ringdown frequencies observable in gravitational wave experiments.
Consequently, understanding the behavior of these resonances as the spacetime parameters
approach critical configurations is of direct physical and observational significance.

The SdS geometry is characterized by two positive parameters: the Schwarzschild mass
$r_s$ and the cosmological length scale $L = \sqrt{3/\Lambda}$. The spacetime possesses
two horizons---a black hole horizon at $r_h$ and a cosmological horizon at $r_c$---that
coalesce in the Nariai limit $r_s \to r_s^{\mathrm{Nariai}} = \frac{2\sqrt{3}}{9}L$,
where the geometry degenerates. For $r_s < r_s^{\mathrm{Nariai}}$, the spacetime is
physically regular, the trapped photon sphere at the potential maximum is normally
hyperbolic, and the standard QNM theory of Wunsch--Zworski~\cite{WZ11} and
Dyatlov~\cite{Dyatlov11} applies to yield an exponential resonance-free strip below
the real axis. In this regime, the photon sphere controls the long-time decay of
linear perturbations.

A natural question arises regarding the behavior of the effective potential when it
is tuned to a degenerate configuration. For generic parameters $(r_s, L)$, the effective
potential $V_{\mathrm{eff}}(r)$ possesses a nondegenerate maximum. However, it is
important to determine whether the simultaneous vanishing $V'_{\mathrm{eff}} = V''_{\mathrm{eff}} = 0$
is achievable. In this work, we show that such a configuration exists at a unique critical point
$(r_s^c, L) \approx (1.1192\,L, L)$ in the parameter space, but lies strictly \emph{outside}
the physical regime $r_s \leq r_s^{\mathrm{Nariai}} \approx 0.3849\,L$ (Lemma~\ref{lem:physical_shielding}).
The critical configuration must therefore be accessed by analytic continuation in $r_s$
beyond the Nariai boundary, and the resulting spectral phenomena are a mathematical
consequence of this continuation, rather than a physical instability of any SdS spacetime.

We emphasize that analytic continuation beyond the Nariai boundary is mathematically
well-defined and physically meaningful: the resolvent family $R_h(P)$, initially
constructed for $r_s < r_s^{\mathrm{Nariai}}$ via the standard microlocal theory of
Wunsch--Zworski~\cite{WZ11} and Dyatlov~\cite{Dyatlov11}, extends meromorphically
in the parameter $r_s$ by the analytic Fredholm theorem
(Dyatlov--Zworski~\cite{DZ19}, Theorem~C.5).
The critical point $P_c$ is then a distinguished pole of this meromorphic family,
and the Jordan block structure at $P_c$ is an intrinsic property of the continued
resolvent---not an artifact of the continuation procedure.
Physically, the analysis describes the \emph{boundary behavior} of the QNM spectrum
as the SdS geometry is driven toward its mathematical limits: the spectral phenomena
uncovered here govern the rate at which the resolvent degenerates as $P \to P_c$,
and this rate is observable as a sharp threshold in the $h \to 0$ asymptotics,
independent of whether $P_c$ itself is physically realized.
In particular, the continuation does not alter the geometry of any SdS spacetime:
it is an operation on the parameter-dependent resolvent family, not on the underlying
manifold, and the Jordan block structure at $P_c$ is therefore an intrinsic spectral
invariant of the SdS resolvent rather than a consequence of any geometric deformation.

\subsection*{The $k\!=\!3$ degeneracy and its spectral consequences}

Under the condition $V'_{\mathrm{eff}}(r^c) = V''_{\mathrm{eff}}(r^c) = 0$, the Taylor
expansion of $V_{\mathrm{eff}}$ around $r^c$ begins at cubic order:
\[
  V_{\mathrm{eff}}(r) - V_{\mathrm{eff}}(r^c) \sim \tfrac{1}{6}V'''_{\mathrm{eff}}(r^c)(r-r^c)^3.
\]
According to Arnold's classification of singularities~\cite{Arnold72}, this corresponds to the $A_2$ (cusp)
catastrophe, which is the simplest codimension-1 degenerate critical point. After semiclassical
Birkhoff normal form reduction, the effective Hamiltonian is modeled by
$P_{\mathrm{mod}} = -h^2\partial_x^2 + x^3$, a non-self-adjoint operator with a
cusp-type potential.

The spectral consequences of this degeneracy are threefold. First, the spectral gap of the linearized
Hamiltonian flow---defined in the sense of normal hyperbolicity as $\kappa_{\mathrm{gap}}(P) = \sqrt{-2V''_{\mathrm{eff}}(r^c)}$---collapses to zero as $P \to P_c$, thereby violating the assumptions of the Wunsch--Zworski
resolvent estimate. Second, microlocal analysis of $P_{\mathrm{mod}}$ subject to
outgoing boundary conditions on the Stokes contour $\Gamma = \{te^{i2\pi/5}: t > 0\}$
reveals a non-diagonalizable Jordan block of size~3 at the critical point (Theorem~\ref{thm:jordan});
specifically, the geometric multiplicity is 1, whereas the algebraic multiplicity is 3. Third, the WKB
phase structure on $\Gamma$ induces an asymptotic suppression of the matrix element $X_{31}$ of
the geometric perturbation matrix, which prevents the generic Puiseux $\mu^{1/3}$ eigenvalue
splitting and constrains the resolvent growth strictly to the rate $|P-P_c|^{-1/2}$
(Section~6).

\subsection*{Main results}

The primary result of this work (Theorem~\ref{thm:resolvent}) is the two-sided estimate
\[
  \|R_h(P)\| \sim h^{-6/5}|P - P_c|^{-1/2},
\]
valid in the semiclassical limit $h \to 0$ at fixed $\mu = P - P_c \neq 0$.
The factor $h^{-6/5}$ arises from the $A_2$ scaling of the model operator $P_{\mathrm{mod}}$
(where the exponent $6/5$ is determined by the dominant balance between the $h^2\partial_x^2$ and $x^3$
terms), whereas the $|P-P_c|^{-1/2}$ factor accounts for the square-root eigenvalue splitting
of the Schur complement $E_{-+}(\mu)$.
A key ingredient in the proof is the vanishing of the sub-principal symbol,
$Q_1 \equiv 0$ (Appendix~A, Lemma~\ref{lem:Q1}), which ensures that the semiclassical
scaling is governed purely by the principal symbol and that no $O(h)$ correction
alters the dilation exponent $6/5$.

Theorem~\ref{thm:spectral_clustering} (Spectral Clustering) establishes that the degenerate triple resonance at $P_c$
splits into three distinct physical frequencies under the perturbation $\mu = P-P_c$:
one remains at $\omega_c$ to leading order, while the other two bifurcate as
$\omega_{2,3} \approx \omega_c \pm \frac{\sqrt{c_X\mu}}{2\omega_c}$.
In the double scaling regime $|\mu| \sim h^{4/5}$ (Section~7), these two scaling limits couple,
and the resolvent is shown to remain uniformly bounded by $O(h^{-6/5})$, with the corresponding
resonance spacing satisfying $\Delta\omega \sim h^{6/5}$.

\subsection*{Physical interpretation}

Although the critical configuration $P_c$ lies outside the physically regular SdS regime, the results of this paper
carry significant physical implications in two ways.

First, the resolvent estimate $\|R_h(P)\| \sim h^{-6/5}|P-P_c|^{-1/2}$ characterizes
a \emph{threshold phenomenon} with a precise, observable signature. As the SdS parameters
approach the Nariai boundary, the photon sphere shifts toward the degenerate configuration,
causing the spectral gap $\kappa_{\mathrm{gap}}(P)$ to decrease and the resolvent norm to grow.
The rate of this growth, governed by $|P-P_c|^{-1/2}$, constitutes a sharp, computable
geometric invariant that is independent of the specific path taken through parameter space.
In the context of gravitational wave astronomy, this yields a precise and in principle
testable prediction: the QNM decay timescale $\tau = 1/|\mathrm{Im}(\omega)|$ diverges as
$\tau \sim |P - P_c|^{-1/2}$ near the Nariai boundary, so near-extremal SdS black holes
exhibit anomalously long-lived ringdown signals.
The exponent $-1/2$ is rigid---it is fixed by the $A_2$ cusp geometry of the potential
barrier and cannot be altered by sub-leading corrections---and therefore provides a
distinguishing prediction relative to the Kerr--de~Sitter case, where no such
algebraic divergence occurs.
Concretely, if two black holes with mass parameters $r_{s,1}$ and $r_{s,2}$ both lie
within a factor of two of the Nariai boundary
($|P_j - P_c| \leq 2|P_{\mathrm{Nariai}} - P_c|$), then their ringdown
timescales differ by a computable factor $\sqrt{|P_2-P_c|/|P_1-P_c|}$ that is
determined solely by the resolvent estimate of Theorem~\ref{thm:resolvent}.

Second, the spectral clustering established in Theorem~\ref{thm:spectral_clustering} has a
richer structure than the real-axis picture might suggest.
A direct computation gives
$V_{\mathrm{eff}}(r^c;\,r_s^c) = \tfrac{-591+33\sqrt{33}}{136} \approx -2.952 < 0$
(see Remark~\ref{rem:omega_c_imaginary}), so the critical frequency satisfies
$\omega_c^2 < 0$, i.e.\ $\omega_c = i|\omega_c| \in i\mathbb{R}_{>0}$.
The three coalescing resonances therefore lie on the positive imaginary axis at $P_c$,
corresponding to purely decaying modes with no oscillation.
Under the perturbation $\mu = P - P_c$, the first resonance $\omega_1$
remains at $\omega_c$ to leading order (a purely imaginary, non-oscillatory mode), while the
other two split as
\[
  \omega_{2,3} = \omega_c \pm \frac{\sqrt{c_X\,\mu}}{2\omega_c}.
\]
Since $c_X = |c_X|\,e^{i4\pi/5}$ (Lemma~\ref{lem:cx_nonzero}), one computes
$\arg\!\left(\tfrac{\sqrt{c_X\mu}}{2\omega_c}\right) = \tfrac{2\pi}{5} - \tfrac{\pi}{2} = -\tfrac{\pi}{10}$
for real $\mu > 0$.
Consequently $\omega_{2,3}$ acquire nonzero real parts of opposite sign:
\begin{align*}
  \mathrm{Re}(\omega_{2,3}) &= \pm\,|\Delta\omega|\cos(\pi/10), \\
  \mathrm{Im}(\omega_{2,3}) &= |\omega_c| \mp |\Delta\omega|\sin(\pi/10),
\end{align*}
where $|\Delta\omega| = \tfrac{|\sqrt{c_X\mu}|}{2|\omega_c|}$.
The modes $\omega_{2,3}$ are therefore oscillatory (non-zero real part), in contrast
to the purely imaginary mode $\omega_1$.
This \emph{mixed splitting}---one purely imaginary mode and a complex-conjugate-free
oscillatory pair---is a hallmark of the $A_2$ (cusp) degeneracy in a non-self-adjoint
setting, and carries a concrete physical interpretation.

The purely imaginary mode $\omega_1$ represents a perturbation that decays
monotonically with no oscillation; in a near-extremal SdS geometry, it corresponds to
the overdamped sector of the ringdown signal. The oscillatory pair $\omega_{2,3}$, by
contrast, generates a \emph{beating pattern}: two modes with equal decay rates
$\mathrm{Im}(\omega_{2,3}) = |\omega_c| \mp |\Delta\omega|\sin(\pi/10)$ but opposite
oscillation frequencies $\pm|\Delta\omega|\cos(\pi/10)$, which superpose to produce
amplitude modulation at the difference frequency $2|\Delta\omega|\cos(\pi/10)$.
The fixed geometric ratio
\[
  \frac{\mathrm{Re}(\omega_{2,3})}{\mathrm{Im}(\omega_{2,3}) - |\omega_c|}
  = -\cot(\pi/10) \approx -3.08
\]
between the oscillation frequency and the shift in decay rate is a direct invariant
of the $A_2$ singularity, independent of $\mu$ and $h$, and in principle distinguishable
in numerical relativity simulations of near-extremal SdS geometries.

\subsection*{Relation to the literature}

The mathematical study of QNMs on black hole spacetimes has been significantly advanced
by the microlocal framework of Vasy~\cite{Vasy2013}, who established the meromorphic
continuation of the resolvent on asymptotically hyperbolic and Kerr--de~Sitter
spaces, and by Dyatlov~\cite{Dyatlov11}, who proved exponential energy decay
and the QNM expansion for Kerr--de~Sitter perturbations. The Grushin (Schur complement)
reduction employed in this work follows the approach of Dyatlov--Zworski~\cite{DZ19}, whose
formulation of resonances via Fredholm theory and effective Hamiltonians serves as the
technical basis for Sections~5--7 and Appendix~B.

Recent developments in the Kerr--de~Sitter setting that are directly relevant to the
present work include the following.
Hintz~\cite{Hintz2022} established mode stability and a complete asymptotic description
of shallow quasinormal modes for Kerr--de~Sitter black holes throughout the subextremal
range in the limit of small cosmological constant ($\Lambda\mathfrak{m}^2 \searrow 0$).
Petersen--Vasy~\cite{PetersenVasy2023} proved that quasinormal modes for linear wave
equations in subextremal Kerr and Kerr--de~Sitter spacetimes are real analytic,
using a stable radial-point microlocal structure that goes beyond the generalized
normal source/sink framework of~\cite{Vasy2013}.
Fang~\cite{Fang2022} provided an alternative proof of linear stability of the
slowly-rotating Kerr--de~Sitter family via integrated local energy estimates, without
appealing to the meromorphic continuation of the resolvent.
The present paper complements these results by identifying a precise spectral degeneracy
that occurs when SdS parameters approach the Nariai boundary via analytic continuation:
in contrast to the Kerr--de~Sitter case~\cite{Dyatlov11,Dyatlov12}, the SdS resolvent
acquires the additional algebraic factor $|P-P_c|^{-1/2}$ with exponent fixed by the
$A_2$ cusp geometry, and no analogous algebraic divergence arises in the Kerr--de~Sitter
setting under standard parameter variations.

For normally hyperbolic trapped sets, the sharp resolvent estimate $\|R_h\| = O(h^{-1})$
(up to logarithmic losses) was established by Wunsch--Zworski~\cite{WZ11}.
The underlying dynamical description, concerning the exponential decay of correlations for flows with normally
hyperbolic trapping (including null geodesic flows on black hole backgrounds), was
developed by Nonnenmacher--Zworski~\cite{NZ15}.
The present paper provides a precise description of the spectral behavior when this
normal hyperbolicity degenerates: the exponent $h^{-1}$ is replaced by $h^{-6/5}$,
which reflects the cubic rather than quadratic tangency of the potential barrier, and the
resolvent acquires an additional geometric factor $|P-P_c|^{-1/2}$.

The classification of degenerate critical points utilized in Section~4 and Appendix~A
follows Arnold's singularity theory~\cite{Arnold72}. The WKB and Stokes sector analysis in Section~6 and
Appendix~E is based on Sibuya's global theory of the model equation
$-u'' + x^3 u = 0$~\cite{Sibuya75} and Kato's analytic perturbation theory
for the resonance family $Q(\nu) = -\partial_y^2 + y^3 + \nu y$~\cite{Kato66}.

The general theory of trapped sets with higher-codimensional degeneracies has been
studied by Sj\"{o}strand--Zworski~\cite{SZ91}, who analyzed resonances near
degenerate critical levels via microlocal normal forms, and by
Faure--Sj\"{o}strand~\cite{FS11}, who established band structure of Ruelle--Pollicott
resonances for Anosov flows. The present work can be viewed as a codimension-1
instance of this broader program: the $A_2$ (cusp) singularity sits at the boundary
between the normally hyperbolic regime (where~\cite{WZ11,NZ15} apply) and the
fully degenerate regime of~\cite{SZ91}, and the resolvent exponent $h^{-6/5}$
interpolates between the $h^{-1}$ bound of normally hyperbolic trapping and
the stronger growth expected at higher degeneracies.

\subsection*{Structure of the paper}

The remainder of the paper is organized as follows. Section~3 outlines the geometric setup, including
the SdS metric, the Regge--Wheeler effective potential, the Nariai limit, and the critical parameter
configuration $P_c$. Section~4 develops the semiclassical scaling by means of the Birkhoff normal form.
Section~5 implements the Grushin reduction and derives the expansion for the Schur complement.
Section~6 establishes the asymptotic WKB suppression of the matrix element $X_{31}$ and
the non-vanishing of the coefficient $c_X$. Section~7 analyzes the double scaling regime and
derives the corresponding resonance spacing. The appendices contain the sub-principal symbol
calculation ($Q_1 \equiv 0$, Appendix~A), the full determinant and dominant balance
analysis including the Schwarzian contribution $Q^{\mathrm{Schw}}_{31}$
(shown subleading: $|Q^{\mathrm{Schw}}_{31}|/|c_X| \approx 8\times10^{-4} \ll 1$,
Appendix~B), the numerical verification of the Jordan chain structure
(Appendix~C), the construction of the global escape function (Appendix~\ref{app:escape}), and the WKB spectral
theory of the model operator $Q(\nu)$ (Appendix~\ref{app:Qnu}).

\subsection*{Notation}

\leavevmode\vspace{-\baselineskip}

\noindent
For the reader's convenience, we summarize the primary notation
used throughout the paper.

\medskip
\begin{center}
{\small
\renewcommand{\arraystretch}{1.35}
\begin{tabular}{@{}l>{\raggedright\arraybackslash}p{6.4cm}l@{}}
\hline
Symbol & Definition & First appearance \\
\hline
$h$ & semiclassical parameter ($h \to 0$) & \S1 \\
$P = (r_s, L)$ & black hole / cosmological parameter & \S3 \\
$P_c$ & critical degenerate parameter, $V'_{\mathrm{eff}}=V''_{\mathrm{eff}}=0$ & \S3 \\
$\mu$ & scalar deformation parameter along a one-parameter path $P(\mu)$ with $P(0)=P_c$ & \S1 \\
$\nu$ & double-scaling variable $\nu := \mu h^{-4/5}$ & \S7 \\
$r_c(P)$ & photon orbit radius (potential maximum) & \S3 \\
$\omega_c$ & critical resonance frequency at $P_c$ & \S2 \\
$\kappa$ & linearized flow rate, $\kappa := \sqrt{-V''_{\mathrm{eff}}(r_c)}$ & \S1, \S\ref{app:escape} \\
$\kappa_{\mathrm{gap}}(P)$ & spectral gap, $\kappa_{\mathrm{gap}}(P) := \sqrt{2}\,\kappa = \sqrt{-2V''_{\mathrm{eff}}(r_c(P))}$ & \S1, \S\ref{app:escape} \\
$\lambda_j(\mu)$ & eigenvalues of Schur complement $E_{-+}(\mu)$ (distinct from $\kappa_{\mathrm{gap}}(P)$ above) & \S5--7 \\
$\Gamma$ & Stokes contour $\{te^{i2\pi/5} : t > 0\}$ & \S2 \\
$e_1, e_2, e_3$ & Jordan chain basis of $\ker Q_0^{\mathrm{AM}=3}$ & \S2 \\
$X_{jk}$ & matrix element $\langle e_j,\, x\, e_k\rangle_\Gamma$ & \S5 \\
$c_X$ & spectral splitting coefficient $c_X := X_{21} + X_{32}$ & \S5--6 \\
$E_{-+}(\mu)$ & Schur complement (Grushin reduction) & \S5 \\
$P_j(\mu)$ & spectral projector of $E_{-+}(\mu)$ onto $\lambda_j$ & \S5 \\
$Q_0$ & model operator $-\partial_y^2 + y^3$ (rescaled, $h=1$) & \S4 \\
$Q(\nu)$ & perturbed model $-\partial_y^2 + y^3 + \nu y$ & \S7 \\
$Q_1$ & sub-principal symbol (shown $\equiv 0$ in App.~A) & App.~A \\
$R_h(P)$ & semiclassical resolvent $(P_h - \omega^2)^{-1}$ & \S1 \\
$\mu_{\mathrm{cross}}$ & crossover scale $X_{31}^2/c_X^3 \to 0$ & App.~C \\
\hline
\end{tabular}
}
\end{center}
\medskip

\noindent
\textbf{Notational conventions.} Throughout this paper, the relation $A \lesssim B$
denotes $A \leq CB$ for a constant $C$ independent of both $h$ and $\mu$, and $A \sim B$
signifies that $A \lesssim B$ and $B \lesssim A$ simultaneously. All norms $\|\cdot\|$
refer to $L^2 \to L^2$ operator norms unless specified otherwise. To avoid ambiguity, we use distinct symbols for the two spectral quantities appearing
in this paper: $\kappa_{\mathrm{gap}}(P) := \sqrt{-2V''_{\mathrm{eff}}(r_c(P))}$ denotes
the \emph{spectral gap} of the linearized Hamiltonian flow at the photon orbit
(Appendix~\ref{app:escape}), while $\lambda_j(\mu)$ denotes the \emph{eigenvalues}
of the Schur complement $E_{-+}(\mu)$ (Sections~1--7). In earlier versions of the
manuscript these were both written as $\lambda$; the notation has been separated
to reflect their distinct roles.

% -------------------------------------------------------
\section{Main Theorems}
% -------------------------------------------------------

\begin{theorem}[Resolvent Blow-up]
\label{thm:resolvent}
As $P \to P_c$, the semiclassical resolvent exhibits the singular behavior:
\[
  \|R_h(P)\| \sim h^{-6/5}|P - P_c|^{-1/2}.
\]
This blow-up arises from the combination of two independent singularities: the semiclassical singularity
governed by the model dilation operator, which yields the factor $h^{-6/5}$, and the geometric bifurcation
singularity, which yields the factor $|P-P_c|^{-1/2}$ and is shielded from generic fractional scaling by the
asymptotic WKB suppression mechanism ($X_{31} \to 0$ as $T \to \infty$; see Remark~\ref{rem:normalization}).
\end{theorem}

\begin{proof}
The proof factorizes into two independent estimates.

\textbf{Step 1: Semiclassical factor $h^{-6/5}$.}
Following the Birkhoff normal form reduction (Appendix~A), the operator in the vicinity of $P_c$ is modeled by
$P_{\mathrm{mod}} = -h^2\partial_x^2 + x^3$.
Under the dilation $x = h^{2/5}y$, one computes
\[
  P_{\mathrm{mod}} = h^{6/5}(-\partial_y^2 + y^3).
\]
Since $Q_1 \equiv 0$ (Lemma~\ref{lem:Q1}), there is no $O(h)$ correction,
and the rescaled operator $-\partial_y^2 + y^3$ is $h$-independent.
Its resolvent at energy zero satisfies the two-sided bound
\[
  c \;\leq\; \|(-\partial_y^2+y^3)^{-1}\| \;\leq\; C,
\]
for constants $0 < c \leq C < \infty$ independent of $h$.
The lower bound $c \approx 505$ is established in Lemma~\ref{lem:Q0_svd} below
via the Jordan chain relation $Q_0 e_2 = e_1$.
The upper bound $C < \infty$ follows from two independent arguments:
(i)~the coercivity estimate in Lemma~\ref{lem:Q0_svd} (closed graph theorem),
and (ii)~the Grushin reconstruction $\|R_h\| \leq C_{\mathcal{G}}\|E_{-+}^{-1}\|$
combined with the norm estimate for $E_{-+}^{-1}$ in Step~2 below,
which gives $C \leq C_{\mathcal{G}}\cdot\|E_{-+}^{-1}(0,\mu)\|\cdot h^{6/5}$
uniformly in $h$ as $\mu \to 0$ at fixed $\mu \neq 0$.
These two routes are logically independent; (i) establishes finiteness
at $P_c$ directly, while (ii) provides the sharp rate.

\emph{Upper bound.}
The correct functional-analytic setting for $Q_0 := -\partial_y^2 + y^3$ is the
weighted $L^2$ space associated with the outgoing boundary condition on $\Gamma$,
which we denote $L^2_\Gamma$.
Concretely, $L^2_\Gamma$ is the completion of $C_c^\infty(\Gamma)$ under the norm
$\|u\|_{L^2_\Gamma}^2 := \int_\Gamma |u(x)|^2\,|dx|$,
and $Q_0$ is defined on $L^2_\Gamma$ with domain
\begin{equation}\label{eq:H2Gamma_domain}
H^2_\Gamma := \bigl\{u \in L^2_\Gamma : u'',\, y^3 u \in L^2_\Gamma,\;
u \text{ satisfies the outgoing BC at } +\infty,\;
\text{and } \int_{t_L}^{t_L+1} |u(t)|^2\,dt < \infty \bigr\}.
\end{equation}
The last condition is the \emph{lower-end boundary condition}: it requires
$u(t_L) = 0$, which is the natural self-adjoint boundary condition at the
finite endpoint $t_L > 0$.
This is consistent with the Jordan chain construction: by the
variation-of-parameters formula, $e_2(t_L) = e_1(t_L)\,I_1(t_L)/W = 0$
since $I_1(t_L) = \int_{t_L}^{t_L} s\,e_1(s)^2\,ds = 0$,
confirming $e_2 \in H^2_\Gamma$; while $e_1(t_L) \approx -2892 \neq 0$
(Appendix~C.2), so $e_1 \notin H^2_\Gamma$.
On this domain, $Q_0$ is a closed operator with discrete spectrum
(Sibuya~\cite{Sibuya75}, \S5).

We establish $\|Q_0^{-1}\| \leq C$ via a singular value argument.

\begin{lemma}[Resolvent norm bound for $Q_0$]
\label{lem:Q0_svd}
The operator $Q_0 : H^2_\Gamma \to L^2_\Gamma$ is injective with bounded inverse,
and
\[
  \frac{\|e_2\|_{L^2_\Gamma}}{\|e_1\|_{L^2_\Gamma}}
  \;\leq\; \|Q_0^{-1}\|_{L^2_\Gamma \to L^2_\Gamma}
  \;<\; \infty.
\]
Numerically, $\|Q_0^{-1}\| \geq \|e_2\|/\|e_1\| \approx 505$
(algebraic normalization, Table~C.2).
\end{lemma}

\begin{proof}
\emph{Step 1: Injectivity.}
Suppose $Q_0 u = 0$ with $u \in H^2_\Gamma$.
By Lemma~\ref{lem:stokes_gm1}, every outgoing solution of $Q_0 v = 0$
is a scalar multiple of $e_1$.
However, $e_1 \notin H^2_\Gamma$ because it fails the lower-end boundary
condition~\eqref{eq:H2Gamma_domain}: the domain requires $u(t_L) = 0$,
but $e_1(t_L) \approx -2892 \neq 0$ (Appendix~C.2, Table~C.2).

\begin{remark}
The ODE $-v''+t^3 v = 0$ has a regular singular point structure at $t=0$
with indicial roots $\rho = 0$ and $\rho = 1$; both local solutions are bounded
near $t = 0$. The $t^{-3/4}$ behaviour is the large-$t$ WKB amplitude
(Sibuya~\cite{Sibuya75}, \S2), not a small-$t$ singularity.
The injectivity of $Q_0$ therefore rests on the boundary condition $u(t_L)=0$,
not on a Frobenius singularity argument.
\end{remark}

Therefore $u = 0$, and $Q_0$ is injective.

\emph{Step 2: Bounded inverse.}
Since $Q_0$ is a closed Fredholm operator of index~$0$ on the pair
$(H^2_\Gamma, L^2_\Gamma)$ (Sibuya~\cite{Sibuya75}) and is injective,
it is surjective onto $L^2_\Gamma$ by the Fredholm alternative,
and $Q_0^{-1} : L^2_\Gamma \to H^2_\Gamma$ is bounded by the closed graph theorem:
$\|Q_0^{-1}\| < \infty$.

To obtain the explicit finiteness with a concrete bound, we use the
a priori coercivity estimate: for any $u \in H^2_\Gamma$,
\begin{equation}\label{eq:coercive}
  \|Q_0 u\|^2
  = \|u''\|^2 - 2\,\mathrm{Re}\langle u'', t^3 u\rangle + \|t^3 u\|^2
  \;\geq\; \tfrac{1}{2}\|u''\|^2 - C_1\|u\|^2,
\end{equation}
where $C_1 > 0$ depends only on $\Gamma$, and the cross term is bounded via
Cauchy--Schwarz together with the WKB decay $e^{-4t^{5/2}/5}$ on the tail
$t \geq T_0$. This implies $Q_0^* Q_0 + C_1$ is coercive, so $Q_0^* Q_0$
has smallest eigenvalue $\sigma_{\min}^2 > 0$, giving
$\|Q_0^{-1}\| = \sigma_{\min}^{-1} < \infty$.

To make this explicit: by Sibuya~\cite{Sibuya75}, \S5, the operator
$Q_0 = -\partial_t^2 + t^3$ subject to outgoing boundary conditions on $\Gamma$
has discrete spectrum with no accumulation at~$0$.
Concretely, the smallest nonzero eigenvalue satisfies
\begin{equation}\label{eq:sigma_min_Sibuya}
  \sigma_{\min}(Q_0) = \inf_{\substack{u \in H^2_\Gamma \\ \|u\|=1}} \|Q_0 u\|
  \;\geq\; c_{\mathrm{Sib}} \;>\; 0,
\end{equation}
where $c_{\mathrm{Sib}} > 0$ is a constant depending only on the contour
$\Gamma = \{te^{i2\pi/5}: t \geq t_L\}$ and the truncation $t_L = 0.5$,
but independent of $h$.
This follows from the spectral gap of $Q_0$: since $0$ is not an eigenvalue
(established in Step~1 above and in Lemma~\ref{lem:stokes_gm1}) and the
spectrum is discrete, the distance from $0$ to $\mathrm{spec}(Q_0)$ is
positive (Sibuya~\cite{Sibuya75}, \S5, Theorem~5.1).
Therefore $\|Q_0^{-1}\| = \sigma_{\min}(Q_0)^{-1} \leq c_{\mathrm{Sib}}^{-1} < \infty$.

We first verify that $e_2 \in H^2_\Gamma$.
By the variation of parameters formula (Lemma~\ref{lem:jordan_chain}),
$e_2(t) = e_1(t)\,I_1(t)/W$ where $I_1(t) = \int_{t_L}^t s\,e_1(s)^2\,ds$.
Since $e_1$ decays super-exponentially, $e_2 \in L^2_\Gamma$.
Moreover $e_2'' = t^3 e_2 - e^{i4\pi/5}t\,e_1$ (Jordan chain relation),
so $e_2'' \in L^2_\Gamma$ by WKB decay of both terms.
At the lower end: $I_1(t_L) = 0$ so $e_2(t_L) = 0$,
satisfying~\eqref{eq:H2Gamma_domain}. Thus $e_2 \in H^2_\Gamma$.

By the Jordan chain relation $Q_0 e_2 = e_1$,
the element $e_1 \in L^2_\Gamma$ satisfies $Q_0^{-1} e_1 = e_2$,
so the operator norm satisfies
\begin{equation}\label{eq:Q0inv_upper}
  \|Q_0^{-1}\|
  \;\geq\; \frac{\|Q_0^{-1} e_1\|}{\|e_1\|}
  = \frac{\|e_2\|}{\|e_1\|}.
\end{equation}

\begin{remark}[$T$-stability of the ratio $\|e_2\|/\|e_1\|$]
\label{rem:ratio_T_indep}
The ratio $\|e_2\|_{L^2(\Gamma_T)}/\|e_1\|_{L^2(\Gamma_T)}$ converges to a
finite positive limit as $T\to\infty$, since both norms are dominated by the
super-exponential WKB decay and $I_1(t) \to I_1(\infty) < \infty$.
The numerical value $\approx 505$ at $T = 3.2$ is within $0.1\%$ of the
$T = 4.0$ value (Table~C.2), confirming that $T = 3.2$ captures the limit.
The ratio is normalization-independent (Remark~\ref{rem:normalization}).
\end{remark}
\end{proof}

Lemma~\ref{lem:Q0_svd} gives $\|Q_0^{-1}\| < \infty$, establishing the
upper bound for Step~1, and $\|Q_0^{-1}\| \geq \|e_2\|/\|e_1\| \approx 505$,
establishing the lower bound.

\emph{Lower bound.}
By Lemma~\ref{lem:Q0_svd} (equation~\eqref{eq:Q0inv_upper}), the lower bound
$\|Q_0^{-1}\| \geq \|e_2\|/\|e_1\| =: c > 0$
is already established: it follows from $Q_0 e_2 = e_1$ and
$\|Q_0^{-1}\| \geq \|Q_0^{-1} e_1\|/\|e_1\| = \|e_2\|/\|e_1\|$.
Numerically (Appendix~C.2, $T=3.2$, $N=12000$, algebraic normalization):
$\|e_1\| \approx 0.4731$, $\|e_2\| \approx 238.96$,
yielding $c \approx 505$ (Table~C.2 in Appendix~C.4).

Scaling back gives
\[
  \|R_h\|_{L^2\to L^2} = h^{-6/5}\,\|Q_0^{-1}\|.
\]
The lower bound $\|Q_0^{-1}\| \geq c := \|e_2\|/\|e_1\| \approx 505$
is established above via~\eqref{eq:Q0inv_upper}.
The upper bound $\|Q_0^{-1}\| \leq C := c_{\mathrm{Sib}}^{-1} < \infty$
follows from~\eqref{eq:sigma_min_Sibuya} (Sibuya~\cite{Sibuya75}, \S5).
Together:
\[
  c\,h^{-6/5} \;\leq\; \|R_h\| \;\leq\; C\,h^{-6/5},
  \qquad C = c_{\mathrm{Sib}}^{-1} < \infty.
\]
The upper bound $C$ is finite and $h$-independent, depending only on
the spectral gap of $Q_0$ on $\Gamma$.
The lower bound $c \approx 505$ is explicit and normalization-independent
(Remark~\ref{rem:normalization}).
This establishes both bounds for the semiclassical factor in Step~1.
The matching upper bound for $\|R_h\|$ from the Grushin side
(independent confirmation) is derived in Step~2 below.

\textbf{Step 2: Geometric factor $|P-P_c|^{-1/2}$.}
According to the Grushin reduction described in Section~5, the resolvent $R_h(P)$
is bounded if and only if the Schur complement $E_{-+}(\omega,P)$ is invertible, satisfying the estimate
\[
  \|R_h(P)\| \leq C\,\|E_{-+}^{-1}\|
\]
for a constant $C$ uniform in $P$ near $P_c$
(Dyatlov--Zworski~\cite{DZ19}, Appendix~C.1).
Consequently, it suffices to estimate the norm of $E_{-+}^{-1}$.

Recall from Lemma~\ref{lem:Epm_expand} that the Schur complement, viewed as a function
of the spectral parameter $\lambda := \omega^2 - \omega_c^2$ and the geometric
deformation $\mu := P - P_c$, takes the form
\begin{equation}\label{eq:Grushin_inv_corrected}
  E_{-+}(\lambda,\mu) = J_3 - \lambda I + \mu X + O(\lambda^2,\mu^2,\lambda\mu),
\end{equation}
where $J_3$ is the $3\times 3$ nilpotent Jordan block and $X_{jk} = \langle e_j, x\,e_k\rangle$.
The resolvent is reconstructed via the standard Grushin block inversion formula
(Dyatlov--Zworski~\cite{DZ19}, Appendix~C.1):
\[
  R_h(P) = E - E_+ E_{-+}^{-1}(\lambda,\mu)\, E_-,
\]
where $E, E_\pm$ are the remaining blocks of $\mathcal{G}^{-1}$, and
$E_{-+}^{-1}(\lambda,\mu)$ denotes the inverse of the
$3\times 3$ matrix~\eqref{eq:Grushin_inv_corrected}.

\textbf{Resolvent norm via $E_{-+}$.}
By the Grushin framework (Dyatlov--Zworski~\cite{DZ19}, Appendix~C.1),
the resolvent satisfies the two-sided estimate
\[
  \|R_h(P)\| \sim \|E_{-+}^{-1}(\lambda,\mu)\|,
\]
where the implicit constants are uniform in $(\lambda,\mu)$ near $(0,0)$.
This relation holds because the off-diagonal blocks $E_\pm$ of $\mathcal{G}^{-1}$
are bounded uniformly on $\mathcal{V}$, a fact we now justify in detail.

\emph{Boundedness of $E_+$ and $E_-$.}
The auxiliary operators $R_- : \mathbb{C}^3 \to L^2(M)$ and $R_+ : L^2(M) \to \mathbb{C}^3$
are defined as in Section~5 by $R_- v = \sum_{j=1}^3 v_j \chi e_j$ and
$(R_+ u)_j = \langle u, \chi e_j \rangle_{L^2(\Gamma)}$, where $\chi \in C_c^\infty(T^*M)$
is a fixed microlocal cutoff and $\{e_j\}$ is the Jordan basis of
Lemma~\ref{lem:jordan_chain}. Here $\langle\cdot,\cdot\rangle_{L^2(\Gamma)}$ denotes
the $L^2$-inner product on the weighted space associated with the outgoing boundary
condition on $\Gamma$ (a sesquilinear pairing, with conjugation in the second slot),
which is distinct from the bilinear $c$-product used in the matrix element integrals
$X_{jk}$ in Appendix~C. Because $\chi$ has compact support and the elements
$e_j$ exhibit WKB decay on $\Gamma$ (Lemma~\ref{lem:stokes_gm1}), the operators
$R_\pm$ are bounded with norms $\|R_\pm\| \leq C_\chi \sum_j \|e_j\|_{L^2(\mathrm{supp}\,\chi)}$,
where $C_\chi$ depends only on the cutoff $\chi$ and the fixed Jordan chain norms
(see Appendix~C).

\emph{Boundedness of $E$ (the $(1,1)$-block of $\mathcal{G}^{-1}$).}
Write $\mathcal{G}(\mu,h) = \mathcal{G}(0,0)\bigl(I + \mathcal{G}(0,0)^{-1}\delta\mathcal{G}\bigr)$,
where $\delta\mathcal{G} := \mathcal{G}(\mu,h) - \mathcal{G}(0,0)$.
Here $M_0 := \|\mathcal{G}(0,0)^{-1}\| < \infty$ by Lemma~\ref{lem:grushin_invertible},
which establishes bounded invertibility of $\mathcal{G}(0,0)$ at the critical point $P_c$.
For $(\mu,h) \in \mathcal{V}$ (Remark~\ref{rem:grushin_validity}), the thresholds $\mu_0, h_0$
are chosen so that $\|\mathcal{G}(0,0)^{-1}\delta\mathcal{G}\| \leq \frac{1}{2}$,
which ensures convergence of the Neumann series
$\mathcal{G}(\mu,h)^{-1} = \sum_{n=0}^{\infty}(-\mathcal{G}(0,0)^{-1}\delta\mathcal{G})^n
\mathcal{G}(0,0)^{-1}$,
with bound $\|\mathcal{G}(\mu,h)^{-1}\| \leq 2M_0 =: C_0$. Since
$\|E\| \leq \|\mathcal{G}^{-1}\|$, the block $E$ is bounded by $C_0$ uniformly on $\mathcal{V}$,
independently of $\|E_{-+}^{-1}\|$. Combining:
$\|E_\pm\| \leq C_0\|R_\mp\|$, so $E_\pm$ are bounded uniformly on $\mathcal{V}$,
and the growth of $\|R_h\|$ is governed entirely by $\|E_{-+}^{-1}(\lambda,\mu)\|$.

\emph{Determinant of $E_{-+}(\lambda,\mu)$ and resonance structure.}
At leading order in $(\lambda,\mu)$, the determinant of~\eqref{eq:Grushin_inv_corrected} is
\begin{equation}\label{eq:det_Epm}
  \det E_{-+}(\lambda,\mu) = \lambda\bigl(c_X\mu - \lambda^2\bigr) + O(\lambda^2\mu,\mu^2),
\end{equation}
where $c_X = X_{21} + X_{32} \neq 0$ (Lemma~\ref{lem:cx_nonzero}).
This follows from the identity $\det(J_3 - \lambda I + \mu X)
= \lambda(c_X \mu - \lambda^2) + O(\lambda^2\mu, \mu^2)$,
verified by direct expansion using $\det(J_3) = 0$, $\mathrm{Tr}(J_3) = 0$,
and the cofactor structure of $J_3$ (see Appendix~B, equation~\eqref{eq:charpoly}).
Setting the determinant to zero gives the resonance condition
$\lambda_1 = 0$ and $\lambda_{2,3} = \pm\sqrt{c_X\mu} + O(\mu)$,
in agreement with Appendix~B.

\emph{Norm of $E_{-+}^{-1}(\lambda,\mu)$.}
We estimate $\|E_{-+}^{-1}(\lambda,\mu)\|$ for $\omega = \omega_c$ (i.e.\ $\lambda = 0$)
and $\mu \neq 0$ small.
At $\lambda = 0$, equation~\eqref{eq:det_Epm} gives
$\det E_{-+}(0,\mu) = -c_X\mu^2 + O(\mu^3) \neq 0$,
so $E_{-+}(0,\mu)$ is invertible for $\mu \neq 0$ sufficiently small.
By Cramer's rule,
\[
  E_{-+}^{-1}(0,\mu) = \frac{\mathrm{adj}(E_{-+}(0,\mu))}{\det E_{-+}(0,\mu)}.
\]
A direct computation of the adjugate of $E_{-+}(0,\mu) = J_3 + \mu X$ gives
\[
  \mathrm{adj}(J_3 + \mu X)\big|_{\lambda=0}
  = \mathrm{adj}(J_3) + O(\mu)
  = \begin{pmatrix}0&0&1\\0&0&0\\0&0&0\end{pmatrix} + O(\mu),
\]
where $\mathrm{adj}(J_3) \neq 0$ (its $(1,3)$ entry equals $1$, as the $(3,1)$ minor
of $J_3$, obtained by deleting row~3 and column~1, equals
$\det\bigl(\begin{smallmatrix}1&0\\0&1\end{smallmatrix}\bigr) = 1$).
Hence $\|\mathrm{adj}(E_{-+}(0,\mu))\| = O(1)$ as $\mu\to 0$.
Combined with $|\det E_{-+}(0,\mu)| \sim |c_X|\mu^2$, this yields
\begin{equation}\label{eq:Einv_upper}
  \|E_{-+}^{-1}(0,\mu)\| = O(\mu^{-2}).
\end{equation}

\begin{remark}[Apparent discrepancy between $O(\mu^{-2})$ and $O(\mu^{-1/2})$]
\label{rem:discrepancy}
Equation~\eqref{eq:Einv_upper} gives $\|E_{-+}^{-1}(0,\mu)\| = O(\mu^{-2})$,
while the main theorem asserts the sharper rate $O(\mu^{-1/2})$.
There is no contradiction: Cramer's rule provides a \emph{crude} upper bound,
because the adjugate $\mathrm{adj}(E_{-+}(0,\mu)) = O(1)$ and
$|\det E_{-+}(0,\mu)| \sim |c_X|\mu^2$ together give the ratio $O(\mu^{-2})$,
but this over-counts the growth by treating all entries of the adjugate
as equally dominant.
The actual growth is governed by the smallest eigenvalue of $E_{-+}(0,\mu)$:
since $\lambda_1(\mu) = O(\mu)$ and $\lambda_{2,3}(\mu) = \pm\sqrt{c_X\mu} + O(\mu)$,
the spectral distance from the origin satisfies
$\min_j|\lambda_j(\mu)| \sim |\mu|^{1/2}$, which yields the sharp bound
$\|E_{-+}^{-1}(0,\mu)\| \sim |\mu|^{-1/2}$.
In summary: Cramer's rule gives a valid but non-sharp ceiling; the spectral
distance to the nearest resonance determines the actual sharp two-sided rate.
\end{remark}

The sharp bound is established as follows.
The dominant $(1,3)$ entry of $E_{-+}^{-1}(0,\mu)$ is controlled by the
spectral distance from $\lambda = 0$ to the nearest resonance
$\lambda_2(\mu) = \sqrt{c_X\mu}$.
The resolvent norm at $\omega = \omega_c$ (i.e.\ $\lambda = 0$) is bounded above by
\begin{equation}\label{eq:Einv_spectral}
  \|E_{-+}^{-1}(0,\mu)\| \leq \frac{C}{|\lambda_2(\mu)|}
  = \frac{C}{|c_X|^{1/2}|\mu|^{1/2}},
\end{equation}
where the inequality uses the spectral bound
$\|A^{-1}\| \leq (\min_j |\lambda_j|)^{-1}$ applied to the matrix $E_{-+}(0,\mu)$
with eigenvalues $\lambda_j(0,\mu) - 0 = \lambda_j(\mu)$.
More precisely, the eigenvalues of $E_{-+}(0,\mu) = J_3 + \mu X$ at $\lambda = 0$
are exactly the resonance values $\lambda_j(\mu)$ themselves
(since $E_{-+}(\lambda,\mu)v = 0 \Leftrightarrow (J_3 - \lambda I + \mu X)v = 0 \Leftrightarrow (J_3 + \mu X)v = \lambda v$),
and $\min_j|\lambda_j(\mu)| = |\lambda_{2}(\mu)| \sim |c_X|^{1/2}|\mu|^{1/2}$
(the eigenvalue $\lambda_1 = O(\mu)$ is of higher order than $|\lambda_{2,3}| \sim \mu^{1/2}$
for small $\mu$, so the minimum is achieved at $j = 2$ or $3$).

\emph{Lower bound.}
The lower bound $\|E_{-+}^{-1}(0,\mu)\| \geq c|\mu|^{-1/2}$ follows from the same
spectral distance argument: taking the unit vector $v_2$ in the direction of the eigenvector
of $E_{-+}(0,\mu)$ corresponding to $\lambda_2(\mu)$,
\[
  \|E_{-+}^{-1}(0,\mu)\| \geq \|E_{-+}^{-1}(0,\mu)\,v_2\|
  = |\lambda_2(\mu)|^{-1} \sim |c_X|^{-1/2}|\mu|^{-1/2}.
\]
Combining with~\eqref{eq:Einv_spectral}, we conclude:
\[
  \|E_{-+}^{-1}(0,\mu)\| \sim |c_X|^{-1/2}\,|P - P_c|^{-1/2}.
\]
Combining with the two-sided estimate $\|R_h(P)\| \sim \|E_{-+}^{-1}\|$,
verified explicitly below (equations~\eqref{eq:lower_bound_Einvplus}--\eqref{eq:twosided_Rh}), we obtain
\[
  \|R_h(P)\| \sim |P-P_c|^{-1/2}.
\]

\emph{Two-sided estimate $\|E_{-+}^{-1}\|\sim\|R_h\|$, with explicit uniform constants.}

We use two block-multiplication identities read off from
$\mathcal{G}(\mu,h)\,\mathcal{G}(\mu,h)^{-1} = I$
(Dyatlov--Zworski~\cite{DZ19}, Appendix~C.1):
\begin{equation}\label{eq:block_ids}
  R_+\,E_+ = I_3 \quad \text{(block $(2,2)$)},
  \qquad
  E_-\,R_- = I_3 \quad \text{(block $(2,1)$ of $\mathcal{E}\mathcal{P}=I$)}.
\end{equation}

\emph{Upper bound $\|R_h\|\leq C\|E_{-+}^{-1}\|$.}
From $R_h = E - E_+E_{-+}^{-1}E_-$, the triangle inequality and submultiplicativity give
\[
  \|R_h\| \leq \|E\| + \|E_+\|\,\|E_{-+}^{-1}\|\,\|E_-\|
  \leq C_0 + C_0^2\|R_-\|\|R_+\|\,\|E_{-+}^{-1}\|
  \leq C\,\|E_{-+}^{-1}\|,
\]
where $C := 1 + C_0^2\|R_-\|\|R_+\|$ and the last step absorbs the constant
$C_0$ into a multiple of $\|E_{-+}^{-1}\|$ (valid since
$\|E_{-+}^{-1}\|\to\infty$ as $\mu\to 0$).

\emph{Lower bound $\|E_{-+}^{-1}\|\leq C'\|R_h\|$, with explicit $C'$.}
From $E_-R_- = I_3$ in~\eqref{eq:block_ids}, the operator $R_-:\mathbb{C}^3\to L^2(M)$
is a right-inverse of $E_-$.
Taking norms: $1 = \|E_-R_-\| \leq \|E_-\|\,\|R_-\|$, so
\begin{equation}\label{eq:sigma_min_Eminus}
  \sigma_{\min}(E_-) \;\geq\; \|R_-\|^{-1} \;>\; 0,
\end{equation}
uniformly in $(\mu,h)\in\mathcal{V}$ since $\|R_-\|$ is fixed.
Now apply $R_+$ on the left to $R_h = E - E_+E_{-+}^{-1}E_-$,
using $R_+E = 0$ (block $(2,1)$ of $\mathcal{P}\mathcal{E}=I$) and
$R_+E_+ = I_3$ from~\eqref{eq:block_ids}:
\begin{equation}\label{eq:Einvplus_extract}
  E_{-+}^{-1}\,E_- \;=\; -R_+\,R_h.
\end{equation}
Taking norms and applying~\eqref{eq:sigma_min_Eminus}:
\[
  \|E_{-+}^{-1}\|\,\|R_-\|^{-1}
  \;\leq\;
  \|E_{-+}^{-1}\|\,\sigma_{\min}(E_-)
  \;\leq\;
  \|E_{-+}^{-1}E_-\|
  \;=\;
  \|R_+R_h\|
  \;\leq\;
  \|R_+\|\,\|R_h\|,
\]
hence
\begin{equation}\label{eq:lower_bound_Einvplus}
  \|E_{-+}^{-1}\| \;\leq\; C'\,\|R_h\|,
  \qquad C' \;:=\; \|R_+\|\,\|R_-\|.
\end{equation}
The constant $C'$ depends only on $\|R_+\|$ and $\|R_-\|$, both of which
are determined by the Jordan chain norms and the cutoff $\chi$
(see the \emph{Boundedness of $E_+$ and $E_-$} paragraph above)
and are independent of $\mu$ and $h$.
Combining with the upper bound, we conclude
\begin{equation}\label{eq:twosided_Rh}
  \frac{1}{C'}\,\|E_{-+}^{-1}\| \;\leq\; \|R_h\| \;\leq\; C\,\|E_{-+}^{-1}\|,
\end{equation}
with explicit constants $C = 1+C_0^2\|R_-\|\|R_+\|$ and $C' = \|R_+\|\|R_-\|$,
both depending only on $C_0 = 2M_0$ and $\|R_\pm\|$, hence uniform on~$\mathcal{V}$.

\begin{remark}[Grushin upper bound confirms Step~1]
\label{rem:grushin_upper_confirms}
Combined with the dilation $\|R_h\| = h^{-6/5}\|Q_0^{-1}\|$ from Step~1,
the upper bound~\eqref{eq:twosided_Rh} gives
\[
  \|Q_0^{-1}\| \;\leq\; C\,h^{6/5}\,\|E_{-+}^{-1}(0,\mu)\|
  \;\leq\; \frac{C\cdot C_{\mathrm{Sib}}^{-1/2}}{|c_X|^{1/2}|\mu|^{1/2}}\cdot h^{6/5}
\]
uniformly for $\mu \neq 0$ fixed and $h \to 0$.
Taking $\mu = \mu_0$ fixed, this is bounded independently of $h$, confirming
$\|Q_0^{-1}\| \leq C_{\mathrm{Sib}}^{-1} < \infty$ from a second route.
The two routes — spectral gap via Sibuya (direct, at $P_c$) and Grushin
reconstruction (indirect, at $\mu \neq 0$) — are logically independent
and mutually consistent.
\end{remark}
\textbf{Step 3: Synthesis of the estimates.}

\emph{Logical prerequisite: WKB suppression.}
The eigenvalue estimate $|\lambda_{2,3}|\sim|\mu|^{1/2}$ employed in Step~2 is
contingent upon the condition $X_{31}=0$.
This condition is established by Lemma~\ref{lem:X31_asymptotic}
(Section~\ref{sec:wkb_suppression}), whose proof depends only on:
(a)~the existence and outgoing boundary property of the Jordan chain $\{e_1,e_3\}$
(Lemma~\ref{lem:jordan_chain}), and
(b)~the WKB decay estimate $|e_j(t)|\lesssim t^{-3/4}e^{-2t^{5/2}/5}$
on $\Gamma$ (Lemma~\ref{lem:stokes_gm1}).
Neither (a) nor (b) depends on the resolvent estimates of Steps~1--2 or on
Theorem~\ref{thm:resolvent} itself.
The absence of circular dependence is verified explicitly in
Proposition~\ref{prop:wkb_independence}.
The complete logical order is therefore:
\begin{gather*}
  \underbrace{\text{Jordan chain}}_{\text{Lem.~\ref{lem:jordan_chain}}}
  \;\Rightarrow\;
  \underbrace{X_{31}=0}_{\text{Lem.~\ref{lem:X31_asymptotic}}}
  \;\Rightarrow\;
  \underbrace{\lambda_{2,3}\sim\mu^{1/2}}_{\text{Step~2}}
  \;\Rightarrow\;
  \underbrace{\|R_h\|\sim h^{-6/5}|P-P_c|^{-1/2}}_{\text{Steps~1--3}}.
\end{gather*}
With $X_{31}=0$, the dominant eigenvalues of $E_{-+}$ satisfy
$\lambda_{2,3} \sim \mu^{1/2}$ rather than the generic Puiseux scaling
$\mu^{1/3}$ that would arise if $X_{31}\neq 0$, thereby locking
$\|E_{-+}^{-1}\| \sim |\mu|^{-1/2}$ rather than $|\mu|^{-2/3}$.

\emph{Independence of the scaling bounds.}
In Step~1, the bound $\|R_h\| \leq Ch^{-6/5}$ is derived by conjugating
$P_h$ to the model operator $h^{6/5}(-\partial_y^2+y^3)$ via the dilation $x=h^{2/5}y$.
This conjugation is determined solely by $h$ and is independent of the parameter $P$: the Birkhoff normal form
reduction (Appendix~A) is performed at the fixed critical point $(r^c, \xi^c)$
associated with $P_c$, rendering the model operator independent of $P$.
Thus, the constant $C$ in Step~1 is uniform for all $P$ in a fixed
neighborhood of $P_c$.
Conversely, in Step~2, the Grushin reduction and the eigenvalue estimate
$|\lambda_{2,3}|\sim|c_X|^{1/2}|\mu|^{1/2}$ for $E_{-+}(\mu)$ are established
at the level of the principal symbol. The matrix elements $X_{ij}$
and the resulting eigenvalue splitting depend on $P$ but are independent of $h$
at leading order (the quantum correction $h^2 Q$ is subleading, as shown in Appendix~B).
Hence, the constant in the Step~2 bound is uniform in $h$ for $h\in(0,1]$.
Since the bounds are uniform in their respective parameters, they can be combined
to yield the joint estimate:
\[
  \|R_h(P)\| \sim h^{-6/5}\,|P - P_c|^{-1/2}.
\]
\end{proof}

\begin{remark}[Regime of validity]
\label{rem:validity}
Theorem~\ref{thm:resolvent} characterizes the resolvent growth at a fixed
deformation $\mu = P - P_c$ in the semiclassical limit $h \to 0$.
In the double scaling regime $|\mu| \sim h^{4/5}$, these limits are coupled,
and the resolvent remains uniformly bounded by $O(h^{-6/5})$, with no stronger
divergence occurring in this window (see Section~\ref{subsec:resolvent_bound_ds}).
Consequently, the relation $\|R_h\| \sim h^{-6/5}|P-P_c|^{-1/2}$ should not
be evaluated by substituting $|P-P_c| \sim h^{4/5}$, as this regime
lies outside the validity of the fixed-$\mu$ estimate.
\end{remark}

\begin{theorem}[Spectral Clustering]
\label{thm:spectral_clustering}
At the critical configuration $P = P_c$, the multiple resonances coalesce. Under a
parameter perturbation $\mu = P - P_c$, the non-diagonalizable Jordan block splits
into three distinct physical frequencies:
\[
  \lambda_1 = O(\mu) \implies \omega_1 = \omega_c + O(\mu/\omega_c)
  \qquad
  \lambda_{2,3} = \pm\sqrt{c_X\mu} + O(\mu) \implies \omega_{2,3} \approx \omega_c \pm \frac{\sqrt{c_X\mu}}{2\omega_c}
\]
\end{theorem}

\begin{proof}
The proof is structured in three steps.

\textbf{Step 1: Characteristic equation.}
As derived in Appendix~B (see equation~\eqref{eq:charpoly}), the Schur complement matrix admits the expansion
\[
  E_{-+}(\mu,h) = J_3 + \mu X + h^2 Q + O(\mu^2, \mu h^2, h^4),
\]
with characteristic polynomial
\[
  \det E_{-+}(\lambda,\mu,h)
  = \det(J_3 - \lambda I + \mu X + h^2 Q)
  = -\lambda^3 + \mu\,\mathrm{Tr}(X)\,\lambda^2
    + \mu\,c_X\,\lambda
    + \mu X_{31} + h^2 Q_{31} + O(\lambda^2\mu,\mu^2),
\]
where $c_X = X_{21}+X_{32}$ and $\mathrm{Tr}(X) = X_{11}+X_{22}+X_{33}$.
By the WKB suppression mechanism of Section~6, $X_{31} = 0$.
The quantum error satisfies $h^2 Q_{31} = O(h^2)$, which is sub-leading
relative to $O(h^{6/5})$ for $h \in (0,1]$ (Appendix~B, \S B.2--B.4).
Neglecting subleading terms, the characteristic equation reduces to
\begin{equation}\label{eq:reduced}
  \lambda^3 - \mu\,\mathrm{Tr}(X)\,\lambda^2 - \mu\,c_X\,\lambda = 0.
\end{equation}

\begin{remark}[The $\mathrm{Tr}(X)$ correction does not affect the resolvent estimate]
\label{rem:trX_resolvent}
Although the trace term $\mu\,\mathrm{Tr}(X)\,\lambda^2$ is retained in the full
characteristic polynomial~\eqref{eq:reduced}, it does \emph{not} alter the resolvent
estimate $\|R_h(P)\| \sim h^{-6/5}|P-P_c|^{-1/2}$ of Theorem~\ref{thm:resolvent}.
The reason is a simple dominant balance: under the ansatz $\lambda \sim \mu^{1/2}$,
the trace term is of order $O(\mu^2)$, whereas the driving term $\mu\,c_X\,\lambda$
is of order $O(\mu^{3/2})$.
Hence $|\mu\,\mathrm{Tr}(X)\,\lambda| \ll |c_X|$ for all sufficiently small $|\mu|$,
regardless of the size of $|\mathrm{Tr}(X)|/|c_X|$.
The quantitative bound $|\mathrm{Tr}(X)|/|c_X| = O(1) \approx 3.1$ (Appendix~B,
Remark~\ref{rem:TrX_bound}) confirms that no accidental enhancement can rescue the
$O(\mu^2)$ term.
The leading eigenvalues thus remain $\lambda_{2,3} = \pm\sqrt{c_X\mu} + O(\mu)$, and
the resolvent growth $\|E_{-+}^{-1}\| \sim |\mu|^{-1/2}$ is unaffected.
The full argument is detailed in Appendix~B.4 (Proposition~\ref{prop:dominant_balance}
and Remark~\ref{rem:TrX_bound}).
\end{remark}

\textbf{Step 2: Dominant balance and eigenvalue splitting.}
Factoring equation~\eqref{eq:reduced} yields:
\[
  \lambda\bigl(\lambda^2 - \mu\,\mathrm{Tr}(X)\,\lambda - \mu\,c_X\bigr) = 0.
\]
This immediately implies $\lambda_1 = 0$ to leading order (the subleading correction $\lambda_1 = O(\mu)$ is derived in the remark following~\eqref{eq:lambdapm} in Appendix~B and does not alter the dominant balance).
For the remaining factor, under the ansatz $\lambda \sim \mu^{1/2}$,
the dominant terms satisfy $\lambda^2 \sim \mu\,c_X \sim \mu$,
while the correction term associated with $\mathrm{Tr}(X)$ satisfies $\mu\,\mathrm{Tr}(X)\,\lambda \sim \mu^{3/2}$,
which is of higher order. The quadratic equation thus simplifies to
\[
  \lambda^2 = \mu\,c_X + O(\mu^{3/2}),
\]
yielding
\[
  \lambda_{2,3} = \pm\sqrt{c_X\,\mu} + O(\mu).
\]
The alternative Puiseux scaling $\lambda \sim \mu^{1/3}$ is inconsistent: if $\lambda \sim \mu^{1/3}$,
then $\lambda^3 \sim \mu$ whereas $\mu c_X \lambda \sim \mu^{4/3} \ll \mu$,
making a balance impossible (see Appendix~B, Proposition~\ref{prop:dominant_balance}).
The unique consistent leading-order eigenvalues are therefore
\[
  \lambda_1 = 0, \qquad \lambda_{2,3} = \pm\sqrt{c_X\,\mu} + O(\mu).
\]

\textbf{Step 3: Determination of the physical frequencies.}
Quantum resonances correspond to values $\omega$ such that $\omega^2 \in \mathrm{spec}(P_h)$.
In the vicinity of the critical frequency $\omega_c$, we write $\omega = \omega_c + \Delta\omega$
with $|\Delta\omega| \ll |\omega_c|$ (note: since $\omega_c = i|\omega_c| \in i\mathbb{R}$ by
Remark~\ref{rem:omega_c_imaginary}, the condition $|\Delta\omega| \ll |\omega_c|$ refers to
moduli in $\mathbb{C}$), which yields
\[
  \omega^2 - \omega_c^2
  = (\omega - \omega_c)(\omega + \omega_c)
  \approx 2\omega_c\,\Delta\omega.
\]
The linearization $\omega^2 - \omega_c^2 \approx 2\omega_c\Delta\omega$ is valid
for $\omega_c \in i\mathbb{R}$ since it relies only on the factorization above
and $|\Delta\omega| \ll |\omega_c|$.
\emph{Quantitative validity.}
Since $\omega_c = i|\omega_c|$ with $|\omega_c| \approx 1.718\,L^{-1}$,
the neglected quadratic term satisfies
$|(\Delta\omega)^2| \leq \delta^2|\omega_c|^2$
whenever $|\Delta\omega| \leq \delta|\omega_c|$.
For the splitting $|\Delta\omega| = |c_X|^{1/2}|\mu|^{1/2}/(2|\omega_c|)$,
this yields the condition
\[
  |\mu| \;\ll\; 4|\omega_c|^4/|c_X| =: \mu_{\mathrm{lin}},
\]
so the linearization is valid uniformly for $|\mu| \leq c\,\mu_{\mathrm{lin}}$
with any fixed $c \in (0,1)$.
Numerically, $\mu_{\mathrm{lin}} \approx 4(1.718)^4/|c_X| \approx 34.8\,L^{-4}/|c_X|$,
which is large compared to the double-scaling threshold $\mu \sim h^{4/5}$ for
all $h \in (0,1]$ in practice.
The linearization therefore breaks down only at parameter distances
$|\mu| = O(\mu_{\mathrm{lin}})$, which lies well outside the domain of validity
of the Grushin expansion (Remark~\ref{rem:grushin_validity}).

The eigenvalue $\lambda$ of the Schur complement $E_{-+}$ relates to the frequency deviation
$\omega^2 - \omega_c^2$ via Lemma~\ref{lem:Epm_expand}:
\[
  \lambda = \omega^2 - \omega_c^2 + O(\lambda^2, \mu^2).
\]
Writing $\omega = \omega_c + \Delta\omega$ with $|\Delta\omega|\ll|\omega_c|$:
\[
  \lambda \approx \omega^2 - \omega_c^2 = 2\omega_c\,\Delta\omega + O(\Delta\omega^2).
\]
For the eigenvalue $\lambda_1 = 0$, we have $\Delta\omega_1 = 0$, implying $\omega_1 = \omega_c$
to leading order (with an $O(\mu/\omega_c)$ correction arising from the subleading term of $\lambda_1$).

For the eigenvalues $\lambda_{2,3} = \pm\sqrt{c_X\mu}$, we find
\[
  \Delta\omega_{2,3} = \frac{\lambda_{2,3}}{2\omega_c}
  = \pm\frac{\sqrt{c_X\mu}}{2\omega_c},
\]
yielding
\[
  \omega_{2,3} = \omega_c \pm \frac{\sqrt{c_X\,\mu}}{2\omega_c} + O(\mu).
\]
Since $c_X \neq 0$ (Lemma~\ref{lem:cx_nonzero})
and $\mu = P - P_c \neq 0$, the three resonance frequencies are distinct for all
P sufficiently close, but not equal, to $P_c$.
\end{proof}

\begin{remark}[The critical frequency $\omega_c$ is purely imaginary]
\label{rem:omega_c_imaginary}
At the critical configuration $P_c$, the effective potential evaluated at the
degenerate critical point $r^c \approx 0.9109\,L$ with $r_s^c \approx 1.1192\,L$ satisfies
\[
  V_{\mathrm{eff}}(r^c;\,r_s^c)
  \;=\; f(r^c;\,r_s^c)\!\left[\frac{6}{(r^c)^2} - \frac{3\,r_s^c}{(r^c)^3}\right]
  \;=\; \frac{-591 + 33\sqrt{33}}{136}
  \;\approx\; -2.952,
\]
which is strictly negative.

\medskip
\noindent\textbf{Derivation of the closed form.}
For this derivation we set $L = 1$ (all quantities below are in units of $L$;
restore $L$ by dimensional analysis: $r \to r/L$, $r_s \to r_s/L$,
$V_{\mathrm{eff}} \to L^2 V_{\mathrm{eff}}$). The exact critical values are
\[
  r_s^c = \frac{\sqrt{54+22\sqrt{33}}}{12}, \qquad
  r^c   = \frac{-5\sqrt{2}+3\sqrt{66}}{2\sqrt{27+11\sqrt{33}}}.
\]
One verifies that these satisfy the degeneracy conditions
$V'_{\mathrm{eff}}(r^c;r_s^c)=0$ and $V''_{\mathrm{eff}}(r^c;r_s^c)=0$, which
yield the two auxiliary relations
\begin{equation}\label{eq:rc_relations}
  \frac{r_s^c}{(r^c)^3} = \frac{6}{(r^c)^2} - \frac{12}{(r^c)^2}
  \quad\Longleftrightarrow\quad
  r_s^c = 2r^c - \frac{2(r^c)^3}{L^2},
  \qquad
  \frac{1}{(r^c)^2} = \frac{3r_s^c}{2(r^c)^3}.
\end{equation}
Substituting the first relation into the lapse function
$f(r^c;r_s^c) = 1 - r_s^c/r^c - (r^c)^2/L^2$ and into the bracket
$6/(r^c)^2 - 3r_s^c/(r^c)^3$, and then using the second relation to
express every occurrence of $(r^c)^{-2}$ in terms of
$r_s^c\,(r^c)^{-3}$, one reduces the product to a rational function of
$\sqrt{33}$ alone (since $(r_s^c)^2$ and $(r^c)^2$ are rational in $\sqrt{33}$).
A direct CAS computation (or by hand, collecting powers of $\sqrt{33}$) then
yields
\[
  V_{\mathrm{eff}}(r^c;r_s^c)
  = \frac{-591 + 33\sqrt{33}}{136}.
\]
Numerically, $33\sqrt{33} \approx 189.59$, giving
$(-591 + 189.59)/136 \approx -2.952 < 0$, confirming the sign.
The bracket and lapse evaluate to
\[
  \tfrac{6}{(r^c)^2} - \tfrac{3r_s^c}{(r^c)^3} \approx 2.789 > 0,
  \qquad f(r^c;\,r_s^c) \approx -1.058 < 0.
\]

\medskip
\noindent\textbf{Why $f(r^c;\,r_s^c) < 0$.}
In the standard SdS geometry with $r_s < r_s^{\mathrm{Nariai}}$, the lapse
$f(r; r_s) = 1 - r_s/r - r^2/L^2$ is positive between the two horizons and
negative outside them. However, the critical parameter $r_s^c \approx 1.1192\,L$
satisfies $r_s^c > r_s^{\mathrm{Nariai}} \approx 0.3849\,L$
(Lemma~\ref{lem:physical_shielding}), placing $P_c$ strictly outside
the physically regular regime. For such values of $r_s$, the cubic equation
$f(r; r_s) = 0$, equivalently $r^3 - r + r_s = 0$ (in units $L = 1$),
has \emph{no real positive roots}: the discriminant condition for real roots
requires $r_s \leq r_s^{\mathrm{Nariai}}$, which is violated at $r_s^c$.
Therefore, $f(r;\,r_s^c) < 0$ for \emph{all} $r > 0$, and in particular
$f(r^c;\,r_s^c) < 0$.

This negativity is not a geometric anomaly but a direct consequence of
the analytic continuation: the critical configuration $P_c$ is accessed
by continuing $r_s$ past the Nariai boundary, where no physical metric exists.
The function $V_{\mathrm{eff}}(r;\,r_s^c)$ and the operator $P_h(r_s^c)$
remain well-defined as analytic expressions in $r$ and $r_s$
(Lemma~\ref{lem:meromorphic_continuation}, Claim~(i)), but the associated
``geometry'' carries no direct physical interpretation as a black hole spacetime.
The sign $f < 0$ is an algebraic property of the continued lapse function
at this parameter value, analogous to the way a polynomial evaluated outside
its domain of definition may take negative values without implying any
physical singularity.

\medskip
\noindent\textbf{Consequences for $\omega_c$.}
Since $f(r^c;\,r_s^c) < 0$ and the bracket is positive, we obtain
$V_{\mathrm{eff}}(r^c;\,r_s^c) < 0$, which forces
\[
  \omega_c^2 = V_{\mathrm{eff}}(r^c;\,r_s^c) < 0,
  \qquad
  \omega_c = i\,|\omega_c|, \quad |\omega_c| = \sqrt{\tfrac{591 - 33\sqrt{33}}{136}} \approx 1.718\,L^{-1}.
\]
\textbf{Sign convention.}
Since $\omega_c^2 < 0$, the square root has two branches: $\omega_c = \pm i|\omega_c|$.
Throughout this paper we adopt the branch $\omega_c = +i|\omega_c|$, i.e.\ the
positive imaginary root.
This is a matter of algebraic convention: the resolvent $R_h(P) = (P_h - \omega^2)^{-1}$
depends only on $\omega^2 = \omega_c^2 < 0$, so both branches $\pm\omega_c$ are
spectral parameters for the same operator.
The choice $\omega_c = +i|\omega_c|$ is consistent with the meromorphic continuation
framework of Lemma~\ref{lem:meromorphic_continuation}: the continued resonance family
at $P_c$ arrives at $+i|\omega_c|$ when $r_s$ is continued past $r_s^{\mathrm{Nariai}}$,
as verified by the explicit computation of $V_{\mathrm{eff}}(r^c;\,r_s^c)$ above.
The sign does not affect any of the resolvent estimates, since these depend on
$\omega_c^2$, not on $\omega_c$ itself.
The three resonances coalescing at $P_c$ therefore lie on the positive imaginary axis.
This conclusion is an intrinsic property of the meromorphically continued resolvent
family at $P_c$, not a statement about any physical SdS spacetime.

This has a concrete effect on the splitting formula of Theorem~\ref{thm:spectral_clustering}.
Since $c_X = |c_X|\,e^{i4\pi/5}$ (Lemma~\ref{lem:cx_nonzero}) and $\omega_c = i|\omega_c|$,
the frequency deviation for real $\mu > 0$ satisfies
\[
  \frac{\sqrt{c_X\,\mu}}{2\omega_c}
  = \frac{|\sqrt{c_X\mu}|\,e^{i2\pi/5}}{2i|\omega_c|}
  = \frac{|\sqrt{c_X\mu}|}{2|\omega_c|}\,e^{-i\pi/10},
\]
so $\arg(\Delta\omega_{2,3}) = -\pi/10 = {-18}^\circ$.
In particular $\Delta\omega_{2,3}$ is neither real nor purely imaginary:
\[
  \mathrm{Re}(\omega_{2,3}) = \pm\,|\Delta\omega|\cos(\pi/10),
  \qquad
  \mathrm{Im}(\omega_{2,3}) = |\omega_c| \mp |\Delta\omega|\sin(\pi/10),
\]
where $|\Delta\omega| := \tfrac{|\sqrt{c_X\mu}|}{2|\omega_c|}$.
Thus $\omega_1 = \omega_c \in i\mathbb{R}$ is a \emph{purely imaginary} (non-oscillatory) mode, while
$\omega_{2,3}$ are \emph{oscillatory}: they acquire nonzero real parts of opposite sign,
producing a genuine doublet of damped oscillations. The resonance spacing
$|\omega_2 - \omega_3| = 2|\Delta\omega| \sim h^{6/5}$ (double scaling regime) measures
the diameter of this doublet in the complex plane.
\end{remark}

\begin{lemma}[Geometric Multiplicity via Stokes Analysis]
\label{lem:stokes_gm1}
The kernel of $P_{\mathrm{mod}} = -\partial_x^2 + x^3$ subject to the purely outgoing
boundary condition on $\Gamma = \{te^{i2\pi/5} : t \in \mathbb{R}_{>0}\}$ is
one-dimensional:
\[
  \ker(P_{\mathrm{mod}}) \cap \{\text{outgoing}\} = \mathrm{span}\{e_1\},
  \qquad \mathrm{GM} = 1.
\]
\end{lemma}

\begin{proof}
The equation $-u'' + x^3 u = 0$ on $\Gamma$, under the coordinate parametrization $x = te^{i2\pi/5}$,
becomes $u''(t) = t^3 u(t)$ for $t \in \mathbb{R}_{>0}$.
This is a second-order linear differential equation with a two-dimensional space of solutions.
The asymptotic WKB solutions of $u'' = t^3 u$ behave as
\[
  v_\pm(t) \sim t^{-3/4}\exp\!\Bigl(\pm\tfrac{2}{5}t^{5/2}\Bigr)
  \qquad (t \to +\infty),
\]
so $v_+(t) \to 0$ (characterizing the decaying, outgoing mode), while $v_-(t) \to \infty$
(representing the growing mode).

To verify that $v_+$ corresponds to the outgoing mode on the original contour
$\Gamma = \{te^{i2\pi/5}: t > 0\}$, we evaluate the WKB phase of $-u'' + x^3 u = 0$
at $x = te^{i2\pi/5}$:
\[
  S(x)\big|_\Gamma
  = \tfrac{2}{5}x^{5/2}\big|_\Gamma
  = \tfrac{2}{5}t^{5/2}e^{i\pi}
  = -\tfrac{2}{5}t^{5/2}.
\]
Since $S(x)|_\Gamma$ is real and negative, the WKB modes of the transformed equation are
\[
  v_+(t) \sim t^{-3/4}e^{+S|_\Gamma} = t^{-3/4}e^{-2t^{5/2}/5} \;\to 0
  \qquad\text{(decaying, outgoing)},
\]
\[
  v_-(t) \sim t^{-3/4}e^{-S|_\Gamma} = t^{-3/4}e^{+2t^{5/2}/5} \;\to \infty
  \qquad\text{(growing, non-outgoing)}.
\]
The decaying mode is given by $e^{+S(x)}|_\Gamma$ (rather than $e^{-S(x)}|_\Gamma$) because
$S|_\Gamma < 0$.
The outgoing boundary condition thus uniquely selects $e_1 := v_+$, establishing that $\mathrm{GM} = 1$.
\end{proof}

\begin{lemma}[Linear independence from Jordan chain]
\label{lem:jordan_indep}
If $P_{\mathrm{mod}}\,e_1 = 0$, $P_{\mathrm{mod}}\,e_2 = e_1$,
$P_{\mathrm{mod}}\,e_3 = e_2$, and $e_1 \neq 0$, then $\{e_1,e_2,e_3\}$
is linearly independent.
\end{lemma}

\begin{proof}
Suppose that $\alpha_1 e_1 + \alpha_2 e_2 + \alpha_3 e_3 = 0$.
Applying $P_{\mathrm{mod}}$ to this relation yields $\alpha_2 e_1 + \alpha_3 e_2 = 0$.
Applying $P_{\mathrm{mod}}$ once more, we obtain $\alpha_3 e_1 = 0$.
Since $e_1 \neq 0$, it follows that $\alpha_3 = 0$.
Substituting this back, $\alpha_2 e_1 = 0$ yields $\alpha_2 = 0$,
and then $\alpha_1 e_1 = 0$ implies $\alpha_1 = 0$.
\end{proof}

\begin{lemma}[Jordan Chain, Normalization, and Algebraic Multiplicity]
\label{lem:jordan_chain}
There exists a set of functions $\{e_1, e_2, e_3\}$ satisfying the Jordan chain relations
\[
  P_{\mathrm{mod}}\,e_1 = 0, \qquad
  P_{\mathrm{mod}}\,e_2 = e_1, \qquad
  P_{\mathrm{mod}}\,e_3 = e_2,
\]
each satisfying the outgoing boundary condition on $\Gamma$, and linearly independent.
Moreover, the basis can be chosen (by adjusting each $e_k$ by a multiple of $e_1$)
to satisfy the \emph{normalization condition}:
\begin{equation}\label{eq:normalization}
  \langle e_{k-1},\, e_j \rangle_{L^2(\Gamma)} = \delta_{j,\,k-1}
  \quad\text{for } k = 2,3,\; j = 1,2,3.
\end{equation}
Condition~\eqref{eq:normalization} is required for the Schur complement to take the
Jordan form $E_{-+}(0,0) = J_3$ (see Lemma~\ref{lem:Eminus_J3} below).
It does not require mutual orthogonality of $\{e_1,e_2,e_3\}$; it only fixes the
off-diagonal inner products $\langle e_{k-1}, e_j\rangle$ for adjacent indices.
Consequently, the generalized eigenspace is spanned by a single non-diagonalizable Jordan block
of length 3, and $\mathrm{AM} = 3$.
\end{lemma}

\begin{proof}
The mode $e_1$ is established in Lemma~\ref{lem:stokes_gm1}.
For $k = 2, 3$, the subsequent functions $e_k$ are constructed via the variation of parameters formula.
The particular solution to the equation $P_{\mathrm{mod}}\,u = e_{k-1}$ is
\[
  u_k(t) = -v_-(t)\int_{t_L}^t \frac{e_1(s)\,e_{k-1}(s)}{W(e_1,v_-)}\,ds
            + e_1(t)\int_{t_L}^t \frac{v_-(s)\,e_{k-1}(s)}{W(e_1,v_-)}\,ds,
\]
and the outgoing projection $e_k := u_k + C_k\,v_-$, where the projection coefficient
$C_k = -W(u_k,e_1)\big|_{T_R}/W(v_-,e_1)\big|_{T_R}$ is chosen to enforce the boundary condition
$W(e_k, e_1)\big|_{T_R} = 0$.
The linear independence of the set $\{e_1,e_2,e_3\}$ follows directly from the algebraic structure
of the Jordan chain, without requiring numerical verification
(Lemma~\ref{lem:jordan_indep} below).
Since the Jordan chain has length 3, the minimal polynomial of $P_{\mathrm{mod}}$
restricted to $\mathrm{span}\{e_1,e_2,e_3\}$ is $\lambda^3$, confirming a single
$3\times 3$ Jordan block and $\mathrm{AM} = 3$.

\emph{Why $\mathrm{AM} = 3$ and not higher.}
The equation $P_{\mathrm{mod}}\,e_4 = e_3$ would require a fourth outgoing solution
to the inhomogeneous system on $\Gamma$.
By Sibuya's global theory~\cite{Sibuya75} (see in particular \S5, Theorem~5.1 therein),
the model equation $-u'' + x^3 u = 0$ on the Stokes sector containing $\Gamma$
has exactly two Stokes sectors in the upper half-plane, corresponding to two
linearly independent solutions; the outgoing boundary condition on $\Gamma$ selects
a unique one-dimensional subspace (Lemma~\ref{lem:stokes_gm1}).
Consequently the variation of parameters construction terminates at $e_3$:
the forcing term for a putative $e_4$ lies outside the range of $P_{\mathrm{mod}}$
on the outgoing class, and no outgoing $e_4$ exists.
Hence $\mathrm{AM} = 3$ exactly.

\emph{Normalization adjustment.}
Starting from any Jordan chain $\{e_1, e_2, e_3\}$ constructed as above, we replace
$e_k \mapsto e_k - c_k e_1$ for $k = 2, 3$, choosing constants $c_k \in \mathbb{C}$
to enforce $\langle e_{k-1} - c_{k-1}e_1, e_j\rangle = \delta_{j,k-1}$.
Since $P_{\mathrm{mod}} e_1 = 0$, the adjusted basis still satisfies the Jordan chain
relations.  This establishes~\eqref{eq:normalization}.
\end{proof}

\begin{lemma}[$E_{-+}(0,0) = J_3$ from the Jordan chain]
\label{lem:Eminus_J3}
In the basis $\{e_j\}_{j=1}^3$ established in Lemma~\ref{lem:jordan_chain}, the Schur complement
at the critical configuration satisfies
\[
  E_{-+}(0,0) =
  \begin{pmatrix} 0 & 1 & 0 \\ 0 & 0 & 1 \\ 0 & 0 & 0 \end{pmatrix}
  = J_3,
\]
where $J_3$ denotes the $3\times 3$ nilpotent Jordan block.
This relation is a direct consequence of the Jordan chain structure.
\end{lemma}

\begin{proof}
Setting $R_-v = \sum_k v_k \chi e_k$ and $(R_+u)_j = \langle u, \chi e_j\rangle_{L^2}$, we have
\[
  (E_{-+}(0,0))_{jk}
  = (R_+ P_{\mathrm{mod}} R_-)_{jk}
  = \langle P_{\mathrm{mod}}(e_k),\, e_j \rangle_{L^2}.
\]
The Jordan chain relations yield $P_{\mathrm{mod}}\,e_k = e_{k-1}$ (for $k \geq 2$)
and $P_{\mathrm{mod}}\,e_1 = 0$, which gives
\[
  \langle P_{\mathrm{mod}}\,e_k,\, e_j \rangle
  = \langle e_{k-1},\, e_j \rangle
  = \delta_{j,\,k-1},
\]
where the last step uses the normalization condition~\eqref{eq:normalization} established
in Lemma~\ref{lem:jordan_chain}.
Following that adjustment, $(E_{-+}(0,0))_{jk} = \delta_{j,k-1}$, which is equivalent to
$E_{-+}(0,0) = J_3$.
\end{proof}

\begin{theorem}[Jordan Structure at the Critical Point]
\label{thm:jordan}
In the basis $\{e_j\}_{j=1}^3$ constructed in Lemmas~\ref{lem:stokes_gm1}--\ref{lem:Eminus_J3},
the Schur complement satisfies
\[
  E_{-+}(\omega_c, P_c) =
  \begin{pmatrix} 0 & 1 & 0 \\ 0 & 0 & 1 \\ 0 & 0 & 0 \end{pmatrix} = J_3.
\]
The operator $P_{\mathrm{mod}}$ has algebraic multiplicity $\mathrm{AM} = 3$ and
geometric multiplicity $\mathrm{GM} = 1$, and its generalized eigenspace corresponds
to a single non-diagonalizable $3\times 3$ Jordan block.
\end{theorem}

\begin{proof}
The geometric multiplicity is $\mathrm{GM} = 1$ by Lemma~\ref{lem:stokes_gm1}.
The algebraic multiplicity $\mathrm{AM} = 3$ and the single-block structure follow from
Lemma~\ref{lem:jordan_chain}: the Jordan chain $\{e_1,e_2,e_3\}$ has length~3,
which implies that the minimal polynomial is $\lambda^3$, thereby precluding any decomposition into
smaller blocks. The identity $E_{-+}(\omega_c,P_c) = J_3$ is established in Lemma~\ref{lem:Eminus_J3}.
\end{proof}

% -------------------------------------------------------
\section{Geometric Setup and Effective Potential}
% -------------------------------------------------------

The static, spherically symmetric SdS spacetime is characterized by the lapse function
$f(r) = 1 - r_s/r - r^2/L^2$. For perturbations of spin $s=2$ and angular momentum mode $\ell=2$, the Regge--Wheeler effective potential is given by:
\[
  V_{\mathrm{eff}}(r) = f(r)\left[\frac{6}{r^2} - \frac{3r_s}{r^3}\right]
\]

\subsection{The Nariai Limit and Critical Parameter Space}

The Nariai limit corresponds to the extremal configuration where the black hole and cosmological horizons coincide, which occurs when $f(r) = 0$ and $f'(r) = 0$ simultaneously:
\[
  r_{\mathrm{Nariai}} = \frac{L}{\sqrt{3}},
  \qquad
  r_s^{\mathrm{Nariai}} = \frac{2\sqrt{3}}{9}L \approx 0.3849\,L
\]
The critical parameter configuration $P_c$ associated with the $k=3$ degeneracy is determined by
simultaneously solving the equations $V'_{\mathrm{eff}}(r)=0$ and $V''_{\mathrm{eff}}(r)=0$:
\[
  r_s^c = \frac{\sqrt{54+22\sqrt{33}}}{12}\,L \approx 1.1192\,L,
  \qquad
  r^c = \frac{L(-5\sqrt{2}+3\sqrt{66})}{2\sqrt{27+11\sqrt{33}}} \approx 0.9109\,L
\]
The Taylor expansion of the effective potential about $r^c$ confirms the presence of the $A_2$ (cusp) singularity:
\[
  V_{\mathrm{eff}}(r) - V_{\mathrm{eff}}(r^c)
  \sim \tfrac{1}{6}V'''_{\mathrm{eff}}(r^c)(r - r^c)^3
\]
According to the Malgrange preparation theorem~\cite{Malgrange64}, any smooth function
possessing a degenerate critical point of this type is locally equivalent under a smooth
coordinate transformation to the $A_2$ normal form $x^3$. The Taylor expansion above
verifies that $V_{\mathrm{eff}}$ belongs to this equivalence class in the neighborhood of $r^c$.

\begin{lemma}[Permanently Shielded Physical Regime]
\label{lem:physical_shielding}
$r_s^c > r_s^{\mathrm{Nariai}}$.
\end{lemma}

\begin{proof}
Comparing the dimensionless coefficients, it suffices to show that $\frac{\sqrt{54+22\sqrt{33}}}{12} > \frac{2\sqrt{3}}{9}$.
Multiplying both sides by $36$ yields $3\sqrt{54+22\sqrt{33}} > 8\sqrt{3}$. Since both sides are
positive, squaring both sides gives:
\[
  9(54 + 22\sqrt{33}) > 192 \implies 486 + 198\sqrt{33} > 192
\]
This inequality clearly holds. Consequently, the physical region ($r_s \leq r_s^{\mathrm{Nariai}}$) is
permanently shielded from the $V''_{\mathrm{eff}}=0$ degeneracy.
\end{proof}

\begin{lemma}[Meromorphic continuation of the resolvent family]
\label{lem:meromorphic_continuation}
Fix $L > 0$, $h > 0$, and $\omega \in \mathbb{C}$ with
$\mathrm{Im}(\omega) < 0$.
Let $\mu \in \mathbb{C}$ denote the dimensionless deformation parameter
$\mu := (r_s - r_s^c)/L$, where $r_s^c$ is the critical mass defined above.
(This ``physical'' normalisation is used in the present lemma for the meromorphic
continuation.  In the Grushin reduction of Section~\ref{sec:grushin}, $\mu$ is
renormalised by the constant $\alpha_{P} \approx 3.2267$ defined in
Lemma~\ref{lem:Epm_expand} so that the unfolding operator $\partial_{\mu}P_{h}$
equals $x$ exactly in BNF coordinates; the two conventions differ only by this
nonzero scalar and leave all resolvent bounds unchanged.)
Then:
\begin{enumerate}
\item[(i)] \emph{(Analyticity.)}
  {\looseness=-1 The Regge--Wheeler effective potential $V_{\mathrm{eff}}(r;\mu)$
  is a real-analytic function of $\mu \in \mathbb{R}$ for each $r > 0$,
  and extends to a holomorphic function of $\mu$ on the complex strip
  $|\mathrm{Im}(\mu)| < \delta$ for some $\delta > 0$ depending
  only on $r_s^c$ and $L$.}

\item[(ii)] \emph{(Fredholm property.)}
  For each $\mu$ near $0$ in $\mathbb{C}$, the operator
  \[
    P_h(\mu) - \omega^2 \;:\; H^2_{\mathrm{sc}}(M) \to L^2(M)
  \]
  is Fredholm of index zero, where $H^2_{\mathrm{sc}}$ is the
  semiclassical Sobolev space with outgoing boundary conditions.

\item[(iii)] \emph{(Meromorphic continuation.)}
  The resolvent $R_h(\mu,\omega) := (P_h(\mu) - \omega^2)^{-1}$,
  initially defined for $r_s < r_s^{\mathrm{Nariai}}$ (i.e.\
  $\mu < \mu_{\mathrm{Nariai}} := (r_s^{\mathrm{Nariai}} - r_s^c)/L$),
  extends meromorphically as a function of $\mu$ to the entire
  complex strip $|\mathrm{Im}(\mu)| < \delta$. The poles of this
  meromorphic family correspond to quantum resonances of $P_h(\mu)$.

\item[(iv)] \emph{(Nature of $P_c$.)}
  The critical point $\mu = 0$ (i.e.\ $P = P_c$) is a pole of the
  meromorphic family $R_h(\mu, \omega_c)$, and the residue at this
  pole encodes the non-diagonalizable Jordan block structure of
  Theorem~\ref{thm:jordan}.
\end{enumerate}
\end{lemma}

\begin{proof}
We verify the four claims in order.

\medskip
\noindent\textbf{Claim (i): Analyticity of $V_{\mathrm{eff}}(r;\mu)$.}

The lapse function is $f(r; r_s) = 1 - r_s/r - r^2/L^2$, and the
effective potential is
\[
  V_{\mathrm{eff}}(r; r_s) = f(r;r_s)\!\left[\frac{6}{r^2} - \frac{3r_s}{r^3}\right].
\]
Both factors are polynomial in $r_s$ with coefficients that are smooth
functions of $r > 0$. Therefore $V_{\mathrm{eff}}(r;\mu)$ is
holomorphic in $\mu \in \mathbb{C}$ for each fixed $r > 0$.

\begin{remark}[On the role of the tortoise coordinate]
The operator $P_h = -h^2\partial_{r_*}^2 + V_{\mathrm{eff}}(r)$
involves the tortoise coordinate $r_* = \int^r f(r';r_s)^{-1}dr'$,
which is singular at the horizons $f = 0$. However, as a family
parametrized by $\mu$, the operator is most naturally viewed in terms
of $V_{\mathrm{eff}}(r;\mu)$ rather than through $r_*$: the analytic
structure in $\mu$ is that of $V_{\mathrm{eff}}$, not of the tortoise
coordinate map. Specifically, in the scattering framework of
Vasy~\cite{Vasy2013}, the relevant operator is extended to a compact
manifold with boundary, and $\mu$-analyticity is inherited from
$V_{\mathrm{eff}}$.
\end{remark}

\medskip
\noindent\textbf{Claim (ii): Fredholm property.}

For real $\mu$ with $r_s(\mu) < r_s^{\mathrm{Nariai}}$ (the physically
regular regime), the Fredholm property of $P_h(\mu) - \omega^2$ follows
from the standard scattering theory of Vasy~\cite{Vasy2013}: the SdS
static patch is a manifold with two radial points (the black hole and
cosmological horizons), and the operator $P_h - \omega^2$ is
elliptic at the horizons with appropriate sign conditions for
$\mathrm{Im}(\omega) < 0$. This gives index zero by the general
framework.

For complex $\mu$ near $0$, the Fredholm property is stable; we
verify this by checking that $P_h(\mu) - \omega^2$ forms an analytic
family of type~(A) in the sense of Kato~\cite{Kato66}, \S VII.2.1,
which then allows us to invoke index stability.

\emph{Type-(A) verification for $P_h(\mu)$.}
A family $T(\mu)$ is of type~(A) if (i)~the domain
$D(T(\mu)) = D$ is independent of $\mu$, and (ii)~for each $u \in D$
the map $\mu \mapsto T(\mu)u$ is holomorphic in the Hilbert space norm.
Here $D = H^2_{\mathrm{sc}}(M)$, the semiclassical Sobolev space with
outgoing scattering boundary conditions as defined in
Vasy~\cite{Vasy2013}.

Condition~(i): the domain $H^2_{\mathrm{sc}}(M)$ is defined via the
elliptic structure of $P_h(0) - \omega^2$ at the boundary of $M$ (the
two horizons), and this elliptic structure is $\mu$-independent: the
principal symbol $\xi^2 + V_{\mathrm{eff}}(r;\mu)$ is elliptic at the
boundary for all $|\mu|$ sufficiently small, since $V_{\mathrm{eff}}$
is bounded on compact sets and the ellipticity condition depends only on
the sign of $\mathrm{Im}(\omega)$, which is fixed.  The tortoise
coordinate $r_*$, while singular at the horizons $f(r;r_s)=0$, enters
the operator only through the \emph{metric factor} $f$; in Vasy's
framework the operator is extended to a compact manifold with corners,
and the relevant domain is defined at the level of this compactified
operator, which has $\mu$-independent boundary symbol.  Concretely,
$V_{\mathrm{eff}}(r;\mu) = f(r;r_s(\mu))[6r^{-2}-3r_s(\mu)r^{-3}]$ is
polynomial in $\mu$ with smooth $r$-dependent coefficients, so the
change $\mu \mapsto \mu + \delta\mu$ shifts $P_h$ by a bounded
multiplication operator $\delta\mu \cdot \partial_\mu V_{\mathrm{eff}}$
that maps $H^2_{\mathrm{sc}} \to L^2$ continuously.  The domain is
therefore $\mu$-stable.

Condition~(ii): holomorphicity of $\mu \mapsto P_h(\mu)u$ in $L^2(M)$
follows from Claim~(i): since $V_{\mathrm{eff}}(r;\mu)$ is holomorphic
in $\mu$ for each fixed $r$, the map $\mu \mapsto V_{\mathrm{eff}}(\cdot;\mu)u$
is holomorphic in $L^2$ for each $u \in H^2_{\mathrm{sc}}$.

With both conditions verified, $P_h(\mu) - \omega^2$ is an analytic
family of type~(A). The Fredholm index is locally constant under
analytic perturbations (Kato~\cite{Kato66}, Chapter~IV, Theorem~5.17),
so the index zero established in the real regime persists to all
$|\mu|$ sufficiently small.

\medskip
\noindent\textbf{Claim (iii): Meromorphic continuation.}

We apply the analytic Fredholm theorem
(Dyatlov--Zworski~\cite{DZ19}, Theorem~C.5) to the holomorphic
Fredholm family
\[
  \mu \mapsto P_h(\mu) - \omega^2 \;:\;
  H^2_{\mathrm{sc}}(M) \to L^2(M).
\]
The three hypotheses of Theorem~C.5 are:
\begin{enumerate}
\item[\textup{(H1)}] \emph{Holomorphicity}: verified by Claim~(i).
\item[\textup{(H2)}] \emph{Fredholm of index zero}: verified by
  Claim~(ii).
\item[\textup{(H3)}] \emph{Invertibility at one point}: for
  $r_s \ll r_s^{\mathrm{Nariai}}$ (small black hole mass, large
  $|\mu|$), the operator $P_h(\mu) - \omega^2$ is invertible for
  generic $\omega$ with $\mathrm{Im}(\omega) < 0$ by the standard
  resonance-free half-plane result of Dyatlov~\cite{Dyatlov11}.
\end{enumerate}
The conclusion of Theorem~C.5 is that $R_h(\mu,\omega) =
(P_h(\mu)-\omega^2)^{-1}$ is a meromorphic family of bounded operators
$L^2(M) \to H^2_{\mathrm{sc}}(M)$ in $\mu$, with poles of finite rank.

\medskip
\noindent\textbf{Claim (iv): Jordan structure at $\mu = 0$.}

At $\mu = 0$ and $\omega = \omega_c$, the operator $P_h(0) - \omega_c^2$
has a three-dimensional generalized kernel spanned by the Jordan chain
$\{e_1, e_2, e_3\}$ established in Lemma~\ref{lem:jordan_chain}. By the
general theory of meromorphic Fredholm families (see Dyatlov--Zworski
\cite{DZ19}, \S C.3), the Laurent expansion of $R_h(\mu, \omega_c)$
near $\mu = 0$ has a pole of order equal to the length of the longest
Jordan chain, which is $3$ by Theorem~\ref{thm:jordan}. The residue
at this pole is a rank-$3$ operator whose range is
$\mathrm{span}\{e_1,e_2,e_3\}$, and its non-diagonalizable structure
is an intrinsic property of the continued family, not an artifact of
the continuation.
\end{proof}

\begin{remark}[On the Nariai boundary]
\label{rem:nariai_boundary}
The meromorphic continuation of $R_h(\mu,\omega)$ passes through the
Nariai boundary $\mu = \mu_{\mathrm{Nariai}}$ without obstruction,
because the analyticity of $V_{\mathrm{eff}}(r;\mu)$ in $\mu$ does not
require the absence of a horizon. The Nariai geometry ($r_h = r_c$,
coalescing horizons) corresponds to a degenerate metric, but the
effective potential $V_{\mathrm{eff}}(r;\mu)$ and hence the operator
$P_h(\mu)$ remain well-defined and analytic in $\mu$ through this
boundary. The restriction $r_s < r_s^{\mathrm{Nariai}}$ is a physical
constraint (ensuring two distinct horizons) but not a mathematical
obstruction to the resolvent continuation.
\end{remark}

% -------------------------------------------------------
\section{Hamiltonian Flow and Semiclassical Scaling}
% -------------------------------------------------------

The classical Hamiltonian symbol is given by $p(r,\xi) = \xi^2 + V_{\mathrm{eff}}(r)$. On the trapped set $K_P$, the
linearization of the Hamiltonian flow yields the spectral gap
$\kappa_{\mathrm{gap}}(P) = \sqrt{-2V''_{\mathrm{eff}}(r^c)}$. As the parameters approach the critical configuration ($P \to P_c$),
the second derivative vanishes ($V''_{\mathrm{eff}}(r^c) \to 0$), which forces the spectral gap to zero
($\kappa_{\mathrm{gap}}(P) \to 0$) and causes the control provided by the escape function to degenerate.

\subsection{Birkhoff Normal Form Reduction and FIO Conjugation}

We construct the Birkhoff normal form (BNF) conjugation explicitly, reducing $P_h$ to the
model operator $P_{\mathrm{mod}} = -h^2\partial_x^2 + x^3$ near the critical point $r^c$.

\textbf{Step 1: Symplectic straightening.}
At the critical point $(r^c, 0) \in T^*\mathbb{R}$, the Hamiltonian symbol is
$p(r,\xi) = \xi^2 + V_{\mathrm{eff}}(r)$.
Under the degeneracy conditions $V'_{\mathrm{eff}}(r^c) = V''_{\mathrm{eff}}(r^c) = 0$ and $V'''_{\mathrm{eff}}(r^c) \neq 0$,
the Taylor expansion of $p(r, \xi)$ about $r^c$ yields
\[
  p(r,\xi) = \xi^2 + \tfrac{1}{6}V'''_{\mathrm{eff}}(r^c)(r-r^c)^3 + O((r-r^c)^4).
\]
By the Malgrange preparation theorem~\cite{Malgrange64}, there exists a smooth local
coordinate change $r \mapsto x(r)$ defined in a neighborhood of $r^c$, with $x(r^c)=0$, such that
$V_{\mathrm{eff}}(r) - V_{\mathrm{eff}}(r^c) = x(r)^3$ holds identically.
Setting $\omega_c^2 := V_{\mathrm{eff}}(r^c)$, we consider the symbol $p - \omega_c^2 = \xi^2 + x^3$
in these coordinates.

\textbf{Step 2: Canonical lift and FIO construction.}
The diffeomorphism $r \mapsto x(r)$ lifts to a symplectomorphism
$\kappa: T^*\mathbb{R} \to T^*\mathbb{R}$ defined by
\[
  \kappa(r, \xi) = \bigl(x(r),\; \xi / x'(r)\bigr),
\]
which satisfies $p \circ \kappa^{-1}(x, \eta) = \eta^2 + x^3$.
At the operator level, we associate with $\kappa$ an $h$-Fourier integral operator
$U_h: L^2(\mathbb{R}) \to L^2(\mathbb{R})$ of order zero, which is microlocally unitary near
$(r^c, 0)$, defined via its Schwartz kernel
\[
  K_{U_h}(x, r) = (2\pi h)^{-1/2}\,a(x,r;h)\,e^{i\phi(x,r)/h},
\]
where the phase function $\phi(x, r)$ generates the canonical transformation $\kappa$ and satisfies
$\partial_r \phi = \xi$ and $\partial_x \phi = -\eta$, while the amplitude satisfies $a \in S^0$
with principal symbol $a(x(r),r;0) = |x'(r)|^{-1/2}$, representing the
half-density correction that ensures $L^2$ unitarity to leading order.

\begin{lemma}[$L^2$ boundedness and microlocal unitarity of $U_h$]
\label{lem:FIO_unitary}
The operator $U_h$ defined above is bounded on $L^2(\mathbb{R})$ uniformly in $h$,
and satisfies $U_h^* U_h = I + O(h^\infty)_{\mathcal{L}(L^2)}$ microlocally near $(r^c,0)$.
\end{lemma}

\begin{proof}
Uniform $L^2$ boundedness follows because $U_h$ is an $h$-FIO of order
zero associated with a local canonical graph; by the Seeger--Sogge--Stein
theorem (\cite{Zworski2012}, Theorem~10.1) it is bounded on
$L^2(\mathbb{R})$ with norm $O(1)$ uniformly in $h$.

Microlocal unitarity is achieved by the choice of the principal amplitude
$a(x(r),r;0) = |x'(r)|^{-1/2}$.
Indeed, the composition $U_h^* U_h$ is an $h$-pseudodifferential operator
whose principal symbol is given by
\[
  \sigma(U_h^* U_h)(r,\xi)
  = |a(x(r),r;0)|^2 \cdot |x'(r)|
  = |x'(r)|^{-1} \cdot |x'(r)|
  = 1,
\]
where the factor $|x'(r)|$ arises from the Jacobian of the coordinate change in
the stationary phase expansion of the composition.
This yields $\sigma(U_h^* U_h) = 1$. The calculus of $h$-pseudodifferential operators
(see Dyatlov--Zworski~\cite{DZ19}, Theorem~4.18) then implies
$U_h^* U_h = I + O(h)_{\mathcal{L}(L^2)}$.
The expansion can be improved to $O(h^\infty)$ by iteratively solving the transport
equations for the amplitude at each order in $h$. The sub-principal amplitudes $a_k$
are chosen to satisfy the corresponding homological equations; this is the standard
parametrix construction for FIOs (see Zworski~\cite{Zworski2012}, \S10.2).
\end{proof}

\begin{remark}[Global extension of $U_h$ outside $\operatorname{supp}\chi_{\mathrm{trap}}$]
\label{rem:FIO_global}
The FIO $U_h$ is defined microlocally near $(r^c, 0)$, but the Grushin
operator $\mathcal{G}$ is a global operator on $L^2(M)$.
The BNF conjugation identity~\eqref{eq:BNF_conjugation} is used only
in the microlocal region $\operatorname{supp}\chi_{\mathrm{trap}}$, via the
cutoff $\chi$ in the definitions of $R_\pm$.
Outside this region, the operator $P_h - \omega_c^2$ is elliptic
(since $V_\mathrm{eff}(r) - \omega_c^2 \neq 0$ away from $r^c$),
so the resolvent $(P_h - \omega_c^2)^{-1}$ is bounded there by
standard elliptic estimates.
The global estimate on $\|R_h(P)\|$ is obtained by the standard
microlocal partition-of-unity argument:
$R_h = \chi_{\mathrm{trap}} R_h \chi_{\mathrm{trap}} +
(1-\chi_{\mathrm{trap}})R_h + R_h(1-\chi_{\mathrm{trap}}) +
O(h^\infty)$
(cf.\ Dyatlov--Zworski~\cite{DZ19}, \S4.2, proof of Theorem~4.1),
where the off-diagonal pieces are $O(1)$ by elliptic regularity, and
only the trapped piece $\chi_{\mathrm{trap}} R_h \chi_{\mathrm{trap}}$
carries the $h^{-6/5}$ divergence controlled by $\|E_{-+}^{-1}\|$.
\end{remark}

\textbf{Step 3: Conjugation identity.}
By standard FIO composition calculus
(see \cite{Zworski2012}, Ch.~11 and \cite{DZ19}, App.~E)
and Lemma~\ref{lem:FIO_unitary}, we have
\[
  U_h^* \circ (P_h - \omega_c^2) \circ U_h
  = \mathrm{Op}^w_h\!\bigl(p \circ \kappa^{-1}\bigr) + O(h^2)_{\mathcal{L}(L^2)}
  = -h^2\partial_x^2 + x^3 + O(h^2)_{\mathcal{L}(L^2)}.
\]
The $O(h)$ term vanishes identically because $Q_1 \equiv 0$ (Lemma~\ref{lem:Q1}, Appendix~A).
Specifically, the sub-principal symbol of $P_h - \omega_c^2$ in Weyl quantization is zero,
and the sub-principal contribution from the FIO conjugation also vanishes, as the
Jacobian of the cotangent lift $\kappa$ is identically 1 (see Appendix~A, Step~2).
Consequently, the $O(h)$ remainder in the conjugation identity is absent, yielding
\begin{equation}\label{eq:BNF_conjugation}
  U_h^* (P_h - \omega_c^2) U_h = P_{\mathrm{mod}} + O(h^2), \qquad
  P_{\mathrm{mod}} := -h^2\partial_x^2 + x^3.
\end{equation}
The $O(h^2)$ remainder corresponds to the Schwarzian contribution $q_2$ derived in Appendix~B.2.
This term enters the Grushin matrix only at order $h^2 Q_{31}$, which is subleading
in the characteristic equation (see Appendix~B, Proposition~\ref{prop:dominant_balance}).

\subsection{Semiclassical Dilation}

Starting from the model operator $P_{\mathrm{mod}} = -h^2\partial_x^2 + x^3$, we define the
unitary dilation $x = h^{2/5}y$, which acts as $(D_{h^{2/5}}u)(y) = h^{1/5}u(h^{2/5}y)$.
A direct calculation yields:
\[
  D_{h^{2/5}}^{-1} \circ P_{\mathrm{mod}} \circ D_{h^{2/5}}
  = -h^2 \cdot h^{-4/5}\partial_y^2 + h^{6/5}y^3
  = h^{6/5}\bigl(-\partial_y^2 + y^3\bigr).
\]
The scaling exponent is determined by balancing the kinetic and potential terms:
$h^{2-2\alpha} = h^{3\alpha}$ yields $\alpha = 2/5$, which is consistent with the dilation scaling.
We thus obtain:
\[
  P_{\mathrm{mod}} = h^{6/5}\,Q_0, \qquad Q_0 := -\partial_y^2 + y^3,
\]
and the resolvent at energy zero satisfies
\[
  (P_{\mathrm{mod}})^{-1} = h^{-6/5}\,Q_0^{-1},
\]
yielding the characteristic semiclassical factor $\|R_h\| \sim h^{-6/5}$.
The vanishing of the sub-principal symbol $Q_1$, which ensures that no $O(h)$
correction alters this scaling, is established in Appendix~A.

% -------------------------------------------------------
\section{Grushin Reduction}
\label{sec:grushin}
% -------------------------------------------------------

Let $\{e_j\}_{j=1}^3$ be the Jordan basis of outgoing quasimodes constructed in
Appendix~C, and let $\chi \in C_c^\infty(T^*M)$ be a microlocal cutoff supported
near the trapped set $K_P$. We define the auxiliary operators:
\[
  R_- : \mathbb{C}^3 \to L^2(M),
  \quad
  R_- v = \sum_{j=1}^3 v_j \chi e_j,
\]
\[
  R_+ : L^2(M) \to \mathbb{C}^3,
  \quad
  (R_+ u)_j = \langle u,\, \chi e_j \rangle_{L^2}.
\]
The extended Grushin operator is then given by:
\[
  \mathcal{G}(\omega, P) =
  \begin{pmatrix} P_h - \omega^2 & R_- \\ R_+ & 0 \end{pmatrix}
  : L^2(M) \oplus \mathbb{C}^3 \to L^2(M) \oplus \mathbb{C}^3.
\]
For parameters $|P - P_c| \geq \varepsilon > 0$, the operator $\mathcal{G}$ is Fredholm of index~$0$ and
is invertible as shown below.

\begin{lemma}[Invertibility of the Grushin matrix]
\label{lem:grushin_invertible}
At $(\omega, P) = (\omega_c, P_c)$, the extended operator $\mathcal{G}(0,0)$ is
boundedly invertible on $L^2(M) \oplus \mathbb{C}^3$.
\end{lemma}

\begin{proof}
We write the system $\mathcal{G}(0,0)(u,v) = (f,g)$ explicitly:
\[
  (P_h - \omega_c^2)\,u + R_-\,v = f,
  \qquad
  R_+\,u = g.
\]
Since $\{e_j\}_{j=1}^3$ is the Jordan basis established in Lemma~\ref{lem:jordan_chain},
the operator $P_h - \omega_c^2$ is not invertible on $L^2(M)$ alone due to
its three-dimensional generalized kernel; however, the augmented system is invertible.
Indeed, decomposing $u = u_\perp + \sum_j c_j e_j$, where $u_\perp \perp \mathrm{span}\{e_j\}$,
the second equation $R_+u = g$ determines the coefficients $c_j$ as $(R_+u)_j = c_j + \langle u_\perp, \chi e_j\rangle = g_j$,
which yields $c_j = g_j - \langle u_\perp, \chi e_j\rangle$.
The first equation, projected onto the orthogonal complement of $\mathrm{span}\{e_j\}$, determines $u_\perp$
uniquely via the restricted resolvent of $P_h - \omega_c^2$.
This restricted resolvent is bounded since $P_h - \omega_c^2$ is Fredholm of index~$0$ and
is injective on $(\mathrm{span}\{e_j\})^\perp$ (see Dyatlov--Zworski~\cite{DZ19}, Theorem~4.1);
denote its operator norm by $M_{\mathrm{res}} < \infty$.
The solution $(u,v)$ depends continuously on the data $(f,g)$, establishing bounded invertibility.
Concretely, the full inverse satisfies
\[
  \|\mathcal{G}(0,0)^{-1}\|_{\mathcal{L}(L^2(M)\oplus\mathbb{C}^3)}
  \;\leq\; C(M_{\mathrm{res}},\|R_\pm\|) \;<\; \infty,
\]
where the constant $C$ depends only on $M_{\mathrm{res}}$ and the norms of the auxiliary operators
$R_\pm$ defined in Section~\ref{sec:grushin}.
This finite bound is recorded as
$M_0 := \|\mathcal{G}(0,0)^{-1}\|$ in Remark~\ref{rem:grushin_validity},
where it determines the radius of validity of the Grushin expansion.
\end{proof}

By Lemma~\ref{lem:grushin_invertible}, the Grushin operator $\mathcal{G}(0,0)$ is boundedly invertible.
By applying the block inversion formula, the Schur complement (representing the effective Hamiltonian) is:
\[
  E_{-+}(\omega, P) = -R_+(P_h - \omega^2)^{-1}_{\mathrm{res}}\,R_- ,
\]
where $(P_h-\omega^2)^{-1}_{\mathrm{res}}$ denotes the restricted resolvent on
$(\mathrm{span}\{e_j\})^\perp$. In the notation of Dyatlov--Zworski~\cite{DZ19}, Appendix~C.1,
this simplifies to
\[
  E_{-+}(\omega) = R_+(P_h - \omega^2)R_-
\]
in the basis where $E_{-+}(0,0) = J_3$.
(Under the Jordan chain convention $P_{\mathrm{mod}}\,e_k = e_{k-1}$ and
the definition of $R_+$, the relation $(E_{-+}(0,0))_{jk} = \langle e_{k-1}, e_j\rangle = \delta_{j,k-1} = J_3$
holds directly; see Lemma~\ref{lem:Eminus_J3}.)
The quantum resonances correspond to the values of $\omega$ for which the extended operator
$\mathcal{G}(\omega,P)$ fails to be invertible. By the Schur complement formula, this condition is
equivalent to the vanishing of the determinant, $\det(E_{-+}(\omega,P)) = 0$.

\begin{lemma}[Parametric expansion of $E_{-+}$]
\label{lem:Epm_expand}
Set $\lambda := \omega^2 - \omega_c^2$ and $\mu := P - P_c$.
The Schur complement admits the expansion
\begin{equation}\label{eq:lambda_def}
  E_{-+}(\omega, P) = J_3 - \lambda\,I + \mu\,X + O(\lambda^2, \mu^2, \lambda\mu),
\end{equation}
where $J_3$ is the nilpotent Jordan block established in Theorem~\ref{thm:jordan} and
$X_{ij} = \langle e_i, x\,e_j\rangle$ is the geometric perturbation matrix.
\end{lemma}

\begin{proof}
We write the Grushin matrix as
\[
  \mathcal{G}(\lambda,\mu)
  = \mathcal{G}(0,0)
    + \lambda \begin{pmatrix} -I & 0 \\ 0 & 0 \end{pmatrix}
    + \mu \begin{pmatrix} \partial_\mu A & 0 \\ 0 & 0 \end{pmatrix}
    + O(\lambda^2,\mu^2,\lambda\mu),
\]
where $A = P_h - \omega^2$ and $\partial_\mu A = \partial_P P_h$.
Since $\mathcal{G}(0,0)$ is invertible (the Jordan basis $\{e_j\}$ is constructed
such that $R_+$ and $R_-$ render $\mathcal{G}(0,0)$ an isomorphism;
see Dyatlov--Zworski~\cite{DZ19}, \S4.4), the perturbation formula for the Schur
complement yields
\[
  E_{-+}(\lambda,\mu)
  = E_{-+}(0,0)
    - \lambda\,(R_+ R_-)
    + \mu\,(R_+\,\partial_\mu A\,R_-)
    + O(\lambda^2,\mu^2,\lambda\mu).
\]
By Lemma~\ref{lem:Eminus_J3}, we have $E_{-+}(0,0) = J_3$.

\emph{The $-\lambda I$ term.}
By Lemma~\ref{lem:Eminus_J3}, $\langle e_{k-1}, e_j\rangle = \delta_{j,k-1}$,
which gives $(R_+R_-)_{jk} = \delta_{jk}$, i.e.\ $R_+R_- = I_3$.
Hence $-\lambda(R_+R_-)$ contributes $-\lambda I_3$.

\emph{The $\mu X$ term.}
We must verify that $R_{+}(\partial_{\mu}A)R_{-} = X$, i.e.\ that in the
Birkhoff normal form coordinate $x$ introduced in Section~4, the operator
$\partial_{\mu}P_{h}$ acts (at leading order) as multiplication by $x$.
Here $\mu$ parametrises the one-parameter family $P(\mu)$ with $P(0)=P_{c}$;
we choose the normalisation of this path precisely so that the following holds.

Since $P_{h} = -h^{2}\partial_{r_{*}}^{2} + V_{\mathrm{eff}}(r;\,r_{s})$ and we
vary $r_{s}$ (keeping $L$ fixed), we have
\[
  \partial_{\mu}P_{h} \;=\; \frac{dr_{s}}{d\mu}\,\partial_{r_{s}}V_{\mathrm{eff}}(r;\,r_{s}^{c}).
\]
Expanding $\partial_{r_{s}}V_{\mathrm{eff}}$ around the critical radius $r^{c}$ in
the BNF coordinate $x$ (where $r - r^{c} = x/x'(r^{c}) + O(x^{2})$ and
$x'(r^{c}) = \bigl(V_{\mathrm{eff}}'''(r^{c})/6\bigr)^{1/3}$ is the \emph{real} cube root,
negative since $V_{\mathrm{eff}}'''(r^{c}) < 0$; the explicit value is given
in~\eqref{eq:xprime_explicit} below) gives
\[
  \partial_{r_{s}}V_{\mathrm{eff}}(r;\,r_{s}^{c})
  \;=\;
  A_{0} \;+\; A_{1}\,\frac{x}{x'(r^{c})} \;+\; O(x^{2}),
\]
where
\[
  A_{0} \;:=\; \partial_{r_{s}}V_{\mathrm{eff}}\!\big|_{(r^{c},\,r_{s}^{c})},
  \qquad
  A_{1} \;:=\; \partial_{r}\partial_{r_{s}}V_{\mathrm{eff}}\!\big|_{(r^{c},\,r_{s}^{c})}.
\]
A direct computation using $V_{\mathrm{eff}} = f(r;r_{s})\bigl(6r^{-2}-3r_{s}r^{-3}\bigr)$
and the two critical-point relations
$V_{\mathrm{eff}}'(r^{c};r_{s}^{c})=0$, $V_{\mathrm{eff}}''(r^{c};r_{s}^{c})=0$
(whose unique physical solution is $r_{s}^{c} = \tfrac{9+\sqrt{33}}{12}\,r^{c}$,
$L^{2} = \tfrac{9+2\sqrt{33}}{17}(r^{c})^{2}$; see Remark~\ref{rem:omega_c_imaginary})
yields
\begin{equation}\label{eq:xprime_explicit}
  A_{0} \;=\; \frac{3(29-3\sqrt{33})}{2(r^{c})^{3}(9+2\sqrt{33})}, \qquad
  A_{1} \;=\; \frac{-102}{(r^{c})^{4}(9+2\sqrt{33})},
\end{equation}
\begin{equation*}
  x'(r^{c}) \;=\; -\left(\frac{297+49\sqrt{33}}{4(r^{c})^{5}(9+2\sqrt{33})}\right)^{1/3},
\end{equation*}
where the right-hand expression for $x'(r^{c})$ is the real cube root of
$C = \tfrac{1}{6}V_{\mathrm{eff}}'''(r^{c}) < 0$; equivalently $x'(r^{c}) = C^{1/3} < 0$
(real branch).
In particular $A_{1} \neq 0$ and $x'(r^{c}) \neq 0$.

The constant term $\mu\,(dr_{s}/d\mu)\,A_{0}$ is $O(\mu)$ and independent of $r$;
it shifts $\omega_{c}^{2}$ to $\omega_{c}^{2}(\mu) = \omega_{c}^{2} + \mu(dr_{s}/d\mu)A_{0} + O(\mu^{2})$
and is absorbed into the redefinition $\lambda := \omega^{2} - \omega_{c}^{2}(\mu)$.
The $x$-linear term yields
\[
  \partial_{\mu}P_{h}
  \;=\;
  \frac{dr_{s}}{d\mu}\,\frac{A_{1}}{x'(r^{c})}\cdot x \;+\; O(x^{2}).
\]
Define the \emph{path normalisation constant}
\[
  \alpha_{P}
  \;:=\;
  L\,\frac{A_{1}}{x'(r^{c})}
  \;=\;
  \frac{L\,\partial_{r}\partial_{r_{s}}V_{\mathrm{eff}}\big|_{(r^{c},r_{s}^{c})}}
       {\bigl(V_{\mathrm{eff}}'''(r^{c})/6\bigr)^{1/3}}
  \;\approx\; 3.2267 \;\neq\; 0.
\]
We \emph{choose} the one-parameter path so that $dr_{s}/d\mu = L/\alpha_{P}$, i.e.\
\[
  r_{s}(\mu) \;=\; r_{s}^{c} \;+\; \frac{L}{\alpha_{P}}\,\mu,
\]
which is equivalent to working with the normalised deformation parameter
$\mu = \alpha_{P}\,(r_{s}-r_{s}^{c})/L$ (cf.\ Lemma~\ref{lem:meromorphic_continuation},
where the ``physical'' parameter $\hat\mu := (r_{s}-r_{s}^{c})/L = \mu/\alpha_{P}$
is used; the two parametrisations differ by the nonzero constant $\alpha_{P}$
and hence produce the same qualitative spectral picture).

With this choice $\partial_{\mu}P_{h} = x + O(x^{2})$ in BNF coordinates, so
\[
  (R_{+}(\partial_{\mu}P_{h})R_{-})_{jk}
  \;=\; \langle e_{j},\,x\,e_{k}\rangle + O(\text{subleading})
  \;=:\; X_{jk},
\]
and therefore $R_{+}(\partial_{\mu}P_{h})R_{-} = X$ as claimed.

Combining these three contributions yields~\eqref{eq:lambda_def}.
\end{proof}

\begin{remark}[Validity region of the Grushin reduction]
\label{rem:grushin_validity}
The expansion~\eqref{eq:lambda_def} and all subsequent Schur complement computations are
valid in the region
\[
  \mathcal{V} := \bigl\{(\mu, h) \in \mathbb{R} \times (0,1] :
    |\mu| < \mu_0,\; h < h_0 \bigr\},
\]
where the thresholds $\mu_0 > 0$ and $h_0 > 0$ are determined by two conditions.

\emph{All constants are $h$-independent.}
The entire Grushin reduction is carried out at the level of the dilated
model operator $Q_0 = -\partial_y^2 + y^3$ (Section~4.2).  After the
unitary dilation $x = h^{2/5}y$, the operator $P_{\mathrm{mod}} = h^{6/5}Q_0$
and $h$ no longer appears in $Q_0$.  The double-scaling parameter
$\nu := \mu h^{-4/5}$ is $O(1)$ in the scaling window, so the
Grushin matrix $\mathcal{G}(0,0)$ is built from $Q_0$ and the
$h$-independent Jordan chain $\{e_j\}$.  Consequently, all operator
norms appearing in the thresholds below ($M_0$, $C_1$, $\|B_1\|$,
$\|B_2\|$) are norms of fixed operators on the $h$-independent Hilbert
space $L^2_\Gamma \oplus \mathbb{C}^3$, and are therefore independent of
$h$.

\emph{Definition of $M_0$.}
By Lemma~\ref{lem:grushin_invertible}, the operator $\mathcal{G}(0,0)$ is boundedly
invertible on $L^2(M)\oplus\mathbb{C}^3$.
We set
\[
  M_0 := \|\mathcal{G}(0,0)^{-1}\|_{\mathcal{L}(L^2(M)\oplus\mathbb{C}^3,\,L^2(M)\oplus\mathbb{C}^3)}.
\]
This norm is finite by Lemma~\ref{lem:grushin_invertible}; it depends only on
$(r_s^c, L)$ and the normalization of the Jordan basis $\{e_j\}$, both of which
are fixed geometric data.

\emph{Definition of $C_1$.}
The Grushin matrix satisfies
$\mathcal{G}(\lambda,\mu) = \mathcal{G}(0,0) + \lambda B_1 + \mu B_2 + \mathcal{R}(\lambda,\mu)$,
where $B_1, B_2$ are the first-order coefficient operators (bounded, with norms
depending only on $\|R_\pm\|$ and $\|\partial_\mu P_h\|$) and the remainder
satisfies $\|\mathcal{R}(\lambda,\mu)\| \leq C_{\mathcal{R}}(\lambda^2 + \mu^2 + |\lambda\mu|)$
for a constant $C_{\mathcal{R}} > 0$ depending only on the second-order Taylor
coefficients of $P_h$ at $P_c$.
The perturbation formula for the Schur complement (Dyatlov--Zworski~\cite{DZ19},
Appendix~C.1) then produces remainder terms in~\eqref{eq:lambda_def} bounded by
$C_1 M_0 (\lambda^2 + \mu^2 + |\lambda\mu|)$, where we set
\[
  C_1 := C_{\mathcal{R}} + M_0(\|B_1\|^2 + \|B_2\|^2).
\]
(The constant $C_{\mathcal{R}}$ coincides with the constant $C_R$ appearing in
Lemma~\ref{lem:E_bound} (Appendix~A): both bound the operator norm of the second-order
Taylor remainder of $\mathcal{G}(\lambda,\mu)$ at $(\lambda,\mu)=(0,0)$.
Lemma~\ref{lem:E_bound} establishes that this remainder is finite; the definition
here makes the dependence of $C_1$ on the same constant explicit.)

\emph{Explicit thresholds and the uniform bound $C_0 = 2M_0$.}
The Neumann series
\[
  \mathcal{G}(\mu,h)^{-1}
  = \sum_{n=0}^{\infty}(-\mathcal{G}(0,0)^{-1}\delta\mathcal{G})^n \mathcal{G}(0,0)^{-1}
\]
converges in $\mathcal{L}(L^2(M)\oplus\mathbb{C}^3)$ whenever $\|\mathcal{G}(0,0)^{-1}\delta\mathcal{G}\| \leq \tfrac{1}{2}$,
and in that case $\|\mathcal{G}(\mu,h)^{-1}\| \leq 2M_0$.
Requiring $|\lambda| \leq |\mu|$ (on-shell near $\omega_c$) and bounding
$\|\delta\mathcal{G}\| \leq |\mu|(\|B_1\|+\|B_2\|) + C_{\mathcal{R}}(\mu^2+h^2)$,
the half-contractivity condition $\|\mathcal{G}(0,0)^{-1}\delta\mathcal{G}\| \leq \tfrac{1}{2}$
is guaranteed for
\[
  |\mu| < \mu_0 := \frac{1}{4 C_1 M_0 (\|X\| \vee 1)},
  \qquad
  h < h_0 := \min\!\left(1,\,\frac{\mu_0^{1/2}}{2C_1 M_0}\right).
\]
Within $\mathcal{V}$, the full inverse $\mathcal{G}(\mu,h)^{-1}$ is therefore bounded by $2M_0$
uniformly in $(\mu,h)$.
Since the $(1,1)$-block satisfies $\|E(\mu,h)\| \leq \|\mathcal{G}(\mu,h)^{-1}\|$,
we obtain $C_0 := 2M_0$ as the explicit uniform bound for $E$.

\emph{Why $\|E\|$ does not grow with $\|E_{-+}^{-1}\|$.}
The operator $E$ is the $(1,1)$-block of the \emph{full bounded} inverse $\mathcal{G}^{-1}$,
not of the Schur complement inverse $E_{-+}^{-1}$.
Although the Schur complement identity $E = (P_h - \omega^2 + R_- E_{-+}^{-1} R_+)^{-1}$
is formally correct, it is not the route used to bound $\|E\|$:
the Neumann series bounds $\|\mathcal{G}^{-1}\|$ at the level of the full $2\times 2$ block
operator, so all four blocks---including $E$---are simultaneously controlled by $2M_0$
(using $\|A_{11}\|\leq\|A\|$ for block operators on a Hilbert direct sum),
independently of how large $\|E_{-+}^{-1}\|$ becomes.

Throughout the double scaling regime ($|\mu| \sim h^{4/5}$, $h \to 0$), the pair $(\mu,h)$
lies inside $\mathcal{V}$ for $h$ sufficiently small: since $\mu_0$ is $h$-independent (as
established above), the condition $|\mu| < \mu_0$ reduces to $Ch^{4/5} < \mu_0$, which holds
for all $h < ({\mu_0}/{C})^{5/4}$.  The thresholds $\mu_0, h_0$ therefore provide a
genuinely uniform validity region, and the Grushin estimates remain valid throughout the
double-scaling window without further smallness conditions on $h$.
Outside $\mathcal{V}$ (i.e.\ for $|\mu| \geq \mu_0$ or $h \geq h_0$), the operator
$\mathcal{G}$ is still Fredholm---specifically, $P_h(\mu) - \omega^2$ remains Fredholm
of index zero by the analytic Fredholm theorem (Dyatlov--Zworski~\cite{DZ19},
Theorem~C.5) applied to the meromorphic family of
Lemma~\ref{lem:meromorphic_continuation}(ii)--(iii)---but the $3\times 3$ finite-dimensional
reduction is no longer sufficient and the full resolvent must be used.
\end{remark}


% -------------------------------------------------------
\subsection{WKB Suppression and Logical Independence of the Proof}
\label{sec:wkb_suppression}

\begin{proposition}[Logical independence of the WKB suppression]
\label{prop:wkb_independence}
The asymptotic vanishing $X_{31} = 0$ established in
Lemma~\ref{lem:X31_asymptotic} below depends only on the following chain
of results, none of which invokes Theorem~\ref{thm:resolvent} or the
resolvent estimates of Steps~1--2 in its proof:
\begin{enumerate}
  \item[\emph{(i)}] \emph{Jordan chain} (Lemma~\ref{lem:jordan_chain},
    Section~2): the functions $e_1, e_3$ exist, satisfy the Jordan chain
    relations, and satisfy the outgoing boundary condition on $\Gamma$.
    This is established from the Stokes sector analysis of
    $P_{\mathrm{mod}} = -\partial_y^2 + y^3$ alone.
  \item[\emph{(ii)}] \emph{WKB decay on $\Gamma$} (Lemma~\ref{lem:stokes_gm1},
    Section~2): every outgoing solution satisfies
    $|e_j(t)| \lesssim t^{-3/4}e^{-2t^{5/2}/5}$ as $t\to\infty$ on $\Gamma$.
    This follows from the asymptotic analysis of the Stokes sector for
    $-u''+t^3 u = 0$ (Sibuya~\cite{Sibuya75}), independently of any
    resolvent or parameter.
  \item[\emph{(iii)}] \emph{Product decay}: the product $e_1 e_3$ carries
    the combined decay $e^{-4t^{5/2}/5}$, which forces the integral
    $X_{31}^{(T)} = e^{i4\pi/5}\int_{t_L}^T t\,e_1 e_3\,dt$ to zero as
    $T \to \infty$.
\end{enumerate}
In particular, the logical order is:
\begin{gather*}
  \text{(i) Jordan chain}
  \;\Rightarrow\;
  \text{(ii) WKB decay}
  \;\Rightarrow\;
  \text{(iii) } X_{31} = 0 \\
  \;\Rightarrow\;
  \lambda_{2,3} \sim \mu^{1/2}
  \;\Rightarrow\;
  \|R_h\| \sim h^{-6/5}|P-P_c|^{-1/2}.
\end{gather*}
No step in the forward direction uses any result from a later step.
The proof of Theorem~\ref{thm:resolvent} therefore contains no circular
dependence.
\end{proposition}

\begin{lemma}[Asymptotic vanishing and effective suppression of $X_{31}$]
\label{lem:X31_asymptotic}
Let $X_{31}^{(T)} := e^{i4\pi/5}\int_{t_L}^{T} t\,e_1(t)\,e_3(t)\,dt$ denote the
truncated matrix element at contour length $T$.
\begin{enumerate}
\item[(i)] \emph{(Not exactly zero at finite $T$).}
  For any finite $T < \infty$, the truncated element satisfies $X_{31}^{(T)} \neq 0$.
  The statement ``$X_{31} = 0$'' throughout this paper refers to the asymptotic limit
  $T \to \infty$, not to an exact algebraic identity.
\item[(ii)] \emph{(Super-exponential decay.)}
  The WKB estimate on $\Gamma$ gives
  \[
    \bigl|X_{31}^{(T)}\bigr| \;\leq\; C\,e^{-4T^{5/2}/5}
  \]
  for a constant $C > 0$ depending only on $t_L$ and the normalization of
  $\{e_j\}$, not on $T$.  In particular,
  $X_{31}^{(T)}/c_X \to 0$ super-exponentially as $T\to\infty$.
\item[(iii)] \emph{(Physical regime: $\mu X_{31}$ is negligible.)}
  For any fixed $\mu > 0$, the crossover scale
  $\mu_{\mathrm{cross}}(T) := (X_{31}^{(T)})^2/c_X^3 \to 0$
  super-exponentially in $T$.
  Since the physical parameter $\mu = P - P_c$ is independent of $T$, we have
  $\mu \gg \mu_{\mathrm{cross}}(T)$ for $T$ large enough, so the term $\mu X_{31}^{(T)}$
  contributes at order $O(\mu_{\mathrm{cross}}^{1/2} \mu^{1/2}) \ll O(\mu^{1/2} c_X^{1/2})$
  in the eigenvalue expansion~\eqref{eq:lambdapm} and does not affect the
  dominant $\lambda \sim \mu^{1/2}$ scaling.
\end{enumerate}
\end{lemma}

\begin{proof}
Part~(i): $e_1(t) > 0$ and $e_3(t) > 0$ on $[t_L, T]$ (Steps~1 and~4 of
Lemma~\ref{lem:cx_nonzero}), so the integrand $t\,e_1(t)\,e_3(t) > 0$
on $(t_L, T)$, yielding $|X_{31}^{(T)}| > 0$.

Part~(ii): On $\Gamma$, both $e_1$ and $e_3$ satisfy the outgoing boundary condition
and hence carry the WKB decay $|e_j(t)| \sim t^{-3/4}e^{-2t^{5/2}/5}$.
Therefore $|e_1(t)\,e_3(t)| \lesssim t^{-3/2}e^{-4t^{5/2}/5}$, and the tail integral
$\int_T^\infty t\,|e_1 e_3|\,dt \leq C' e^{-4T^{5/2}/5}$
for a constant $C' > 0$.  Since the integral on $[t_L, T]$ is bounded and the
tail decays super-exponentially, the result follows.

Part~(iii): Follows directly from Part~(ii) and the bound $|c_X| \geq c_0 > 0$
(Lemma~\ref{lem:cx_nonzero}).
\end{proof}

\begin{remark}[Terminology]
\label{rem:X31_terminology}
When we write $X_{31} = 0$ in the characteristic equation analysis (e.g.\
Sections~6--7 and Appendix~B), this is shorthand for the statement that the
$T\to\infty$ limit of $X_{31}^{(T)}$ is zero, and that for any fixed $\mu > 0$ the
contribution $\mu X_{31}^{(T)}$ is negligible relative to the dominant term
$\mu c_X \lambda \sim \mu^{3/2}$ for $T$ large.  This is an \emph{asymptotic
suppression}, not an exact algebraic cancellation.  The distinction is made precise
by the crossover analysis in Part~(iii) above and Remark~\ref{rem:dominant_balance}.
\end{remark}


By Lemma~\ref{lem:X31_asymptotic}, the asymptotic vanishing
\[
  X_{31} = \langle e_3, x e_1 \rangle
  = \int_\Gamma x\, e_1(x)\, e_3(x)\, dx \;\xrightarrow{T\to\infty}\; 0
\]
is a consequence of the outgoing WKB boundary conditions: both $e_1$ and $e_3$ carry
the decaying envelope $e^{-2t^{5/2}/5}$, so their product decays as $e^{-4t^{5/2}/5}$,
which is not countered by the polynomial weight $x \sim t$.
Since the contour $\Gamma = \{te^{i2\pi/5} : t > 0\}$ is a ray rather than a symmetric
domain, classical parity-based cancellation arguments do not apply here; the suppression
is an asymptotic WKB effect (see Remark~\ref{rem:X31_terminology}).
Under the parametrization $x = te^{i2\pi/5}$, we have
$X_{31} = e^{i4\pi/5}\int_0^\infty t\,e_1(t)\,e_3(t)\,dt$,
consistent with the representation for $X_{jk}$ in Appendix~C.
By contrast, the cross terms $e_1 e_2$ and $e_2 e_3$ carry
different net WKB phases---one factor lies in a decaying Stokes sector,
while the other is in an intermediate sector---which yields
non-zero sub-diagonal elements $c_X = X_{21} + X_{32} \neq 0$,
as established by the following lemma.

\begin{lemma}[$c_X \neq 0$: analytic proof, corrected]
\label{lem:cx_nonzero}
Let $\{e_1, e_2, e_3\}$ be the Jordan chain established in Lemma~\ref{lem:jordan_chain}.
Then $c_X = X_{21} + X_{32} \neq 0$.
\end{lemma}

\begin{proof}
We establish $X_{21} \neq 0$ and $X_{32} \neq 0$ separately, then
show $c_X \neq 0$.

\medskip
\noindent\textbf{Step 1: Real-phase structure of $e_1$.}

The equation $-u'' + t^3 u = 0$ has real coefficients. Its WKB
solutions satisfy
\[
  v_\pm(t) \;\sim\; t^{-3/4}\exp\!\Bigl(\pm\tfrac{2}{5}t^{5/2}\Bigr)
  \bigl(1 + O(t^{-5/2})\bigr), \qquad t \to +\infty.
\]
The outgoing mode $e_1 = v_+$ is selected by integrating backward
from $t = T$ with initial data
$e_1(T) = T^{-3/4}e^{-2T^{5/2}/5}$ and
$e_1'(T) = -T^{3/2}\,T^{-3/4}e^{-2T^{5/2}/5}(1+O(T^{-5/2}))$,
both of which are real. Since the ODE has real coefficients and the
initial data are real, the solution $e_1(t)$ is \emph{real-valued}
for all $t \in [t_L, T]$. Positivity $e_1(t) > 0$ follows because
$e_1(T) > 0$ and the Stokes structure of $-u''+t^3u=0$ guarantees
that $v_+$ is nodeless on the positive real axis
(Sibuya~\cite{Sibuya75}, Chapter~2). In particular,
\begin{equation}\label{eq:e1_real}
  e_1(t) \in \mathbb{R}_{>0} \quad \text{for all } t \in [t_L, T].
\end{equation}

\medskip
\noindent\textbf{Step 2: Sign of the Wronskian.}

The growing mode $v_-$ satisfies $v_-(T) = T^{-3/4}e^{+2T^{5/2}/5}(1+O(T^{-5/2})) > 0$.
The Wronskian $W := W(e_1, v_-)\big|_T$ is \emph{independent of $T$}
by Abel's identity: since the equation $-u''+t^3u=0$ has no first-derivative
term, its coefficient of $u'$ is zero, so
$\tfrac{d}{dt}W(e_1,v_-) = 0$ identically, and $W$ is the same constant
for every evaluation point $T > t_L$.
Substituting the WKB asymptotics at $T\gg 1$:
\[
  W = e_1(T)\,v_-'(T) - e_1'(T)\,v_-(T)
    = 2T^{-1/2}(1 + O(T^{-5/2})),
\]
so $W > 0$ for $T$ sufficiently large. (Numerical verification across
$T \in [2.5, 4.0]$ gives $W \approx 0.847$ stable to $2\times 10^{-4}$;
see Appendix~C, Table~C.3.)

\medskip
\noindent\textbf{Step 3: Positivity of $X_{21}$.}

Under the coordinate change $x = te^{i2\pi/5}$ on $\Gamma$, the
matrix element is
\[
  X_{21} = e^{i4\pi/5}\int_{t_L}^\infty t\,e_1(t)\,e_2(t)\,dt.
\]
From the variation of parameters formula:
\[
  e_2(t) = \frac{e_1(t)}{W}\int_{t_L}^t s\,e_1(s)^2\,ds
          =: \frac{e_1(t)}{W}\,I_1(t).
\]
Since $e_1 > 0$ and $W > 0$, we have $e_2(t) > 0$. Setting
$A_1 := \int_{t_L}^\infty t\,e_1(t)^2\,I_1(t)\,dt > 0$,
\[
  X_{21} = \frac{e^{i4\pi/5}\,A_1}{W} \neq 0.
\]

\medskip
\noindent\textbf{Step 4: Positivity of $X_{32}$.}

Applying variation of parameters once more, with source
$\mathrm{src}_2 = e^{i4\pi/5}\,t\,e_2(t)$:
\[
  e_3(t) = \frac{e_1(t)}{W}\int_{t_L}^t s\,e_2(s)\,e_1(s)\,ds
          - \frac{v_-(t)}{W}\int_{t_L}^t s\,e_2(s)\,e_1(s)^2/W\,ds + \cdots
\]
Substituting $e_2 = e_1 I_1/W$ and collecting, the outgoing projection gives:
\[
  e_3(t) = \frac{e_1(t)}{W^2}\,I_2(t), \qquad
  I_2(t) := \int_{t_L}^t s\,e_1(s)^2\,I_1(s)\,ds > 0.
\]
Here $I_2 > 0$ since the integrand is strictly positive ($e_1 > 0$, $I_1 \geq 0$
with $I_1$ strictly increasing, $W > 0$), so $e_3(t) > 0$.
Setting $A_2 := \int_{t_L}^\infty t\,e_1(t)^2\,I_1(t)\,I_2(t)\,dt > 0$
and using $e_2 = e_1 I_1/W$, $e_3 = e_1 I_2/W^2$:
\[
  X_{32} = e^{i4\pi/5}\int_{t_L}^\infty t\,e_2(t)\,e_3(t)\,dt
         = e^{i4\pi/5}\int_{t_L}^\infty t\,\frac{e_1 I_1}{W}\cdot\frac{e_1 I_2}{W^2}\,dt
         = \frac{e^{i4\pi/5}\,A_2}{W^3} \neq 0.
\]

\medskip
\noindent\textbf{Step 5: Non-vanishing of $c_X$.}

Since $W > 0$, $A_1 > 0$, and $A_2 > 0$:
\[
  c_X = X_{21} + X_{32}
      = e^{i4\pi/5}\!\left(\frac{A_1}{W} + \frac{A_2}{W^3}\right).
\]
The key observation is that \emph{both $X_{21}$ and $X_{32}$ carry
the same phase factor $e^{i4\pi/5}$}, so their sum is
$e^{i4\pi/5}$ times a positive real number; the combination is
not two independent complex numbers whose cancellation must be
ruled out.  Concretely, the quantity in parentheses equals
$A_1 W^{-1} + A_2 W^{-3} > 0$ since each term is strictly
positive (Steps~1--4), and $e^{i4\pi/5} \cdot r \neq 0$ for
any $r > 0$.  Since $\arg(c_X) = 4\pi/5$, in particular
$c_X \notin \mathbb{R}_{<0}$ and $c_X \notin \{0\}$:
we conclude $c_X \neq 0$, with the explicit phase
$c_X = |c_X|\,e^{i4\pi/5}$.
\end{proof}

Numerical verification of this result is presented in Appendix~C, which confirms both the sign
structure and the exponential decay of the ratio $|X_{31}|/|c_X|$ with respect to~$T$.
Under the suppression condition $X_{31} = 0$, the characteristic equation reduces to:
\[
  \lambda(\lambda^2 - c_X \mu) = 0
  \implies \lambda_1 = 0,\quad \lambda_{2,3} = \pm\sqrt{c_X \mu}.
\]
(The full determinant analysis, including the suppression of the trace and quantum error terms,
is detailed in Appendix~B.)

\begin{remark}[Structural Rigidity and Phase Robustness]
The suppression of the matrix element $X_{31} = \langle e_3, x\, e_1 \rangle$ is a robust
consequence of the WKB phase structure rather than an accidental cancellation.
Two numerical tests verify this property.

\emph{Operator sensitivity.}
Replacing the operator $P_{\mathrm{mod}} = -\partial_x^2 + x^3$ with
$\tilde{P} = -\partial_x^2 + x^2$ and rebuilding the Jordan chain independently
(Stokes angle $\theta = \pi/3$, source phase $e^{i2\pi/3}$) gives a suppression
ratio $|\tilde{X}_{31}|/|\tilde{c}_X| \approx 0.33$ at $T = 3.2$,
$0.14$ at $T = 3.5$, and $0.15$ at $T = 4.0$ --- values of order $O(1)$ with no
significant decay. By contrast, the $x^3$ operator yields
$|X_{31}|/|c_X| \approx 1.98\times10^{-3}$ at $T=3.2$,
$8.51\times10^{-5}$ at $T=3.5$, and $1.22\times10^{-7}$ at $T=4.0$.
The ratio for $x^3$ is thus smaller than the $x^2$ ratio by a factor of
$1.3\times10^6$ at $T=4.0$, confirming that the suppression is specific to the WKB
phase pairing induced by the odd potential $x^3$, and not a generic feature of the
Jordan chain structure.

\emph{Phase robustness.}
Under a rotation of the contour angle by $\delta = 0.05$ rad
(i.e.\ $\theta \to \theta + \delta$ with the full ODE coefficient
$e^{i5(\theta+\delta)} t^3$ updated accordingly), the suppression ratio
$|X_{31}|/|c_X|$ remains at the same asymptotic scale as the standard contour:
$9.85\times10^{-5}$ (rotated) vs $8.51\times10^{-5}$ (standard) at $T = 3.5$,
and $1.49\times10^{-7}$ (rotated) vs $1.22\times10^{-7}$ (standard) at $T = 4.0$.
The suppression is thus robust against small phase deformations.

\emph{WKB decay rate.}
A linear fit to the data of Table~C.1 yields
\[
  \log|X_{31}/c_X| \approx -0.711\,T^{5/2} + 6.84,
\]
consistent with the theoretical WKB bound
$\log|X_{31}/c_X| \leq -\frac{4}{5}T^{5/2} + C$ (Appendix~C, Remark~C.3).
The fit residuals are within $8\%$ across $T \in [3.2, 4.0]$, confirming
the super-exponential asymptotic rate.

A dominant balance analysis confirming that the term $\mu X_{31}$ is negligible
in the physical regime $\mu \gg \mu_{\mathrm{cross}}$ is provided in
Remark~\ref{rem:dominant_balance} (Appendix~C).
\end{remark}

% -------------------------------------------------------
\section{Double Scaling Regime}
% -------------------------------------------------------

The two independent limits $h \to 0$ and $P \to P_c$ interact non-trivially when
$|P - P_c| \sim h^{4/5}$. In this regime, neither limit dominates,
and the semiclassical and geometric singularities must be treated simultaneously.

\subsection{Rescaled Operator}

The unfolded model operator is
\[
  P_{\mathrm{mod}}(\mu) = -h^2\partial_x^2 + x^3 + \mu x,
  \qquad \mu = P - P_c.
\]
This represents the standard $A_2$ (cusp) unfolding of $-h^2\partial_x^2 + x^3$,
with $\mu x$ acting as the single unfolding parameter (see Arnold~\cite{Arnold72}).
Under the dilation $x = h^{2/5}y$, the kinetic and cubic terms balance as detailed in
Section~4, and the linear perturbation transforms as:
\[
  P_{\mathrm{mod}}(\mu)
  = h^{6/5}\Bigl(-\partial_y^2 + y^3 + \nu\, y\Bigr)
  =: h^{6/5}\, Q(\nu),
  \qquad
  \nu := \mu\, h^{-4/5}.
\]
The rescaled operator $Q(\nu) = -\partial_y^2 + y^3 + \nu y$ is $h$-independent.
The \emph{double scaling regime} is defined by $\nu = O(1)$, i.e.\ $\mu \sim h^{4/5}$,
which is precisely the region where the cubic and linear terms in $Q(\nu)$ are
of the same order.

\subsection{Spectral Theory of $Q(\nu)$}

For each fixed $\nu \in \mathbb{R}$, the operator $Q(\nu)$ defined on $L^2(\mathbb{R})$
subject to outgoing boundary conditions on the contour $\Gamma$ possesses a discrete
spectrum of resonances $\{\varepsilon_n(\nu)\}_{n \geq 1}$.
The resonance eigenvalues $\varepsilon_n(\nu)$ depend smoothly on $\nu$
and satisfy the uniform bound $\varepsilon_n(\nu) = O(1)$ for $\nu$ in any bounded set;
this is established in Appendix~E (Proposition~\ref{prop:Qnu_O1})
by means of WKB and Stokes sector analysis.
At $\nu = 0$, $Q(0) = -\partial_y^2 + y^3$ is the cusp model of
Lemma~\ref{lem:stokes_gm1}, whose resolvent at energy zero is bounded.
As $|\nu| \to \infty$, the spectrum transitions continuously to the
standard Airy regime.

The bifurcation structure of $Q(\nu)$ reflects the underlying $A_2$ catastrophe.
For $\nu > 0$, the potential $y^3 + \nu y$ has no real critical points, and the
resonances are strictly non-real. For $\nu < 0$, two real turning points
emerge, and the resonances split into pairs that are symmetric about the imaginary axis.

\subsection{Resolvent Bound in the Double Scaling Window}
\label{subsec:resolvent_bound_ds}

Since $P_{\mathrm{mod}}(\mu) = h^{6/5}Q(\nu)$, the resolvent satisfies
\[
  R_h(\mu) = \bigl(P_{\mathrm{mod}}(\mu)\bigr)^{-1}
  = h^{-6/5}\,Q(\nu)^{-1}.
\]
For $\nu = O(1)$ (i.e.\ $\mu \sim h^{4/5}$) and under the spectral condition $0 \notin \mathrm{spec}(Q(\nu))$,
the uniform bound $\|Q(\nu)^{-1}\| = O(1)$ holds throughout the regime by
Corollary~\ref{cor:Qnu_resolvent} (Appendix~E), which yields:
\[
  \|R_h(\mu)\| = h^{-6/5}\,\|Q(\nu)^{-1}\| = O(h^{-6/5})
  \qquad (\nu = O(1)).
\]
There is no $|\nu|^{-1/2}$ divergence of $\|Q(\nu)^{-1}\|$ as $\nu \to 0$
within this regime. Indeed, by Remark~\ref{rem:specgap_nu0}, we have $0 \notin \mathrm{spec}(Q(\nu))$
for all $|\nu| \leq C$, implying that $(Q(\nu))^{-1}$ is uniformly bounded on $[-C,C]$.
The factor $|\nu|^{-1/2}$ appearing in Theorem~\ref{thm:resolvent} is
associated with the finite-dimensional Schur complement $E_{-+}(\nu)$ at fixed $h$.
In the double scaling regime, $\nu$ and $h$ are coupled, and the full operator
$Q(\nu)$ regularizes the singularity, remaining uniformly invertible.

Consequently, substituting $|\mu| = |\nu| h^{4/5}$ with $|\nu| \leq C$, we obtain:
\[
  \|R_h(\mu)\| \leq C_0\, h^{-6/5}
  \qquad \text{uniformly for } |\nu| \leq C.
\]
\emph{Matching lower bound.}
We establish the lower bound $\|R_h(\mu)\| \geq c\,h^{-6/5}$ with the same semiclassical
scaling, showing that the upper bound is sharp in the double scaling regime.
Since $R_h(\mu) = h^{-6/5} Q(\nu)^{-1}$, it suffices to show that
$\|Q(\nu)^{-1}\| \geq c > 0$ uniformly for $|\nu| \leq C$.
Because $Q(\nu)$ is a closed operator with compact resolvent on the weighted $L^2$ space,
its spectrum consists of discrete resonances. For a fixed $\nu \in [-C,C]$, let
$\varepsilon_1(\nu)$ be the resonance closest to the origin. By Remark~\ref{rem:specgap_nu0},
we have $0 \notin \mathrm{spec}(Q(\nu))$, which implies $\varepsilon_1(\nu) \neq 0$.
Furthermore, by Proposition~\ref{prop:Qnu_O1}, we have $|\varepsilon_1(\nu)| \leq C_{\mathrm{spec}}$
uniformly for some constant $C_{\mathrm{spec}} > 0$.
The norm of the resolvent is bounded below by the inverse of the spectral distance to the origin:
\[
  \|Q(\nu)^{-1}\| \geq |\varepsilon_1(\nu)|^{-1} \geq C_{\mathrm{spec}}^{-1} =: c > 0.
\]
Combining this with the semiclassical scaling yields:
\[
  \|R_h(\mu)\| = h^{-6/5}\|Q(\nu)^{-1}\| \geq c\,h^{-6/5}
  \qquad \text{uniformly for } |\nu| \leq C.
\]
Together with the upper bound, this yields the two-sided estimate:
\[
  c\,h^{-6/5} \leq \|R_h(\mu)\| \leq C_0\,h^{-6/5}
  \qquad \text{uniformly for } |\mu| \leq C h^{4/5}.
\]
\begin{remark}[Consistency with Theorem~\ref{thm:resolvent}]
\label{rem:consistency}
Evaluating the formula in Theorem~\ref{thm:resolvent} by substituting $|P - P_c| \sim h^{4/5}$
would yield $\|R_h\| \sim h^{-6/5} h^{-2/5} = h^{-8/5}$. Such a
substitution is invalid because Theorem~\ref{thm:resolvent} is derived at a fixed deformation $\mu$
as $h \to 0$, whereas the double scaling regime couples $\mu$ and $h$
via $\mu = \nu h^{4/5}$ with $\nu = O(1)$.

In this coupled regime, the resolvent is analyzed directly via the rescaled operator
$Q(\nu) = -\partial_y^2 + y^3 + \nu y$ rather than the Schur complement eigenvalues
of Theorem~\ref{thm:resolvent}.
The uniform bound $\|Q(\nu)^{-1}\| = O(1)$ for $\nu$ in a compact neighborhood
of $0$ follows from Corollary~\ref{cor:Qnu_resolvent} (Appendix~E).
Since the resonances $\varepsilon_n(\nu)$ are bounded and bounded away from zero,
the resolvent norm at zero energy is controlled by the spectral gap.
This is a property of the $L^2$-operator norm of $Q(\nu)^{-1}$, whereas the Schur
complement eigenvalues $\lambda_j(\nu)$ characterize the deviation of $E_{-+}(\nu)$
from the nilpotent Jordan block. Since these quantities represent different spectral
properties, no contradiction arises.

In the double scaling regime, the net semiclassical scaling remains $h^{-6/5}$,
which is identical to the behavior at the critical point $P = P_c$.
Consequently, the double scaling window does not produce a stronger divergence;
it represents the regime where the semiclassical and parameter-bifurcation effects balance.
\end{remark}

\subsection{Resonance Spacing}

The physical resonance energies satisfy $E_n = h^{6/5}\,\varepsilon_n(\nu)$.
Since $\varepsilon_n(\nu) = O(1)$ uniformly in the double scaling regime, consecutive
resonances are spaced by $\Delta E_n = h^{6/5}\,\Delta\varepsilon_n$, where
$\Delta\varepsilon_n = O(1)$.

To express this in terms of the physical frequencies, we recall that in the vicinity of $\omega_c$:
\[
  E = \omega^2 - \omega_c^2 \approx 2\omega_c\,\Delta\omega,
\]
where the critical frequency $\omega_c$ is determined solely by the geometry at $P_c$
and is independent of~$h$ (see Section~4).
This yields:
\[
  \Delta\omega \sim \frac{\Delta E_n}{2\omega_c}
  \sim \frac{h^{6/5}\,\Delta\varepsilon_n}{2\omega_c}
  \sim h^{6/5},
\]
which constitutes the fundamental resonance spacing in the double scaling regime.

\subsection*{Further Directions}

\textbf{Kerr--de~Sitter generalization.}
The Regge--Wheeler operator on SdS is spherically symmetric; the analogous problem
on Kerr--de~Sitter (KdS) backgrounds introduces a non-trivial angular momentum
parameter $a$. The trapped set on KdS is a normally hyperbolic submanifold,
and the spectral gap $\kappa_{\mathrm{gap}}(P,a) := \sqrt{-2V_{\mathrm{eff}}''(r_c(P,a))}$
depends on both $P$ and $a$.
A natural question is whether the simultaneous vanishing
$\kappa_{\mathrm{gap}}(P,a) = 0$ is achievable for parameters $(P,a)$ in the physical regime,
and whether the analytic continuation argument developed here extends to the KdS setting.
The Grushin framework of Dyatlov--Zworski~\cite{DZ19} remains applicable,
and the $A_2$ normal form analysis is expected to carry over, though the
multiparameter unfolding and the WKB suppression mechanism require
independent verification. The asymptotic distribution of QNMs for KdS black holes
is treated in Dyatlov~\cite{Dyatlov12}, and global quasilinear stability in this
geometry is established in Hintz--Vasy~\cite{HintzVasy2016}.
The degenerate analysis presented here complements these results at the boundary
of the normally hyperbolic regime. Further details on the semiclassical machinery
employed in this work can be found in Zworski~\cite{Zworski2012}.

\textbf{Higher-order degeneracies.}
The present paper treats the $k=3$ ($A_2$, cusp) degeneracy, which is the simplest
codimension-1 case in Arnold's classification~\cite{Arnold72}.
The lower case $k=2$ ($A_1$, Airy) is analogous: the model operator $-h^2\partial_x^2 + x^2$
near a nondegenerate maximum gives a Jordan block of size 2, and the corresponding
resolvent bound $\|R_h\| \sim h^{-1}|P-P_c|^{-1}$ follows from the standard
normally hyperbolic trapping theory~\cite{WZ11,NZ15};
see also Zworski~\cite{Zworski2012} for the semiclassical framework.
The next case, $k=4$ ($A_3$, swallowtail), requires the simultaneous vanishing
$V'_{\mathrm{eff}} = V''_{\mathrm{eff}} = V'''_{\mathrm{eff}} = 0$ and is
codimension-2 in the parameter space.
The model operator is given by $-h^2\partial_x^2 + x^4$, which scales as $h^{8/5}$ under
the dilation $x = h^{2/5}y$, suggesting a resolvent bound of the form
$\|R_h\| \sim h^{-8/5}|P-P_c|^{-\beta}$ for an exponent $\beta$ determined
by the Schur complement of the corresponding $4\times 4$ Jordan block.
Whether the WKB suppression mechanism generalizes---namely, whether the dominant
off-diagonal matrix element $X_{41}$ vanishes asymptotically---remains an open
question that depends on the Stokes sector structure of $-u'' + x^4 u = 0$.

\textbf{Observational signatures.}
The resonance spacing $\Delta\omega \sim h^{6/5}$ in the double scaling regime
differs qualitatively from the spacing $\Delta\omega \sim h$ in the normally
hyperbolic regime. To make this concrete, recall that in the Regge--Wheeler
context the semiclassical parameter is identified as $h = \ell^{-1}$, where
$\ell$ is the angular momentum mode number (the large-$\ell$ limit being the
relevant semiclassical regime). Alternatively, for a black hole of mass $M$, the
natural scale is $h \sim (M\omega_{\mathrm{QNM}})^{-1}$.
Under either identification, the modified spacing law $\Delta\omega \sim h^{6/5}$
produces a fractional deviation $\Delta\omega/\Delta\omega_{\mathrm{NH}} \sim h^{1/5}$
from the normally hyperbolic prediction, which grows as $h \to 0$ (i.e.\ at large
$\ell$ or large $M$). Near-future gravitational wave detectors such as LISA or the
Einstein Telescope, operating at high signal-to-noise ratios in the ringdown phase,
may in principle be sensitive to such deviations for nearly-extremal SdS configurations.
Making this connection fully precise---including the relationship between the
double-scaling window $|\mu| \sim h^{4/5}$ and the physical nearness to the Nariai
boundary---is a direction left for future work.

% -------------------------------------------------------
% APPENDICES
% -------------------------------------------------------
\appendix

% -------------------------------------------------------
\section{Sub-Principal Symbol: $Q_1 \equiv 0$}
% -------------------------------------------------------

\subsection*{A.1.~Context}

In the initial step of the Birkhoff normal form (BNF) procedure, the FIO conjugation
described in Section~4 yields the identity~\eqref{eq:BNF_conjugation}. We expand this as
\[
  U_h^* P_h U_h = -h^2\partial_x^2 + x^3 + h\,Q_1(x,\xi;h) + O(h^2),
\]
where $Q_1$ is the sub-principal error term. If $Q_1$ were non-vanishing, an additional homological
equation would need to be solved at the zeroth BNF step, and an $O(h)$ term would
enter the Grushin matrix. The following lemma shows that $Q_1$ vanishes identically for
the Regge--Wheeler operator.

\subsection*{A.2.~Lemma}

\begin{lemma}[Uniqueness of $a_0$]
\label{lem:a0_unique}
Let $a_0 \in C^\infty(\operatorname{supp}\chi_{\mathrm{trap}})$ satisfy
\begin{equation}\label{eq:transport}
  H_{p_{\mathrm{mod}}} a_0 = 0
  \qquad \text{on } \operatorname{supp}\chi_{\mathrm{trap}},
\end{equation}
where $H_{p_{\mathrm{mod}}} = 2\eta\,\partial_x - 3x^2\,\partial_\eta$,
with the normalization $a_0(0,0) = 1$.
Then $a_0 \equiv 1$ on $\operatorname{supp}\chi_{\mathrm{trap}}$.
\end{lemma}

\begin{proof}
Equation~\eqref{eq:transport} is a first-order linear homogeneous PDE.
Its characteristics are the integral curves of $H_{p_{\mathrm{mod}}}$,
i.e.\ the solutions $t \mapsto (x(t), \eta(t))$ of
\[
  \dot x = 2\eta, \qquad \dot\eta = -3x^2.
\]
Along each characteristic, $a_0$ is constant by the method of characteristics.
It suffices to show that every point of $\operatorname{supp}\chi_{\mathrm{trap}}$
lies on a characteristic that passes through (or accumulates at) $(0,0)$.

The critical point $(0,0)$ of $H_{p_{\mathrm{mod}}}$ is the unique point of
$T^*\mathbb{R}$ where $H_{p_{\mathrm{mod}}} = 0$.
The cutoff $\chi_{\mathrm{trap}}$ is supported in a conic neighborhood
$\mathcal{U}$ of the trapped set $K_P = \{(0,0)\}$
(Appendix~\ref{app:escape}, Proposition~\ref{prop:global_escape}).
We split $\operatorname{supp}\chi_{\mathrm{trap}}$ into two parts.

\emph{Points on the zero-energy surface $\{p_{\mathrm{mod}} = 0\}$.}
On $\{\eta^2 + x^3 = 0\}$, we have $x \leq 0$.
The Hamiltonian flow on this invariant surface has $(0,0)$ as its unique
$\omega$-limit point: the conserved quantity $p_{\mathrm{mod}} = \eta^2 + x^3 = 0$
confines the orbit to $x = -\eta^{2/3} \to 0^-$ and $\eta \to 0$, so
every orbit on $\{p = 0\} \cap \mathcal{U}$ accumulates at $(0,0)$ as $t \to +\infty$.

\emph{Points off the zero-energy surface.}
For $\rho \in \operatorname{supp}\chi_{\mathrm{trap}}$ with $p_{\mathrm{mod}}(\rho) \neq 0$,
we choose $\mathcal{U}$ small enough (depending only on the support of $\chi_{\mathrm{trap}}$)
so that every orbit starting in $\mathcal{U}$ is captured: specifically, by the escape
function estimate of Proposition~\ref{prop:global_escape}, any trajectory that enters
$\mathcal{U}$ from outside must exit $\mathcal{U}$ in finite backward time without
passing through $(0,0)$, so the backward orbit of $\rho$ exits $\mathcal{U}$ and
enters a region where $\chi_{\mathrm{trap}} = 0$; equivalently, the forward orbit
of any such $\rho$ entered $\mathcal{U}$ from outside.
However, because $\mathcal{U}$ is an \emph{invariant neighborhood} of $K_P$ for the
trapped set (i.e.\ every orbit in $\mathcal{U}$ with $p = 0$ stays in $\mathcal{U}$),
and since $p_{\mathrm{mod}}$ is a flow-invariant of $H_{p_{\mathrm{mod}}}$,
any orbit starting at $\rho \in \operatorname{supp}\chi_{\mathrm{trap}}$
with $p_{\mathrm{mod}}(\rho) \neq 0$ must exit $\mathcal{U}$ in finite forward time
(as the escape function $H_p G < 0$ on $\mathcal{U} \setminus K_P$).
For such a trajectory the equation $H_{p_{\mathrm{mod}}} a_0 = 0$ propagates
the value $a_0(\rho)$ backward along the orbit; the orbit crosses a boundary
of $\operatorname{supp}\chi_{\mathrm{trap}}$ where the transport equation
no longer constrains $a_0$. The boundary condition $a_0(0,0) = 1$ together
with the normalization of $U_h$ (which fixes $a_0$ on the zero-energy surface
through $\operatorname{supp}\chi_{\mathrm{trap}}$) then propagates to
a tubular neighborhood of $K_P$ by continuity.

In both cases, continuity of $a_0$ and the boundary value $a_0(0,0) = 1$
give $a_0 \equiv 1$ on $\operatorname{supp}\chi_{\mathrm{trap}}$.
\end{proof}

\begin{lemma}[Vanishing of the sub-principal symbol]
\label{lem:Q1}
Let $P_h = -h^2\partial_{r_*}^2 + V_{\mathrm{eff}}(r_*)$, where $V_{\mathrm{eff}}$ is the
Regge--Wheeler effective potential and $r_*$ is the tortoise coordinate. Let
$F: T^*\mathbb{R} \to T^*\mathbb{R}$ be the cotangent lift defined along the Hamiltonian flow of
the tortoise coordinate. In Weyl quantization, the sub-principal symbol
\[
  q_1(x,\xi) :=
  \frac{1}{2i}\bigl\{p_2,\, a_0\bigr\}(x,\xi)
  + \bigl(p^{\mathrm{orig}}_{\mathrm{sub}} \circ F^{-1}\bigr)(x,\xi) \equiv 0
\]
vanishes identically on $\operatorname{supp}\chi_{\mathrm{trap}}$. Consequently,
\[
  U_h^* P_h U_h = -h^2\partial_x^2 + x^3 + O(h^2)
\]
at the level of the principal symbol; the $hQ_1$ term is absent.
\end{lemma}

\begin{proof}
We decompose the proof into three parts.

\textbf{Step~1: The original sub-principal symbol satisfies $p^{\mathrm{orig}}_{\mathrm{sub}} \equiv 0$.}
In tortoise coordinates, the operator takes the form
$P_h = -h^2\partial_{r_*}^2 + V_{\mathrm{eff}}(r_*)$.
Since this expression contains no first-order derivative terms, the Weyl
symbol calculus yields
$P_h = \mathrm{Op}^w_h\bigl(p_2 + h\,p^{\mathrm{orig}}_{\mathrm{sub}} + O(h^2)\bigr)$,
with $p^{\mathrm{orig}}_{\mathrm{sub}}(r_*,\xi) = 0$. Indeed, in Weyl quantization, the sub-principal symbol
corresponds to $i/2$ times the coefficient of the first-order derivative term, which is
absent in~$P_h$.

\textbf{Step~2: Vanishing of the Jacobian contribution.}
The cotangent lift $F$ is a symplectic diffeomorphism of the cotangent bundle:
$(r_*,\xi) \mapsto (x,\eta)$ with $x = x(r_*)$ and $\eta = (dx/dr_*)^{-1}\xi$. Since
this transformation is a symplectomorphism, its Jacobian is identically 1 (as
$\partial x/\partial r_* \cdot \partial\eta/\partial\xi = 1$), yielding:
\[
  \log\bigl|\det(dF)\bigr| = 0 \implies \text{Jacobian contribution} = 0.
\]

\textbf{Step~3: $a_0 \equiv 1$ is the unique smooth solution of the transport equation.}
The principal symbol of the FIO amplitude satisfies
$a_0(x(r),\eta;0) = |x'(r)|^{-1/2}$ (Section~4, Step~2).
In BNF coordinates, $x'(r^c) = C^{1/3}$ (real cube root) with
$C = \tfrac{1}{6}V'''_{\mathrm{eff}}(r^c) < 0$,
so the normalization of $U_h$ fixes $a_0(0,0) = |C|^{-1/3}$.
By rescaling $a_0 \mapsto |C|^{1/3} a_0$ we normalize $a_0(0,0) = 1$,
and Lemma~\ref{lem:a0_unique} then gives $a_0 \equiv 1$
on $\operatorname{supp}\chi_{\mathrm{trap}}$, which yields:
\[
  \frac{1}{2i}\{p_2, a_0\} = \frac{1}{2i}\{p_2, 1\} = 0.
\]

\textbf{Conclusion.} {\looseness=-1 Since all contributions vanish, we conclude that
$q_1(x,\xi) = 0$ on $\operatorname{supp}\chi_{\mathrm{trap}}$.}
\end{proof}

\begin{remark}
This result does not appear as a named theorem in Dyatlov--Zworski~\cite{DZ19}; it is
standard in the microlocal analysis literature. However, because it is critical for
establishing that the first non-trivial quantum correction to the Grushin matrix is of order $O(h^2)$,
it is recorded as a lemma here.
\end{remark}

\begin{remark}[Order of the first quantum error]
Since $Q_1 \equiv 0$, the first non-trivial quantum contribution to $E_{-+}(\mu,h)$ appears
only at order $h^2$:
\[
  E_{-+}(\mu,h) = J_3 + \mu X + h^2 Q_2 + O(h^3, \mu^2, \mu h^2),
\]
where $Q_2$ arises from the Schwarzian derivative in the BNF diffeomorphism
(see Appendix~B, Subsection~B.2).
\end{remark}

\begin{lemma}[Operator norm bound on $E_{-+}(\mu,h)$]
\label{lem:E_bound}
There exists a constant $C_0 = C_0(\mu_0, h_0) > 0$, depending on $(\mu_0, h_0)$ only through
the Jordan basis norms $\|e_j\|_{L^2(\Gamma)}$ and the Schwarzian coefficient
$|\mathcal{S}_x(r^c)|$, such that
\[
  \|E_{-+}(\mu,h)\|_{\mathcal{M}_{3\times 3}} \;\leq\; C_0
  \qquad \text{for all } (\mu,h) \in \mathcal{V},
\]
where $\mathcal{V}$ is the validity region of Remark~\ref{rem:grushin_validity}.
Explicitly,
\[
  C_0(\mu_0, h_0) := \|J_3\| + \mu_0\|X\| + h_0^2\|Q_2\| + \delta_0,
\]
{\looseness=-1 with $\|J_3\| = 1$, $\|X\| \leq C_X := 3\max_{j,k}|X_{jk}|$,
$\|Q_2\| \leq C_Q$, and $\delta_0 > 0$ bounding the
$O(\mu^2, \mu h^2, h^3)$ remainder.
The $\mu$- and $h$-dependencies are \emph{linear}: $C_0$ grows
at most as $\mu_0 + h_0^2$ for small $(\mu_0, h_0)$.}
\end{lemma}

\begin{proof}
By the expansion $E_{-+}(\mu,h) = J_3 + \mu X + h^2 Q_2 + R$, where the remainder
satisfies $\|R\| \leq C_R(\mu^2 + \mu h^2 + h^3)$ for some $C_R > 0$, the triangle
inequality yields
\[
  \|E_{-+}(\mu,h)\| \leq \|J_3\| + |\mu|\,\|X\| + h^2\|Q_2\| + C_R(\mu_0^2 + \mu_0 h_0^2 + h_0^3).
\]
Setting $\delta_0 := C_R(\mu_0^2 + \mu_0 h_0^2 + h_0^3)$ gives the claimed bound with
$\|J_3\| = 1$ and $\|Q_2\|$ fixed geometric constants.

It remains to show $\|X\| < \infty$.
The entries $X_{jk} = e^{i4\pi/5}\int_0^\infty t\,e_j(t)\,e_k(t)\,dt$
are defined on the non-compact contour $\Gamma$, so a direct $L^\infty$ bound on $x$
is unavailable.
Instead, split the integral at the truncation point $T_0 := t_L \vee 1$:
\[
  |X_{jk}|
  \leq \int_{t_L}^{T_0} t\,|e_j||e_k|\,dt
     + \int_{T_0}^\infty t\,|e_j(t)||e_k(t)|\,dt
  =: I_1 + I_2.
\]
\emph{Bound on $I_1$.}
On the compact interval $[t_L, T_0]$, the functions $e_j$ are smooth solutions
of a regular ODE and hence bounded: $\sup_{[t_L,T_0]}|e_j| \leq M_j < \infty$.
Therefore
\[
  I_1 \leq T_0^2\,M_j M_k \cdot (T_0 - t_L) < \infty.
\]

\emph{Bound on $I_2$.}
Both $e_j$ and $e_k$ satisfy the outgoing boundary condition and hence the
WKB decay estimate $|e_j(t)| \leq A_j\,t^{-3/4}e^{-2t^{5/2}/5}$ for $t \geq T_0$
(Lemma~\ref{lem:stokes_gm1} and its analogue for $e_2, e_3$ from
Lemma~\ref{lem:jordan_chain}). Therefore
\[
  I_2
  \leq A_j A_k \int_{T_0}^\infty t^{-1/2} e^{-4t^{5/2}/5}\,dt
  < \infty,
\]
since the integrand decays super-exponentially.
Combining, $|X_{jk}| \leq I_1 + I_2 < \infty$ for each $(j,k)$, so
$\|X\| \leq C_X := 3\max_{j,k}|X_{jk}| < \infty$.
\end{proof}


% -------------------------------------------------------
\section{Full Determinant Analysis: Suppression of $\mathrm{Tr}(X)$ and $h^2 Q_{31}$}
\label{app:B}
% -------------------------------------------------------

\subsection*{B.1.~Full perturbation matrix}

Since $Q_1 \equiv 0$ (Lemma~\ref{lem:Q1}), the Schur complement matrix admits the expansion
\[
  E_{-+}(\mu,h) = J_3 + \mu X + h^2 Q + O(\mu^2, \mu h^2, h^4),
\]
where $\mu = P - P_c$, $J_3$ is the $3\times 3$ Jordan block at the degenerate point,
$X_{ij} = \langle e_i, x\,e_j\rangle$ is the geometric perturbation matrix, and
$Q_{ij} = \langle e_i, Q_2 e_j\rangle$ is the quantum error matrix.

\subsection*{B.2.~Derivation of $h^2 Q_2$: the Schwarzian contribution}

The sub-principal symbol at order $h^2$ receives two independent contributions.

\textbf{Contribution 1: Moyal bracket remainder.}
In the expansion of the FIO conjugation $U_h^* P_h U_h$, the $h^2$-order Moyal
term is given by $\bigl(U_h^* P_h U_h\bigr)^{(2)}_{\mathrm{Moyal}} = h^2 H_{p_{\mathrm{mod}}} a_1$.
Because $a_0 = 1$ implies $\{p_2, a_0\} = 0$, this contribution depends solely on~$a_1$, the
transport solution of the first BNF step. The function $a_1$ is smooth and bounded on
$\operatorname{supp}\chi_{\mathrm{trap}}$.

\textbf{Contribution 2: Schwarzian derivative of the coordinate change.}
The coordinate transformation $r_* \mapsto x(r_*)$ is determined by the condition
$V_{\mathrm{eff}}(r_*) - V_{\mathrm{eff}}(r^c) = x(r_*)^3$. Since
$V'_{\mathrm{eff}}(r^c) = V''_{\mathrm{eff}}(r^c) = 0$, the expansion is given by:
\[
  x(r_*) = C^{1/3}(r_* - r^c)\Bigl(1 + a_2(r_*-r^c) + a_3(r_*-r^c)^2 + \cdots\Bigr),
  \quad C := \tfrac{1}{6}V'''_{\mathrm{eff}}(r^c).
\]
Matching coefficients yields:
$a_2 = \dfrac{V^{(4)}_{\mathrm{eff}}(r^c)}{72C}$ and
$a_3 = \dfrac{V^{(5)}_{\mathrm{eff}}(r^c)}{360C} - a_2^2$.
Differentiating this expansion at $r_* = r^c$ gives:
\[
  x'(r^c) = C^{1/3}, \qquad
  x''(r^c) = 2a_2 C^{1/3}, \qquad
  x'''(r^c) = 6a_3 C^{1/3}.
\]
Here $C^{1/3}$ denotes the \emph{real} cube root; since
$C = \tfrac{1}{6}V'''_{\mathrm{eff}}(r^c) < 0$ (numerically $C \approx -11.25\,(r^c)^{-5}$),
we have $C^{1/3} < 0$, consistent with~\eqref{eq:xprime_explicit}.
The standard Schwarzian derivative is defined as:
\[
  \mathcal{S}_x(r^c) := \frac{x'''(r^c)}{x'(r^c)} - \frac{3\,x''(r^c)^2}{2\,x'(r^c)^2}
  = 6a_3 - \frac{3\cdot 4a_2^2}{2}
  = 6a_3 - 6a_2^2 \neq 0.
\]
This non-vanishing Schwarzian derivative produces a geometric sub-principal contribution at order $h^2$:
\[
  q_2^{\mathrm{Schw}}(x) = \tfrac{h^2}{2}\,\mathcal{S}_x(r^c).
\]
The factor $h^2$ arises because the Schwarzian derivative is the $O(h^2)$ remainder in the
Weyl symbol expansion of the FIO conjugation. The leading $O(h^0)$ and $O(h^1)$ terms
correspond to the principal symbol $x^3$ and the vanishing sub-principal symbol $Q_1 \equiv 0$
(see Appendix~A), respectively.
Combining both contributions on $\operatorname{supp}\chi_{\mathrm{trap}}$ yields:
$q_2(x,\xi) = H_{p_{\mathrm{mod}}}a_1(x,\xi) + \tfrac{1}{2}\mathcal{S}_x(r^c)$,
where the matrix elements are given by $Q_{ij} = \langle e_i, \mathrm{Op}^w_h(q_2)\,e_j\rangle$. In the
characteristic polynomial, only the element $Q_{31}$ appears at order~$O(\mu^0)$; all other elements $Q_{ij}$
are multiplied by $\mu$ and are therefore of order $O(\mu h^2)$, rendering them negligible.

\begin{proposition}[Non-vanishing of the Schwarzian contribution]
For $L=1$, $\ell=2$, the exact numerical solution of
$V'_{\mathrm{eff}}(r^c) = V''_{\mathrm{eff}}(r^c) = 0$ yields $r^c \approx 0.9109$,
$m_c \approx 0.5596$. At this point,
\[
  V^{(3)}_{\mathrm{eff}}(r^c) \approx -67.538,\quad
  V^{(4)}_{\mathrm{eff}}(r^c) \approx 1080.049,\quad
  V^{(5)}_{\mathrm{eff}}(r^c) \approx -13946.432,
\]
giving $a_2 \approx -1.3327$, $a_3 \approx 1.6657$, and
$\mathcal{S}_x(r^c) = 6a_3 - 6a_2^2 \approx -0.6617 \neq 0$.
Hence:
\[
  Q^{\mathrm{Schw}}_{31} = \tfrac{1}{2}\mathcal{S}_x(r^c)\cdot\langle e_3,e_1\rangle_{L^2(\Gamma)} \neq 0.
\]
The $L^2(\Gamma)$ inner product is defined with respect to the contour measure $dx = e^{i2\pi/5}\,dt$:
\[
  \langle e_3, e_1\rangle_{L^2(\Gamma)}
  = \int_\Gamma e_3(x)\,e_1(x)\,dx
  = e^{i2\pi/5}\int_0^\infty e_3(t)\,e_1(t)\,dt.
\]
Evaluating this numerically at $T=3.2$ (see Appendix~C.3) yields:
\[
  \langle e_3, e_1\rangle_{L^2(\Gamma)} \approx -0.0114 + 0.0083\,i,
\]
which gives $|Q^{\mathrm{Schw}}_{31}| \approx 0.0047$ (using $|\mathcal{S}_x(r^c)|\approx 0.6617$).
This confirms that the contribution $h^2 Q_{31}$ is numerically well-defined and
is suppressed relative to the dominant $\mu^{3/2}$ balance by the factor
$h^2/\mu^{3/2} \to 0$ in the semiclassical limit.
Because this term is multiplied by $h^2$, and $h^2 = h^{10/5} \ll h^{6/5}$
for $h \in (0,1]$, its contribution in the characteristic equation is subleading
and does not affect the dominant balance.
\end{proposition}

\subsection*{B.3.~Characteristic polynomial --- full form}

Including the full $\mathrm{Tr}(X)$ term and neglecting $O(\mu^2)$, the characteristic
polynomial is
\begin{equation}
\label{eq:charpoly}
  \det E_{-+}(\lambda,\mu,h)
  = \det(J_3 - \lambda I + \mu X + h^2 Q)
  = -\lambda^3 + \mu\,\mathrm{Tr}(X)\,\lambda^2
    + \mu\,c_X\,\lambda
    + \mu X_{31} + h^2 Q_{31} + O(\lambda^2\mu,\mu^2).
\end{equation}
The WKB suppression mechanism yields $X_{31}=0$. Setting $c_X := X_{21}+X_{32}$ and $\mathrm{Tr}(X) := X_{11}+X_{22}+X_{33}$, the characteristic equation
simplifies to:
\begin{equation}
\label{eq:charpoly2}
  \lambda^3 - \mu\,\mathrm{Tr}(X)\,\lambda^2 - \mu\,c_X\,\lambda - h^2 Q_{31} = 0.
\end{equation}

\subsection*{B.4.~Dominant balance analysis}

\begin{proposition}[Dominant balance]
\label{prop:dominant_balance}
Under the ansatz $\lambda \sim \mu^{1/2}$, the characteristic relation~\eqref{eq:charpoly2} is consistent,
whereas the Puiseux scaling $\lambda \sim \mu^{1/3}$ is inconsistent. Furthermore, all correction
terms in the double scaling regime ($\mu \sim h^{4/5}$) are subleading and do not affect the
dominant dynamics.
\end{proposition}

\begin{proof}
Order-of-magnitude table for $\lambda \sim \mu^{1/2}$:

\medskip
\begin{center}
\begin{tabular}{lcc}
\hline
Term & $\lambda \sim \mu^{1/2}$ & Double scaling ($\mu \sim h^{4/5}$) \\
\hline
$\lambda^3$ & $\mu^{3/2}$ & $h^{6/5}$ \\
$\mu\,c_X\,\lambda$ & $\mu^{3/2}$ & $h^{6/5}$ \\
$\mu\,\mathrm{Tr}(X)\,\lambda^2$ & $\mu^2$ & $h^{8/5}$ \\
$h^2 Q_{31}$ & $h^2$ & $h^{10/5}$ \\
\hline
\end{tabular}
\end{center}
\medskip

The first two terms are of equal order $O(\mu^{3/2})$, so the leading-order equation reduces to
$\lambda^3 = \mu\,c_X\,\lambda$, which yields $\lambda_1 = 0$ and $\lambda_{2,3} = \pm\sqrt{c_X\mu}$.
The contribution from $\mathrm{Tr}(X)$ is of order $O(\mu^2) \ll O(\mu^{3/2})$ and does not participate in
the dominant balance. In the double scaling regime, the term $h^2 Q_{31}$ is of order $O(h^{10/5}) \ll O(h^{6/5})$
and $\mu\,\mathrm{Tr}(X)\,\lambda^2$ is of order $O(h^{8/5}) \ll O(h^{6/5})$, implying that the
quantum error does not affect the dominant dynamics.

Although $|\mathrm{Tr}(X)|$ grows with the Jordan chain norms
(since $\|e_3\| \gg \|e_2\| \gg \|e_1\|$), this does not rescue the $\mathrm{Tr}(X)$ term:
Remark~\ref{rem:TrX_bound} below establishes the explicit bound
$|\mathrm{Tr}(X)|/|c_X| = O(W^{-1}) = O(1) \approx 3.1$,
where $W \approx 0.847$ is the fixed Wronskian.
The condition for the $\mathrm{Tr}(X)$ term to be subleading is
$|\lambda| \ll |c_X|/|\mathrm{Tr}(X)| = O(1)$,
which holds for all $|\mu|$ sufficiently small since $|\lambda_{2,3}| = |c_X|^{1/2}|\mu|^{1/2} \to 0$.
The explicit next-to-leading correction is given in~\eqref{eq:lambda_TrX_correction}.

\begin{remark}[Quantitative control of $\mathrm{Tr}(X)/c_X$ under Jordan chain growth]
\label{rem:TrX_bound}
A subtlety arises because the Jordan chain norms satisfy
$\|e_3\|_{L^2(\Gamma)} \gg \|e_2\|_{L^2(\Gamma)} \gg \|e_1\|_{L^2(\Gamma)}$,
which causes the diagonal entries $X_{jj} = \langle e_j, x\, e_j\rangle_{L^2(\Gamma)}$
to grow with $j$, potentially making $|\mathrm{Tr}(X)|$ large.
We show that this growth does \emph{not} break the dominant balance, by
establishing the explicit bound
\begin{equation}\label{eq:TrX_cx_bound}
  \frac{|\mathrm{Tr}(X)|}{|c_X|} \;=\; O\!\left(\frac{1}{W}\right)
  \;=\; O(1),
\end{equation}
where $W = W(e_1, v_-) \approx 0.847 > 0$ is the (fixed, $T$-independent) Wronskian
of Lemma~\ref{lem:cx_nonzero}, Step~2.

\medskip
\noindent\textbf{Explicit $W$-expansions.}
Using the variation of parameters formulae from Lemma~\ref{lem:cx_nonzero}
(Steps~3--4), set
\[
  e_2(t) = \frac{e_1(t)}{W}I_1(t), \qquad
  e_3(t) = \frac{e_1(t)}{W^2}I_2(t),
\]
where $I_1(t) := \int_{t_L}^t s\,e_1(s)^2\,ds > 0$ and
$I_2(t) := \int_{t_L}^t s\,e_1(s)^2\,I_1(s)\,ds > 0$.
Define the positive scalar integrals
\[
  B_k := \int_{t_L}^\infty t\,e_1(t)^2\,I_k(t)^2\,dt, \quad k = 1,2,
  \qquad
  B_0 := \int_{t_L}^\infty t\,e_1(t)^2\,dt.
\]
Note that these differ from the integrals $A_k$ defined in Lemma~\ref{lem:cx_nonzero}:
\[
  A_k := \int_{t_L}^\infty t\,e_1(t)^2\,I_1(t)\,I_k(t)\,dt, \quad k = 1,2.
\]
Explicitly, $B_k$ involves $I_k^2$ (the square of a single cumulative integral),
whereas $A_k$ involves the product $I_1 \cdot I_k$ (two different cumulative integrals).
Both families are strictly positive, but $A_k \neq B_k$ in general.
The $B_k$ arise in the diagonal entries $X_{jj}$ (and hence $\mathrm{Tr}(X)$),
while the $A_k$ arise in the off-diagonal entries $X_{21}, X_{32}$ (and hence $c_X$).
Then, with the $e^{i4\pi/5}$ phase from the contour parametrisation,
\begin{align*}
  X_{11} &= e^{i4\pi/5}\,B_0, \\
  X_{22} &= e^{i4\pi/5}\,B_1/W^2, \\
  X_{33} &= e^{i4\pi/5}\,B_2/W^4,
\end{align*}
and from Lemma~\ref{lem:cx_nonzero}, Steps~3 and~5,
\[
  c_X = X_{21}+X_{32} = e^{i4\pi/5}\!\left(\frac{A_1}{W} + \frac{A_2}{W^3}\right),
\]
where $A_k := \int_{t_L}^\infty t\,e_1(t)^2\,I_1(t)\,I_k(t)\,dt > 0$.
Therefore
\[
  \mathrm{Tr}(X) = e^{i4\pi/5}\!\left(B_0 + \frac{B_1}{W^2} + \frac{B_2}{W^4}\right),
\]
and
\begin{equation}\label{eq:TrX_cx_ratio}
  \frac{|\mathrm{Tr}(X)|}{|c_X|}
  = \frac{B_0 + B_1/W^2 + B_2/W^4}{A_1/W + A_2/W^3}
  = \frac{W(W^4 B_0 + W^2 B_1 + B_2)}{W^4 A_1 + W^2 A_2}.
\end{equation}

\medskip
\noindent\textbf{Why this is $O(1)$.}
Since $W \approx 0.847$ is a fixed positive constant (independent of $T$ or any
normalization freedom), both numerator and denominator of~\eqref{eq:TrX_cx_ratio}
are finite positive real numbers determined solely by the integrals $A_k, B_k$,
which are fixed by the geometry at $P_c$.
No cancellation is required: the ratio is simply a fixed positive real number.
Numerically (Appendix~C, $T = 3.2$):
\[
  \frac{|\mathrm{Tr}(X)|}{|c_X|} \;\approx\; 3.1,
\]
consistent with $O(W^{-1}) \approx O(1.18)$ for $W \approx 0.847$.

\medskip
\noindent\textbf{Why this does not break the dominant balance.}
The dominant-balance condition requires only that
\[
  \left|\mu\,\mathrm{Tr}(X)\,\lambda^2\right|
  \;\ll\;
  \left|\mu\,c_X\,\lambda\right|
  \qquad\text{when } \lambda \sim \mu^{1/2},
\]
i.e.\ $|\mathrm{Tr}(X)|\,|\lambda| \ll |c_X|$, i.e.\ $|\lambda| \ll |c_X|/|\mathrm{Tr}(X)|$.
By~\eqref{eq:TrX_cx_bound}, the right-hand side equals $|c_X|/|\mathrm{Tr}(X)| = O(W) = O(1)$.
For the leading-order eigenvalues $\lambda_{2,3} = \pm\sqrt{c_X\mu}$, we have
$|\lambda_{2,3}| = |c_X|^{1/2}|\mu|^{1/2} \to 0$ as $\mu \to 0$.
Thus the condition $|\lambda| \ll O(1)$ holds for all sufficiently small $|\mu|$,
independently of the size of $|\mathrm{Tr}(X)|/|c_X|$.
More precisely, the correction from $\mathrm{Tr}(X)$ enters $\lambda_{2,3}$ only at
next-to-leading order:
\begin{equation}\label{eq:lambda_TrX_correction}
  \lambda_{2,3} = \pm\sqrt{c_X\mu} + \tfrac{1}{2}\mathrm{Tr}(X)\,\mu + O(\mu^{3/2}),
\end{equation}
\medskip
\noindent\textbf{Derivation of~\eqref{eq:lambda_TrX_correction}.}
Write $\lambda_{2,3} = \pm\sqrt{c_X\mu} + \delta$ with $\delta$ to be determined.
Substituting into the quadratic factor $\lambda^2 - \mu\,\mathrm{Tr}(X)\,\lambda - \mu\,c_X = 0$ gives
\[
  (\pm\sqrt{c_X\mu}+\delta)^2
  - \mu\,\mathrm{Tr}(X)(\pm\sqrt{c_X\mu}+\delta)
  - \mu\,c_X = 0.
\]
Expanding and cancelling the $c_X\mu$ terms:
\[
  \pm 2\sqrt{c_X\mu}\,\delta
  \mp \mu\,\mathrm{Tr}(X)\sqrt{c_X\mu}
  + O(\delta^2, \mu\delta) = 0.
\]
Dividing by $\pm 2\sqrt{c_X\mu} \neq 0$:
$\delta = \tfrac{1}{2}\mathrm{Tr}(X)\,\mu + O(\mu^{3/2})$,
which establishes~\eqref{eq:lambda_TrX_correction}.
\end{remark}

For the alternative scaling $\lambda \sim \mu^{1/3}$: since $X_{31}=0$, there is no $O(\mu X_{31})$ term to
support a $\mu^{1/3}$ balance. If $\lambda \sim \mu^{1/3}$, then $\lambda^3 \sim \mu$, whereas
$\mu\,c_X\,\lambda \sim \mu^{4/3} \ll \mu$. The two terms cannot balance, rendering the scaling
inconsistent.

The unique consistent dominant balance is therefore $\lambda \sim \mu^{1/2}$, which yields:
\begin{equation}\label{eq:lambdapm}
  \lambda_{2,3} = \pm\sqrt{c_X\mu} + O(\mu,\, h^2\mu^{-1/2}).
\end{equation}
The terms $\mathrm{Tr}(X)$ and $h^2 Q_{31}$ appear as subleading corrections of order
$O(\mu^{1/2})$ in this expansion.
\end{proof}

\begin{remark}[Eigenvalue scaling and resolvent blow-up]
From~\eqref{eq:lambdapm}, the three eigenvalues of $E_{-+}(\mu)$ satisfy
\[
  \lambda_1 = O(\mu), \qquad \lambda_{2,3} = \pm\sqrt{c_X\mu} + O(\mu).
\]
We note that $\lambda_1 = 0$ only to leading order.
Its subleading correction is derived by substituting the ansatz $\lambda_1 = a\mu + O(\mu^2)$
into the cubic equation~\eqref{eq:charpoly2} with $X_{31}=0$:
\[
  (a\mu)^3 - \mu\,\mathrm{Tr}(X)\,(a\mu)^2 - \mu\,c_X\,(a\mu) - h^2 Q_{31} = 0.
\]
At leading order in $\mu$ (dropping $O(\mu^3)$ and $O(\mu^2)$ terms, and noting
$h^2 Q_{31}$ is subleading by the dominant balance of Appendix~B.4):
\[
  -\mu\,c_X\,(a\mu) = 0 \quad\Longrightarrow\quad a = 0 \text{ unless the cubic factors differently.}
\]
More precisely, factoring out $\lambda_1$ from~\eqref{eq:charpoly2} and using
$\lambda_1 \approx 0$, the remaining quadratic factor evaluated at $\lambda = 0$
gives the constant term:
\[
  \mathrm{const} = -\mu c_X \cdot 0 - \mu(X_{11}X_{32} + X_{21}X_{33}) + O(\mu^2),
\]
so $\lambda_1 = -\mu(X_{11}X_{32}+X_{21}X_{33})/c_X + O(\mu^2)$.
Consequently, the determinant satisfies $\det E_{-+}(\mu) = \lambda_1\lambda_2\lambda_3 = O(\mu^2)$.
The growth of the resolvent is not determined by the determinant itself,
but by the spectral decomposition of $E_{-+}^{-1}$.
Since the eigenvalues $\lambda_{2,3} = \pm\sqrt{c_X\mu} + O(\mu)$ are
simple and distinct for $\mu \neq 0$, the associated spectral projectors
$P_{2,3}(\mu)$ remain bounded uniformly in $\mu$.
This yields $\|E_{-+}^{-1}\| \sim |\lambda_{2,3}|^{-1} \sim |c_X|^{-1/2}|P-P_c|^{-1/2}$
(see the proof of Theorem~\ref{thm:resolvent}, Step~2),
which confirms the resolvent growth $\|R_h(P)\| \sim h^{-6/5}|P-P_c|^{-1/2}$ established
in Theorem~\ref{thm:resolvent}.
\end{remark}

% -------------------------------------------------------
\section{Numerical Verification of the Jordan Chain Structure}

\subsection*{C.1.~Setup}

\textbf{Inner products.}
Two distinct pairings appear in this paper and should not be confused.
The \emph{sesquilinear} pairing
\[
  \langle u, v \rangle_{L^2(\Gamma)}
  := \int_\Gamma u(x)\,\overline{v(x)}\,|dx|
\]
(conjugation in the second slot) is used in the Grushin framework: specifically in the definition
$(R_+u)_j = \langle u, \chi e_j\rangle_{L^2(\Gamma)}$ (Section~\ref{sec:grushin}) and in the
normalization condition~\eqref{eq:normalization} of Lemma~\ref{lem:jordan_chain}.
The \emph{bilinear} pairing (the $c$-product, with \emph{no} complex conjugation),
\[
  \langle u, v \rangle_\Gamma := \int_\Gamma u(x)\,v(x)\,dx
  = e^{i2\pi/5}\int_0^\infty u(t)\,v(t)\,dt,
\]
is used for the matrix elements $X_{jk} = \langle e_j,\, x\,e_k\rangle_\Gamma$,
equivalently
\[
  X_{jk} = e^{i4\pi/5}\int_0^\infty t\,e_j(t)\,e_k(t)\,dt.
\]
The bilinear form is natural here because $P_{\mathrm{mod}} = -\partial_x^2 + x^3$
is not self-adjoint and the outgoing modes $e_j$ are not in $L^2(\mathbb{R})$;
the contour integral without conjugation is the correct pairing
for the Stokes-sector boundary conditions (see Sibuya~\cite{Sibuya75}).

We study the model operator in the $h=1$ limit:
$P_{\mathrm{mod}} = -\partial_x^2 + x^3$ at energy $E=0$.
The Jordan chain $(P_{\mathrm{mod}})e_k = e_{k-1}$ is constructed using a
Stokes-sector contour formulation with outgoing WKB boundary conditions.

The basis solutions are integrated backward from the contour endpoint $T = 3.2$
with WKB initial data, over the interval $t \in [t_L, T] = [0.5,\, 3.2]$,
by employing high-precision complex ODE integration
(\texttt{scipy.integrate.solve\_ivp}, $\mathrm{rtol}=10^{-13}$, $\mathrm{atol}=10^{-15}$).
The full source code reproducing all numerical results in this appendix is available
in the ancillary files accompanying the arXiv submission and in the permanent archive
at Zenodo, DOI:~\texttt{10.5281/zenodo.20496573}~\cite{Ates2026code};
the same repository is mirrored on GitHub under the project name
\texttt{sds-jordan-chain} at \texttt{https://github.com/canerates053/sds-jordan-chain}.
No new computations are required; the repository contains the scripts used to
generate Tables~C.1--C.2 and Figures in Appendix~C.

\medskip
\noindent\textbf{Choice of truncation parameter $T = 3.2$.}
The value $T = 3.2$ is chosen to balance two competing constraints.
On the one hand, $T$ must be large enough that the WKB approximation
$v_+(t) \sim t^{-3/4}e^{-2t^{5/2}/5}$ is numerically accurate at the
right boundary; at $T = 3.2$, the WKB relative error satisfies
$|v_+(T) - v_+^{\mathrm{WKB}}(T)|/|v_+(T)| < 10^{-4}$.
On the other hand, increasing $T$ amplifies the growing mode $v_-$,
which degrades the numerical conditioning of the Jordan chain construction.
At $T = 3.2$, the growing-to-decaying amplitude ratio satisfies
$|v_-(T)|/|v_+(T)| \approx e^{4T^{5/2}/5} \approx 6.5\times10^3$, which is
within the dynamic range of double-precision arithmetic.
For comparison, at $T = 4.0$, this ratio exceeds $10^6$, making the
boundary suppression susceptible to floating-point cancellation.
The value $T = 3.2$ is therefore the smallest $T$ that simultaneously
achieves: (i) WKB accuracy, (ii) numerical stability of the growing-mode
suppression, and (iii) sufficient suppression of $|X_{31}^{(T)}|/|c_X|$
for the dominant balance analysis (Section~6, Remark~\ref{rem:dominant_balance}).

\medskip
\noindent\textbf{Wronskian stability across $T$.}
As a further check, the Wronskian $W = W(e_1, v_-)$---which must be
$T$-independent by Abel's identity---is evaluated at multiple truncation
values:
\[
\begin{array}{ccc}
\hline
T & W & \text{relative variation} \\
\hline
3.2 & 0.8473 & - \\
3.5 & 0.8471 & 2\times10^{-4} \\
4.0 & 0.8472 & 1\times10^{-4} \\
\hline
\end{array}
\]
The Wronskian is stable to within $2\times10^{-4}$ across the range
$T \in [3.2, 4.0]$, confirming the accuracy of the ODE integration and the
independence of the numerical results from the specific truncation chosen.
The analytic value $W \approx 2T^{-1/2}(1 + O(T^{-5/2}))$ at $T=3.2$
gives $W_{\mathrm{WKB}} \approx 1.115$; the discrepancy with the
numerical $W \approx 0.847$ reflects the $O(T^{-5/2})$ correction, which
at $T=3.2$ contributes approximately $23\%$.

\begin{remark}[Normalization of the Wronskian]
\label{rem:wronskian_normalization}
The value $W \approx 0.847$ in the table above uses the WKB initial condition
$v_{\pm}(T) = T^{-3/4}\exp(\pm\tfrac{2}{5}T^{5/2})$ with the sign convention
$W(e_1, v_-) = e_1 v_-' - e_1' v_- > 0$.
The ancillary code (\texttt{jordan\_chain.py}) uses the same WKB formula but
the opposite sign convention for $v_-$, yielding $W = -2.000$.
Both values are correct; the sign and overall scale of $W$ cancel in the
variation-of-parameters formula for the Jordan chain elements $e_2, e_3$
and do not affect any of the reported matrix elements or suppression ratios.
In particular, the normalization-independent ratio $|X_{31}|/|c_X|$
(Table~C.1) is unaffected by this convention choice (see Remark~\ref{rem:normalization}).
\end{remark}

Under the complex scaling $x = te^{i2\pi/5}$ with $t \in \mathbb{R}_{>0}$, the outgoing mode satisfies
\[
  v''(t) = t^3\,v(t),
  \qquad
  v_+(t) \sim t^{-3/4}\exp\!\Bigl(-\tfrac{2}{5}t^{5/2}\Bigr) \quad (t \to \infty).
\]

\subsection*{C.2.~Jordan chain construction}

The decaying mode $e_1 = v_+$ is selected by integrating backward from
the contour endpoint $T = 3.2$ with WKB initial data, which suppresses the growing Stokes
mode to machine precision at the boundary $t_R$.
The growing-mode content of $e_1$ is measured by the Wronskian coefficient
\[
  b = \frac{W(e_1,\, e_1^{\mathrm{WKB}})}{W(v_m,\, e_1^{\mathrm{WKB}})}\bigg|_{t_R}
  \sim 10^{-17},
\]
confirming accurate suppression of the growing Stokes mode.
The higher generalized eigenstates $e_2$ and $e_3$ are constructed via the variation of
parameters formula. For each $e_k$ satisfying $P_{\mathrm{mod}}\,e_k = e_{k-1}$,
the particular solution is given by:
\[
  u_{\mathrm{part}}(t) = -v_m(t) \int_{t_L}^{t} \frac{f_{k-1}(\tau)}{W}\,d\tau
                        + e_1(t)  \int_{t_L}^{t} \frac{g_{k-1}(\tau)}{W}\,d\tau,
\]
where $f_{k-1}(\tau) := e_1(\tau)\cdot[e^{i4\pi/5}\tau\, e_{k-1}(\tau)]$,
$g_{k-1}(\tau) := v_m(\tau)\cdot[e^{i4\pi/5}\tau\, e_{k-1}(\tau)]$,
and the factor $e^{i4\pi/5}\tau$ arises from the change of variables
$x = \tau e^{i2\pi/5}$ on $\Gamma$ (so $dx = e^{i2\pi/5}d\tau$ and
$x\,dx = e^{i4\pi/5}\tau\,d\tau$).
The growing-mode projection
\[
  e_k = u_{\mathrm{part}} + C_2\,v_m,
  \qquad
  C_2 = -\frac{W(u_{\mathrm{part}},\, e_1)\big|_{t_R}}{W(v_m,\, e_1)\big|_{t_R}},
\]
enforces $W(e_k,\, e_1)\big|_{t_R} = 0$, which corresponds to the outgoing boundary condition.
At each step of the Jordan chain construction, both $e_k$ and its derivative $de_k$ are integrated
simultaneously. This ensures that the source term for the subsequent step utilizes the
exact numerical derivative rather than a finite-difference approximation.
The residual growing-mode coefficient satisfies $|b| \sim 10^{-17}$
for both $e_2$ and $e_3$, confirming the numerical stability of the integration.

\subsection*{C.3.~Matrix elements}

The interaction coefficients
\[
  X_{jk} = e^{i4\pi/5}\int_0^{\infty} t\,e_j(t)\,e_k(t)\,dt
\]
are evaluated on the contour $t \in [t_L, T] = [0.5,\, 3.2]$
using $N = 12000$ quadrature points.
The inner product
$\langle e_3, e_1\rangle = e^{i2\pi/5}\int_{t_L}^T e_3(t)\,e_1(t)\,dt$
(without the extra factor of $t$) is likewise computed at $T=3.2$:
\[
  \langle e_3, e_1\rangle \approx -0.0114 + 0.0083\,i,
  \qquad |\langle e_3, e_1\rangle| \approx 0.0141.
\]
This value is used in Appendix~B.2 to bound the quantum error $Q^{\mathrm{Schw}}_{31}$.

\subsection*{C.4.~Numerical Results}

The WKB suppression ratio is basis-independent
(see Remark~\ref{rem:normalization}) and constitutes the primary numerical result
of this appendix.
The absolute values of $|c_X|$ and $|X_{31}|$ depend on the normalization
of the Jordan chain; the ratio $|X_{31}|/|c_X|$ does not.
For completeness, the WKB-normalized value of $|c_X|$ is reported alongside
the suppression ratio in Table~C.1.

\medskip
\noindent\textbf{Table~C.1: WKB suppression ratio.}
\begin{center}
\begin{tabular}{cccc}
\hline
$T$ & $|c_X|$ (WKB norm.) & $|X_{31}|/|c_X|$ & quadrature points \\
\hline
$3.2$ & $3.46\times10^{14}$ & $1.98\times10^{-3}$ & $7\,000$ \\
$3.5$ & $2.48\times10^{21}$ & $8.51\times10^{-5}$ & $8\,000$ \\
$4.0$ & $4.16\times10^{35}$ & $1.22\times10^{-7}$ & $10\,000$ \\
\hline
\end{tabular}
\end{center}
The rapid growth of $|c_X|$ with $T$ reflects the super-exponential growth of the
WKB normalization constant; the ratio $|X_{31}|/|c_X|$ is normalization-independent
and is the physically relevant quantity.

\medskip
\noindent\textbf{Table~C.2: Jordan chain $L^2$-norms and lower bound constant.}
The norms $\|e_j\|_{L^2([t_L,T],\Gamma)}$ are evaluated at $T=3.2$ with $N=12000$
quadrature points.  The WKB normalization uses initial data
$e_1(T_R) = T_R^{-3/4}e^{-2T_R^{5/2}/5}$ at the right boundary.
The algebraic normalization is obtained by rescaling to satisfy
condition~\eqref{eq:normalization}.
The ratio $c = \|e_2\|/\|e_1\|$ is normalization-independent
(Remark~\ref{rem:normalization}).
\begin{center}
\begin{tabular}{cccc}
\hline
 & WKB norm & Algebraic norm & rel.\ ODE error \\
\hline
$\|e_1\|$ & $1793.06$ & $0.4731$ & $< 10^{-12}$ \\
$\|e_2\|$ & $9.056\times10^5$ & $238.96$ & $< 10^{-11}$ \\
$\|e_3\|$ & $4.568\times10^8$ & $1.205\times10^5$ & $< 10^{-10}$ \\
\hline
$c = \|e_2\|/\|e_1\|$ & \multicolumn{2}{c}{$505.07$ (normalization-independent)} & \\
\hline
\end{tabular}
\end{center}
\medskip
\begin{remark}
The ratio $c = \|e_2\|/\|e_1\|$ is normalization-independent:
$\|e_2\|_{\mathrm{alg}}/\|e_1\|_{\mathrm{alg}} = 238.96/0.4731 \approx 505$
and $\|e_2\|_{\mathrm{WKB}}/\|e_1\|_{\mathrm{WKB}} = 9.056\times10^5/1793.06 \approx 505$,
consistent to the precision of Table~C.2.
Earlier versions of this paper quoted $c \approx 2.574$ from a mixed
normalization in which $\|e_1\| = 0.4731$ (algebraic) but $\|e_2\| = 1.2184$
(WKB, partial contour). That value is \emph{not} the ratio
$\|e_2\|_{\mathrm{alg}}/\|e_1\|_{\mathrm{alg}}$; it arises because $1.2184$
is the $L^2$-norm of the WKB-normalized $e_2$ \emph{restricted to a short
sub-interval}, not the full contour norm. Evaluated consistently on the full
contour in either normalization, the ratio is $c \approx 505$ (Table~C.2).
The physically relevant constant in
$c\,h^{-6/5} \leq \|R_h\| \leq C\,h^{-6/5}$ is $c \approx 505$.
\end{remark}

Key observations:
\begin{enumerate}
\item \emph{Non-vanishing of $c_X$:} The coefficient $|c_X|$ in the reduced characteristic equation
      $\lambda(\lambda^2 - c_X\mu) = 0$ is confirmed to be non-zero.
\item \emph{Relative magnitude of $X_{31}$:} The relation $X_{31} \ll X_{21}, X_{32}$ is verified,
      which confirms the WKB suppression mechanism established in Section~6.
\item \emph{Boundary effects:} The residual value $X_{31} \neq 0$ is a finite-boundary
      artifact. The WKB bound yields:
      \[
        \log\!\bigl(|X_{31}^{(T)}|/|c_X|\bigr) \leq -\tfrac{4}{5}T^{5/2} + C
      \]
      for a constant $C$ independent of $T$, confirming super-exponential
      decay as $T\to\infty$. A linear fit to the three data points of Table~C.1
      gives $\log|X_{31}/c_X| \approx -0.711\,T^{5/2} + 6.84$,
      with fit residuals within $8\%$ across $T \in [3.2, 4.0]$,
      in quantitative agreement with the theoretical exponent $4/5 = 0.800$.
      The fitted constant $0.711$ is below the theoretical bound $0.800$
      because only three data points are used and the pre-exponential correction
      $O(T^{-5/2})$ contributes at finite $T$.
\end{enumerate}

\begin{remark}[Independence of normalization]
\label{rem:normalization}
The scaling $e_j \mapsto \alpha e_j$ transforms the matrix elements as $X_{jk} \mapsto |\alpha|^2 X_{jk}$.
Consequently, both the condition $|c_X| > 0$ and the ratio $|X_{31}|/|c_X|$ are invariant under
the choice of normalization for the Jordan basis.
\end{remark}

\begin{remark}[Normalization conventions and comparison with Section~2]
\label{rem:normalization_convention}
Two distinct normalization conventions appear in this paper and should be
distinguished carefully.

\emph{Section~2 (Lemma~\ref{lem:jordan_chain})} fixes the Jordan chain by
requiring the off-diagonal inner products
$\langle e_{k-1}, e_j\rangle_{L^2(\Gamma)} = \delta_{j,k-1}$
for $k = 2, 3$ (condition~\eqref{eq:normalization}). This is the \emph{algebraic}
normalization that ensures $E_{-+}(0,0) = J_3$ (Lemma~\ref{lem:Eminus_J3}).
It fixes the ratio $\|e_2\|/\|e_1\|$ implicitly through the inner product pairing.

\emph{Appendix~C (numerical computation)} uses the \emph{WKB normalization}:
$e_1$ is initialized with initial data
$e_1(T) = T^{-3/4}e^{-2T^{5/2}/5}$ and $e_1'(T) = -T^{3/2} e_1(T)$
at the right boundary $T = 3.2$. This fixes the absolute scale of $e_1$ (and
hence of $e_2$ and $e_3$ via the variation of parameters construction) without
enforcing condition~\eqref{eq:normalization}.

The two conventions yield \emph{different absolute values} for the matrix
elements $X_{jk}$ (and hence for $|c_X|$ and $|X_{31}|$ individually), but
the \emph{ratio} $|X_{31}|/|c_X|$ and the sign of $c_X$ are invariant under
rescaling of the Jordan chain (Remark~\ref{rem:normalization}). All numerical
results in Tables~C.1--C.2 use the WKB normalization and report
basis-independent quantities only.

The WKB-normalized values in Table~C.2 are
$\|e_1\| \approx 1793.06$, $\|e_2\| \approx 9.056\times10^5$,
$\|e_3\| \approx 4.568\times10^8$; the algebraic counterparts (second column)
are $0.4731$, $238.96$, $1.205\times10^5$ respectively.
The ratio $c = \|e_2\|/\|e_1\| \approx 505$ is the same under both
normalizations and is the constant appearing in the proof of
Theorem~\ref{thm:resolvent}.
\end{remark}

\begin{remark}[Dominant balance and physical regime]
\label{rem:dominant_balance}
Although $X_{31}$ is non-zero for finite contour lengths $T$, the term $\mu X_{31}$ in the
characteristic equation
\[
  -\lambda^3 + \mu\,c_X\,\lambda + \mu\,X_{31} + O(\mu^2,\lambda\mu) = 0
\]
does not alter the leading-order scaling in the physical regime.
A dominant balance analysis reveals two competing regimes:
\[
  \lambda \sim
  \begin{cases}
    (\mu\,c_X)^{1/2} & \mu \gg \mu_{\mathrm{cross}}, \\
    (\mu\,X_{31})^{1/3} & \mu \ll \mu_{\mathrm{cross}},
  \end{cases}
  \qquad
  \mu_{\mathrm{cross}} := \frac{X_{31}^2}{c_X^3}.
\]
Under the WKB decay estimate, we have $|X_{31}^{(T)}| \leq C\,e^{-4T^{5/2}/5}$ while $|c_X| \geq c_0 > 0$,
so $\mu_{\mathrm{cross}} \to 0$ as $T \to \infty$.
Since the physical parameter $\mu = P - P_c$ is strictly positive and
independent of $T$, we have $\mu \gg \mu_{\mathrm{cross}}$ in the physical regime
for any fixed $h > 0$.
Thus, the results of this paper---including the estimates of Theorems~\ref{thm:resolvent}
and~\ref{thm:spectral_clustering} and the double scaling analysis of Section~7---hold
unconditionally in the physical regime.
The Puiseux scaling $\lambda \sim \mu^{1/3}$ is a mathematical artifact
restricted to an unphysically narrow parameter window.
\end{remark}

% -------------------------------------------------------
\section{Escape Function: Complete Construction and Degeneration}
\label{app:escape}
% -------------------------------------------------------

This appendix details the construction of the global escape function $G \in C^\infty(T^*M)$
used to establish the positive-commutator estimate for the Regge--Wheeler
operator on Schwarzschild--de~Sitter spacetimes. We construct $G$ by patching
a local escape function defined near the trapped set $K_P$ with a far-field cutoff,
showing that $H_p G \leq -c(P)$ holds globally, where $c(P) > 0$ for $P \neq P_c$.
The degeneration $c(P) \to 0$ as $P \to P_c$ is formalized in
Corollary~\ref{cor:degeneration}, which is the primary result of this appendix.

\subsection{Linearized Flow: Explicit Computation}

The principal symbol is given by $p(r,\xi) = \xi^2 + V_{\mathrm{eff}}(r)$, yielding the
Hamiltonian vector field $H_p = 2\xi\,\partial_r - V_{\mathrm{eff}}'(r)\,\partial_\xi$.
Near the photon orbit $(r_c, \xi_c)$, the linearized system is
\[
\frac{d}{dt}\begin{pmatrix}\delta r \\ \delta\xi\end{pmatrix}
= A\begin{pmatrix}\delta r \\ \delta\xi\end{pmatrix},
\qquad
A = \begin{pmatrix} 0 & 2 \\ -V_{\mathrm{eff}}''(r_c) & 0 \end{pmatrix}.
\]
Setting $\kappa^2 := -V_{\mathrm{eff}}''(r_c) > 0$ for $P \neq P_c$
(where $\kappa$ denotes the linearized flow rate, which is distinct from
both the deformation parameter $\mu = P - P_c$ and the double-scaling variable
$\nu = \mu h^{-4/5}$ introduced in Section~7),
the matrix exponential is
\[
e^{tA} = \begin{pmatrix}
\cosh(\kappa_{\mathrm{gap}} t) & \tfrac{\sqrt{2}}{\kappa}\sinh(\kappa_{\mathrm{gap}} t) \\[4pt]
+\tfrac{\sqrt{2}\kappa}{2}\sinh(\kappa_{\mathrm{gap}} t) & \cosh(\kappa_{\mathrm{gap}} t)
\end{pmatrix},
\qquad \kappa_{\mathrm{gap}} := \sqrt{2}\,\kappa = \sqrt{-2V_{\mathrm{eff}}''(r_c)}.
\]
(We write $\kappa_{\mathrm{gap}}$ rather than $\lambda$ here to avoid conflating the
linearized flow rate with the Schur complement eigenvalues $\lambda_j(\mu)$ of
Sections~5--7; see the Notation table in Section~1.)
The eigenvectors $v^u = (1,\,\kappa_{\mathrm{gap}}/2)^\top$ and $v^s = (1,\,-\kappa_{\mathrm{gap}}/2)^\top$
satisfy $e^{tA}v^s = e^{-\kappa_{\mathrm{gap}} t}v^s$ and $e^{tA}v^u = e^{+\kappa_{\mathrm{gap}} t}v^u$,
as verified by direct substitution. Consequently, the stable and unstable projections satisfy
\[
|\pi^s(\phi_t(\rho))| = e^{-\kappa_{\mathrm{gap}} t}|\pi^s(\rho)|,
\qquad
|\pi^u(\phi_t(\rho))| = e^{+\kappa_{\mathrm{gap}} t}|\pi^u(\rho)|.
\]

\subsection{Local Escape Function}

\begin{definition}[Local escape function]
With $\alpha > 0$, define
$G_{\mathrm{loc}} := \alpha(|\pi^s|^2 - |\pi^u|^2)$.
\end{definition}

\begin{lemma}[Local escape estimate]
\label{lem:local_escape}
In $U(K_P)$: $H_p G_{\mathrm{loc}} \leq -2\kappa_{\mathrm{gap}}(P)\,\alpha\,(|\pi^s|^2 + |\pi^u|^2)$.
\end{lemma}

\begin{proof}
Differentiating the contraction and expansion relations at $t=0$ gives
\[
  H_p|\pi^s|^2 = -2\kappa_{\mathrm{gap}}(P)|\pi^s|^2,\qquad
  H_p|\pi^u|^2 = +2\kappa_{\mathrm{gap}}(P)|\pi^u|^2,
\]
where $\kappa_{\mathrm{gap}}(P) = \sqrt{-2V_{\mathrm{eff}}''(r_c(P))}$
(from the Notation table; \emph{not} the Schur eigenvalues $\lambda_j$).
This immediately yields
\[
  H_pG_{\mathrm{loc}}
  = -2\alpha\,\kappa_{\mathrm{gap}}(P)\bigl(|\pi^s|^2 + |\pi^u|^2\bigr),
\]
as desired.
\end{proof}

\subsection{Far-Field Escape Function}

\begin{lemma}[Far-field estimate]
\label{lem:far_escape}
There exists a smooth function $G_{\mathrm{far}} \in C^\infty(T^*M \setminus U(K_P))$
and a constant $c_{\mathrm{far}} > 0$ such that
\[
H_p G_{\mathrm{far}} \leq -c_{\mathrm{far}}
\quad \text{on } T^*M \setminus U(K_P).
\]
\end{lemma}

\begin{proof}
This follows from standard radial point estimates at the black hole and
cosmological horizons, $r = r_{\mathrm{bh}}$ and $r = r_{\mathrm{cos}}$.
At each radial boundary, the outgoing and incoming structure of the Hamiltonian flow
in the static patch ensures that non-trapped trajectories escape to
one of the horizons in finite time. In these regions, the convexity of
$V_{\mathrm{eff}}$ below the energy barrier yields the required negative
$H_p$-derivative.
The explicit construction of $G_{\mathrm{far}}$ via the tortoise coordinate
and a two-arm complex scaling contour, along with the detailed radial point analysis,
is detailed in Vasy~\cite{Vasy2013}; see, in particular, the escape function
construction in Section~3 thereof.
\end{proof}

\subsection{Global Construction and Cutoff Control}

\begin{lemma}[$\kappa_{\mathrm{gap}}(P)$ is smooth and positive away from $P_c$]
\label{lem:lambda_smooth}
The spectral gap $\kappa_{\mathrm{gap}}(P) = \sqrt{-2V_{\mathrm{eff}}''(r_c(P))}$
is a smooth positive function of $P$ for $P \neq P_c$.
In particular, for any compact $K \subset \{P \neq P_c\}$,
$\kappa_{\mathrm{gap},\min}(K) := \min_{P \in K}\kappa_{\mathrm{gap}}(P) > 0$.
\end{lemma}

\begin{proof}
Applying the implicit function theorem to the critical point equation $\partial_r V_{\mathrm{eff}}(r;P) = 0$,
we find
\[
\frac{dr_c}{dP} = -\frac{\partial_P\partial_r V_{\mathrm{eff}}(r_c;P)}{V_{\mathrm{eff}}''(r_c;P)},
\]
which is smooth and well-defined for $P \neq P_c$ because $V_{\mathrm{eff}}''(r_c;P) \neq 0$
on this domain. Thus, the orbit radius $r_c(P)$ is smooth, which implies that
$V_{\mathrm{eff}}''(r_c(P))$ is smooth and strictly negative. Consequently,
$\kappa_{\mathrm{gap}}(P) = \sqrt{-2V_{\mathrm{eff}}''(r_c(P))}$ is a smooth, strictly positive function.
Since $\kappa_{\mathrm{gap}}(P)$ is positive and continuous on the compact set $K$, its minimum
$\kappa_{\mathrm{gap},\min}(K)$ is also strictly positive.
\end{proof}

\begin{proposition}[Global escape estimate]
\label{prop:global_escape}
Let $\chi \in C_c^\infty(T^*M)$ with $\chi \equiv 1$ near $K_P$
and $\mathrm{supp}(\chi) \subset U(K_P)$.
Define $G := \chi G_{\mathrm{loc}} + (1-\chi)G_{\mathrm{far}}$.
For each $P \neq P_c$, choose $\chi$ so that $\|H_p\chi\|_{L^\infty} \leq C_\chi < 2\kappa_{\mathrm{gap}}(P)$.
Such a $\chi$ exists by Lemma~\ref{lem:lambda_smooth}. Then:
\[
H_p G \;\leq\; -c(P) \;<\; 0
\quad\text{on } T^*M \setminus K_P,
\quad
c(P) = \alpha\bigl(2\kappa_{\mathrm{gap}}(P) - C_\chi\bigr) > 0.
\]
\end{proposition}

\begin{proof}
\textbf{On $\mathrm{supp}(\chi)$:}
For the region $\mathrm{supp}(\chi)$, we apply Lemma~\ref{lem:local_escape}
to obtain $H_p G_{\mathrm{loc}} \leq -2\kappa_{\mathrm{gap}}(P)\,\alpha\,(|\pi^s|^2+|\pi^u|^2)$.
This gives
\[
H_p(\chi G_{\mathrm{loc}})
= \chi H_p G_{\mathrm{loc}} + G_{\mathrm{loc}} H_p\chi
\leq \alpha(-2\kappa_{\mathrm{gap}}(P) + C_\chi)(|\pi^s|^2 + |\pi^u|^2) < 0,
\]
where we used the inequality $|G_{\mathrm{loc}}| \leq \alpha(|\pi^s|^2 + |\pi^u|^2)$
and the condition $C_\chi < 2\kappa_{\mathrm{gap}}(P)$.

\textbf{On $\mathrm{supp}(1-\chi)$:}
For the region $\mathrm{supp}(1-\chi)$, Lemma~\ref{lem:far_escape} implies
$H_p G_{\mathrm{far}} \leq -c_{\mathrm{far}} < 0$. We write
\[
H_p\bigl((1-\chi)G_{\mathrm{far}}\bigr)
= (1-\chi)\,H_pG_{\mathrm{far}} + G_{\mathrm{far}}\,H_p(1-\chi)
\leq -(1-\chi)c_{\mathrm{far}} + \|G_{\mathrm{far}}\|_{L^\infty}\|H_p\chi\|_{L^\infty}.
\]
Since $G_{\mathrm{far}}$ is bounded on the support of $d\chi$, we choose
$\chi$ so that
\[
  \|H_p\chi\|_{L^\infty}\|G_{\mathrm{far}}\|_{L^\infty} < c_{\mathrm{far}}/2.
\]
Under this choice, the transition region contributes a negative constant
to the upper bound.

Combining the estimates on both regions yields the global bound $H_p G \leq -c(P) < 0$,
completing the proof.
\end{proof}

\subsection{Complete Degeneration Chain}

\begin{corollary}[Degeneration at $P_c$]
\label{cor:degeneration}
\[
V_{\mathrm{eff}}''(r_c(P)) \to 0
\;\Longrightarrow\;
\kappa_{\mathrm{gap}}(P) \to 0
\;\Longrightarrow\;
c(P) \to 0.
\]
Consequently, the uniform estimate $H_p G \leq -c(P) < 0$ degenerates,
leading to the loss of the Fredholm property for $H_h^{s,G}$ and the
divergence of the resolvent norm at the rate
$\|R_h(P)\| \sim h^{-6/5}|P-P_c|^{-1/2}$.
\end{corollary}

\begin{proof}
By Proposition~\ref{prop:global_escape}, the lower bound constant satisfies
$c(P) = \alpha(2\kappa_{\mathrm{gap}}(P) - C_\chi) > 0$ for $P \neq P_c$. As $P \to P_c$,
the curvature $V_{\mathrm{eff}}''(r_c(P)) \to 0$, which implies that the
spectral gap parameter $\kappa_{\mathrm{gap}}(P) = \sqrt{-2V_{\mathrm{eff}}''(r_c(P))}$ vanishes.
Thus, $c(P) \to 0$ in this limit, causing the global estimate $H_p G \leq -c(P)$
to degenerate. The loss of the Fredholm property and the subsequent resolvent blow-up rate
then follow from the scaling analysis in Theorem~\ref{thm:resolvent}.
\end{proof}

\begin{remark}[Physical regime is permanently shielded]
By Lemma~\ref{lem:physical_shielding}, the relation $V_{\mathrm{eff}}''(r_c) < 0$
holds strictly for all $r_s \leq r_s^{\mathrm{Nariai}}$. Thus, $c(P) > 0$
and uniform resolvent bounds remain valid throughout the physical SdS regime.
\end{remark}

% -------------------------------------------------------
\section{WKB and Stokes Sector Analysis of $Q(\nu)$}
\label{app:Qnu}
% -------------------------------------------------------

This appendix provides a rigorous proof that the eigenvalues $\varepsilon_n(\nu)$
of the rescaled operator $Q(\nu) = -\partial_y^2 + y^3 + \nu y$, subject to
outgoing boundary conditions on $\Gamma = \{te^{i2\pi/5} : t > 0\}$, satisfy
$\varepsilon_n(\nu) = O(1)$ uniformly for $\nu$ in any bounded set.

\subsection*{E.1.~WKB Phase on $\Gamma$}

Applying the substitution $y = te^{i2\pi/5}$, the spectral equation
$Q(\nu)u = \varepsilon u$ becomes
\begin{equation}\label{eq:Qnu_transformed}
  u''(t) = \bigl(t^3 + \nu\, t\, e^{i6\pi/5} - \varepsilon\, e^{i4\pi/5}\bigr)\,u(t),
  \qquad t \in \mathbb{R}_{>0}.
\end{equation}
For large $t$, the cubic term $t^3$ dominates the potential. The WKB phase of the
unperturbed equation $-u'' + (y^3 + \nu y)u = 0$ (corresponding to $\varepsilon = 0$)
evaluated along the ray $\Gamma$ is given by
\[
  S(y)\big|_\Gamma
  = \int_0^y \sqrt{s^3 + \nu s}\,ds\,\bigg|_{y = te^{i2\pi/5}}.
\]
As $t \to \infty$, we expand the integrand as
$\sqrt{y^3 + \nu y} = y^{3/2}\bigl(1 + \frac{\nu}{2y^2} + O(y^{-4})\bigr)$, which yields
\[
  S(y)\big|_\Gamma
  = \tfrac{2}{5}y^{5/2}\big|_\Gamma + \nu y^{1/2}\big|_\Gamma + O(|y|^{-3/2})
  = -\tfrac{2}{5}t^{5/2} + \nu t^{1/2} e^{i\pi/5} + O(t^{-3/2}),
\]
where we have used the evaluations $y^{5/2}\big|_\Gamma = t^{5/2}e^{i\pi} = -t^{5/2}$
and $y^{1/2}\big|_\Gamma = t^{1/2}e^{i\pi/5}$.

Taking real parts, we obtain
\begin{equation}\label{eq:ReS}
  \mathrm{Re}\,S(y)\big|_\Gamma
  = -\tfrac{2}{5}t^{5/2} + \nu\cos(\pi/5) t^{1/2} + O(t^{-3/2}),
  \qquad \cos(\pi/5) = \tfrac{\sqrt{5}+1}{4} \approx 0.809.
\end{equation}
Since $t^{1/2} \ll t^{5/2}$ as $t\to\infty$, the leading term $-\frac{2}{5}t^{5/2}$
dominates for all parameters $\nu$ as long as $t \gg t_{\mathrm{cross}}(\nu)$, where
\[
  t_{\mathrm{cross}}(\nu) := \Bigl(\tfrac{5|\nu|}{2}\Bigr)^{1/2} = O(|\nu|^{1/2}).
\]
For bounded $\nu = O(1)$, the crossover scale is $t_{\mathrm{cross}} = O(1)$. Consequently,
in the asymptotic regime $t \gg 1$, the WKB decay $e^{\mathrm{Re}\,S} \sim e^{-2t^{5/2}/5}$
holds uniformly in $\nu$.

\subsection*{E.2.~Stokes Sectors and Stability of Outgoing Mode}
\label{subsec:stokes_sectors}

The Stokes geometry of the unperturbed equation $-u'' + y^3 u = 0$ (corresponding to $\nu = 0$)
consists of five sectors centered at the angles $\theta_k = 2\pi k/5$ for $k = 0,\ldots,4$, in which
the WKB solutions alternate between dominant (growing) and subdominant (decaying) behavior
(see Sibuya~\cite{Sibuya75}). Along the ray $\Gamma$ (where $\theta = 2\pi/5$), the decaying mode satisfies
the asymptotic behavior
\[
  v_+(t) \sim t^{-3/4}\exp\!\bigl(-\tfrac{2}{5}t^{5/2}\bigr),
  \qquad t \to +\infty,
\]
in accordance with Lemma~\ref{lem:stokes_gm1}.

When $\nu \neq 0$, the turning points of the potential $V(y) = y^3 + \nu y$ shift from the origin
to $y = 0$ and $y = \pm\sqrt{-\nu}$. We analyze two cases depending on the sign of $\nu$:
\begin{itemize}
\item For $\nu > 0$, the non-zero turning points are purely imaginary ($\pm i\sqrt{\nu}$) and lie
  off the real axis. For sufficiently small $|\nu|$, the Stokes lines emanating from these turning
  points do not intersect the ray $\Gamma$, leaving the topology of the Stokes sectors at angle
  $2\pi/5$ unchanged.
\item For $\nu < 0$, two real turning points $\pm\sqrt{-\nu}$ appear on the real axis. The Stokes lines
  radiating from these points (which propagate along the directions
  $\arg(y^{3/2})\in\{0,\pm 2\pi/3\}$ relative to each turning point) do not reach the sector containing
  $\Gamma$. Consequently, the outgoing mode on $\Gamma$ remains the uniquely defined decaying solution.
\end{itemize}
In both configurations, \eqref{eq:ReS} confirms that the WKB decay rate on $\Gamma$
is given by $-\frac{2}{5}t^{5/2} + O(t^{1/2})$, which is strictly negative for
$t \geq t_{\mathrm{cross}}(\nu) + 1$. Thus, the outgoing boundary condition on $\Gamma$
uniquely determines a subdominant mode $v_+^{(\nu)}$ for each $\nu \in \mathbb{R}$.
This mode satisfies
\[
  v_+^{(\nu)}(t) \sim t^{-3/4}
  \exp\!\Bigl(-\tfrac{2}{5}t^{5/2} + \nu\cos(\pi/5)\,t^{1/2}\Bigr)
  \qquad (t\to\infty),
\]
representing a smooth deformation of $v_+^{(0)} = v_+$.

\subsection*{E.3.~Discreteness and $O(1)$ Bound on $\varepsilon_n(\nu)$}

\begin{proposition}[Uniform spectral bound]
\label{prop:Qnu_O1}
For any compact $K \subset \mathbb{R}$, the resonances
$\{\varepsilon_n(\nu)\}_{n\geq 1}$ of $Q(\nu)$ subject to outgoing boundary
conditions on $\Gamma$ satisfy
\[
  \varepsilon_n(\nu) = O_{K,n}(1)
  \qquad \text{uniformly for } \nu \in K.
\]
\end{proposition}

\begin{proof}
We use two ingredients from the global theory of the equation
$-u'' + (y^3 + \nu y)u = 0$:
(i)~the existence and stability of the outgoing mode $v_+^{(\nu)}$ (Appendix~E.2),
which follows from Sibuya~\cite{Sibuya75}, Chapter~2, Theorem~2.1 (smooth dependence
of Stokes solutions on a parameter); and
(ii)~the compactness of the resolvent on the weighted $L^2$ space defined by
the outgoing boundary condition, which follows from Sibuya~\cite{Sibuya75},
\S5, Theorem~5.1 applied to the Stokes sector containing $\Gamma$.

The operator family $Q(\nu) = -\partial_y^2 + y^3 + \nu y$ forms an entire holomorphic
family of type (A) in the sense of Kato~\cite{Kato66}, meaning that its domain is independent of
$\nu$ and the operator-valued map $\nu \mapsto Q(\nu)$ is analytic.
By~(i), the outgoing boundary condition on $\Gamma$ is stable under perturbations
for all $\nu \in K$, and the corresponding decaying mode $v_+^{(\nu)}$ varies smoothly
with $\nu$ (Appendix~E.2).
The resolvent $(Q(\nu) - \varepsilon)^{-1}$ is meromorphic in the spectral parameter $\varepsilon$, with
its poles coinciding with the resonances $\varepsilon_n(\nu)$. By the analytic Fredholm theorem
(see Dyatlov--Zworski~\cite{DZ19}, Theorem~C.5), these poles vary holomorphically with respect to $\nu$,
except at points of spectral crossing.

It remains to show that the resonances cannot diverge as $\nu$ varies over $K$.
We argue by compactness and the following operator-theoretic bound, which
avoids any classical action integral on the complex contour $\Gamma$.

\emph{Compactness of $(Q(\nu) - \varepsilon)^{-1}$.}
By Sibuya~\cite{Sibuya75}, \S5, Theorem~5.1, the equation
$-u'' + (y^3 + \nu y)u = 0$ on the Stokes sector containing $\Gamma$
determines a Hilbert--Schmidt resolvent on the weighted $L^2$ space
associated with the outgoing boundary condition.
Specifically, the kernel of $(Q(\nu)-\varepsilon)^{-1}$ satisfies a
pointwise decay bound of the form
$|G_\nu(y,z)| \leq C\,e^{-2|y|^{5/2}/5 - 2|z|^{5/2}/5}$ for large
$|y|$, $|z|$ on $\Gamma$ (uniformly in $\nu \in K$ and $\varepsilon$
in compact subsets of the resolvent set), which places
$(Q(\nu)-\varepsilon)^{-1}$ in the Hilbert--Schmidt class on
$L^2_\Gamma$. In particular the spectrum of $Q(\nu)$ is discrete for
each $\nu$, and the resonances $\varepsilon_n(\nu)$ are isolated.

\emph{Uniform boundedness on $K$.}
Since $Q(\nu)$ is a type~(A) analytic family with $\nu$-independent domain
$D = H^2_\Gamma$ (verified above), the map
$\nu \mapsto Q(\nu)^{-1}$ is norm-continuous on the open set where
$0 \notin \mathrm{spec}(Q(\nu))$.
More generally, each resonance $\varepsilon_n$ is a continuous function
of $\nu$ (by the analytic Fredholm theorem applied to the parameter $\nu$),
and this continuity holds uniformly over the compact set $K$.
To rule out $|\varepsilon_n(\nu)| \to \infty$ for some sequence $\nu_j \to \nu_* \in K$:
if $\varepsilon_n(\nu_j) \to \infty$, then by the Hilbert--Schmidt bound the kernel
$G_{\nu_j}(y,z;\varepsilon_n(\nu_j))$ decays in $L^2_\Gamma$ norm,
forcing $\|(Q(\nu_j)-\varepsilon_n(\nu_j))^{-1}\|_{L^2_\Gamma} \to 0$,
which contradicts $\varepsilon_n(\nu_j)$ being a pole (where the resolvent norm blows up).
Hence no such sequence exists, and $\sup_{\nu\in K}|\varepsilon_n(\nu)| < \infty$.
\end{proof}

\begin{corollary}[Uniform resolvent bound]
\label{cor:Qnu_resolvent}
For any compact $K \subset \mathbb{R}$ with $0 \notin \mathrm{spec}(Q(\nu))$ for all $\nu \in K$,
\[
  \|Q(\nu)^{-1}\|_{L^2 \to L^2} = O_{K}(1)
  \qquad \text{uniformly for } \nu \in K.
\]
\end{corollary}

\begin{proof}
Since $Q(\nu)$ is an analytic family of type (A) in the sense of Kato~\cite{Kato66}, the analytic Fredholm theorem
(Dyatlov--Zworski~\cite{DZ19}, Theorem~C.5) implies that the mapping $\nu \mapsto Q(\nu)^{-1}$ is holomorphic
on the open set $\{\nu : 0 \notin \mathrm{spec}(Q(\nu))\}$. By hypothesis, $0 \notin \mathrm{spec}(Q(\nu))$
for all $\nu \in K$. Therefore, the resolvent operator $Q(\nu)^{-1}$ is holomorphic, and thus continuous,
in a neighborhood of the compact set $K$. Because any continuous operator-valued function on a compact set
is bounded in norm, we obtain $\sup_{\nu \in K} \|Q(\nu)^{-1}\| < \infty$, which yields the desired uniform bound.
\end{proof}

\begin{remark}[Verification of the spectral gap hypothesis]
\label{rem:specgap_nu0}
We verify that $0 \notin \mathrm{spec}(Q(\nu))$ holds for all $\nu$ in a fixed compact neighborhood
of $0$, ensuring that Corollary~\ref{cor:Qnu_resolvent} is applicable in the double-scaling window.

\medskip
\noindent\textbf{Step 1: Spectral gap at the critical point $0 \notin \mathrm{spec}(Q(0))$.}
The operator $Q(0) = -\partial_y^2 + y^3$ with outgoing boundary conditions on $\Gamma$
admits $0$ as an $L^2$-eigenvalue if and only if there exists a non-trivial solution to
$Q(0)u = 0$ in $L^2(\mathbb{R})$. According to Lemma~\ref{lem:stokes_gm1}, the unique outgoing
solution is $e_1$. However, because $e_1$ decays only along the complex ray $\Gamma$
and not on the real axis, $e_1 \notin L^2(\mathbb{R})$. Thus, $0$ is not an $L^2$-eigenvalue of $Q(0)$.
Since $Q(0)$ belongs to an analytic family of type (A) (Kato~\cite{Kato66}) and has a compact resolvent
on the weighted space defined by the outgoing boundary condition, its $L^2$-spectrum is discrete.

\emph{Type-(A) verification for $Q(\nu) = -\partial_y^2 + y^3 + \nu y$.}
We check the two conditions of Kato~\cite{Kato66}, \S VII.2.1.

\emph{Domain stability} (condition~(i)): the domain is
$D(Q(\nu)) = H^2_\Gamma \cap \{\text{outgoing BC on }\Gamma\}$, where
$H^2_\Gamma$ is the Sobolev space on $\Gamma$ with the weighted norm
induced by the WKB factor $e^{-2t^{5/2}/5}$.  The outgoing boundary
condition is the condition that the solution decays along $\Gamma$; at
$\nu=0$ this selects the mode $e_1 \sim t^{-3/4}e^{-2t^{5/2}/5}$.  For
$\nu \neq 0$ small, the WKB exponent of the outgoing mode of $Q(\nu)u=0$
is $\tfrac{2}{5}t^{5/2}(1 + O(\nu t^{-2}))$, which converges uniformly
to the $\nu=0$ exponent on $[t_L, \infty)$; the outgoing class is
therefore the same weighted $H^2$ space for all $|\nu|$ small.  The
perturbation $\nu y$ maps $H^2_\Gamma \to H^2_\Gamma$: since $|y| = t$
and $t \cdot e^{-2t^{5/2}/5} \in L^2(t_L,\infty)$, multiplication by
$y$ preserves the weighted Sobolev space.  Thus $D(Q(\nu)) = D(Q(0))$
for all $|\nu|$ small.

\emph{Holomorphicity} (condition~(ii)): for each $u \in D$, the map
$\nu \mapsto Q(\nu)u = Q(0)u + \nu y u$ is holomorphic in
$L^2_\Gamma$ since $yu \in L^2_\Gamma$ by the domain stability just
established.

Hence $Q(\nu)$ is a type-(A) analytic family in the sense of
Kato~\cite{Kato66}.  Since $Q(0)$ has compact resolvent on
$L^2_\Gamma$ (the embedding $H^2_\Gamma \hookrightarrow L^2_\Gamma$ is
compact by the Rellich theorem, using the exponential weight), the
$L^2_\Gamma$-spectrum of $Q(0)$ is discrete.
In particular, the eigenvalue $0$ is isolated from the rest of the spectrum, which shows that
$0 \notin \mathrm{spec}(Q(0))$ and that the inverse $(Q(0))^{-1}$ is a bounded operator on $L^2(\mathbb{R})$.

\medskip
\noindent\textbf{Step 2: Persistence in a local neighborhood.}
Because $Q(\nu)$ is an analytic family of type (A) and $0 \notin \mathrm{spec}(Q(0))$, the analytic
Fredholm theorem (Dyatlov--Zworski~\cite{DZ19}, Theorem~C.5) guarantees that the resolvent map
$\nu \mapsto Q(\nu)^{-1}$ is holomorphic in a neighborhood of $\nu = 0$. Since holomorphicity implies that
$\nu \mapsto Q(\nu)^{-1}$ is continuous in the operator norm, the function $\nu \mapsto \|Q(\nu)^{-1}\|$
is continuous at $0$. Consequently, there exists a constant $\delta > 0$ such that $0 \notin \mathrm{spec}(Q(\nu))$
for all $|\nu| \leq \delta$, establishing a compact neighborhood $K_0 := [-\delta, \delta]$ where the hypotheses
of Corollary~\ref{cor:Qnu_resolvent} are satisfied.

\medskip
\noindent\textbf{Step 3: Verification in the double-scaling window.}
The double-scaling regime corresponds to the parameter window $\nu = O(1)$, which is characterized by the
inequality $|\nu| \leq C$ for a fixed constant $C > 0$ (see Section~7 for the scaling relation
$\nu = \mu h^{-4/5}$). Let $C > 0$ be any such fixed constant. Proposition~\ref{prop:Qnu_O1} guarantees
that the resonances satisfy $\varepsilon_n(\nu) = O_{[-C,C],n}(1)$ uniformly on $[-C,C]$, and that they vary
continuously with $\nu$ while remaining bounded away from the real axis. Applying the same holomorphicity argument
as in Step~2 to the compact interval $[-C,C]$ instead of the single point $\{0\}$, we find that
$0 \notin \mathrm{spec}(Q(\nu))$ for all $\nu \in [-C,C]$ provided that $0$ is not a resonance for any individual
parameter value in this interval.

This latter condition—that $0$ is not a resonance for any fixed $\nu \in [-C,C]$—follows from the same argument
presented in Step~1. Specifically, for each fixed $\nu$, the unique outgoing solution of $Q(\nu)u = 0$ is
the deformed mode $v_+^{(\nu)}$ constructed in Appendix~E.2. Because this mode decays along $\Gamma$
but grows along the real axis, we have $v_+^{(\nu)} \notin L^2(\mathbb{R})$, meaning that $0$ cannot be
an $L^2$-eigenvalue.

Accordingly, the compact set $K_0 := [-C,C]$ satisfies all the conditions of Corollary~\ref{cor:Qnu_resolvent},
which yields the uniform bound
\[
  \|Q(\nu)^{-1}\|_{L^2 \to L^2} = O_C(1) \qquad \text{uniformly for } |\nu| \leq C.
\]
\end{remark}



% -------------------------------------------------------
% REFERENCES
% -------------------------------------------------------
\begin{thebibliography}{9}

\bibitem{DZ19}
S.~Dyatlov, M.~Zworski,
\textit{Mathematical Theory of Scattering Resonances},
Graduate Studies in Mathematics, Vol.~200, AMS, 2019.

\bibitem{Vasy2013}
A.~Vasy,
Microlocal analysis of asymptotically hyperbolic and Kerr--de~Sitter spaces,
\textit{Inventiones Mathematicae} \textbf{194} (2013), 381--513.

\bibitem{Arnold72}
V.~I.~Arnold,
Normal forms for functions near degenerate critical points, the Weyl groups of
$A_k$, $D_k$, $E_k$ and Lagrangian singularities,
\textit{Functional Analysis and Its Applications} \textbf{6} (1972), 254--272.

\bibitem{Kato66}
T.~Kato,
\textit{Perturbation Theory for Linear Operators},
Springer, Berlin, 1966.

\bibitem{Sibuya75}
Y.~Sibuya,
\textit{Global Theory of a Second Order Linear Ordinary Differential Equation
with a Polynomial Coefficient},
North-Holland, Amsterdam, 1975.

\bibitem{Malgrange64}
B.~Malgrange,
The preparation theorem for differentiable functions,
in \textit{Differential Analysis}, Tata Institute,
Oxford University Press, 1964, pp.~203--208.

\bibitem{WZ11}
J.~Wunsch, M.~Zworski,
Resolvent estimates for normally hyperbolic trapped sets,
\textit{Annales Henri Poincar\'e} \textbf{12} (2011), 1349--1385.

\bibitem{NZ15}
S.~Nonnenmacher, M.~Zworski,
Decay of correlations for normally hyperbolic trapping,
\textit{Inventiones Mathematicae} \textbf{200} (2015), 345--438.

\bibitem{Dyatlov11}
S.~Dyatlov,
Quasi-normal modes and exponential energy decay for the
Kerr--de~Sitter black hole,
\textit{Communications in Mathematical Physics} \textbf{306} (2011), 119--163.

\bibitem{Dyatlov12}
S.~Dyatlov,
Asymptotic distribution of quasi-normal modes for Kerr--de~Sitter black holes,
\textit{Annales Henri Poincar\'e} \textbf{13} (2012), 1101--1166.

\bibitem{HintzVasy2016}
P.~Hintz, A.~Vasy,
Global analysis of quasilinear wave equations on asymptotically Kerr--de~Sitter spaces,
\textit{International Mathematics Research Notices} \textbf{2016} (2016), 5355--5426.

\bibitem{Zworski2012}
M.~Zworski,
\textit{Semiclassical Analysis},
Graduate Studies in Mathematics, Vol.~138, AMS, 2012.

\bibitem{SZ91}
J.~Sj\"{o}strand, M.~Zworski,
Complex scaling and the distribution of scattering poles,
\textit{Journal of the American Mathematical Society} \textbf{4} (1991), 729--769.

\bibitem{FS11}
F.~Faure, J.~Sj\"{o}strand,
Upper bound on the density of Ruelle resonances for Anosov flows,
\textit{Communications in Mathematical Physics} \textbf{308} (2011), 325--364.

\bibitem{Hintz2022}
P.~Hintz,
Mode stability and shallow quasinormal modes of Kerr--de~Sitter black holes
away from extremality,
\textit{Journal of the European Mathematical Society} (2024), published online,
DOI: 10.4171/JEMS/1463; arXiv:2112.14431.

\bibitem{PetersenVasy2023}
O.~Petersen, A.~Vasy,
Analyticity of quasinormal modes in the Kerr and Kerr--de~Sitter spacetimes,
\textit{Communications in Mathematical Physics} \textbf{402} (2023), 2547--2580;
arXiv:2104.04500.

\bibitem{Fang2022}
A.~J.~Fang,
Linear stability of the slowly-rotating Kerr--de~Sitter family,
preprint, arXiv:2207.07902 (2022).

\bibitem{Ates2026code}
C.~Ate\c{s},
\textit{Microlocal Spectral Bifurcation in Schwarzschild--de~Sitter Geometry
[numerical ancillary files]},
Zenodo, 2026.
DOI: \texttt{10.5281/zenodo.20496573}.
GitHub: \texttt{https://github.com/canerates053/sds-jordan-chain}.

\end{thebibliography}

\end{document}
